
\Data\ID{1003032107}
\Updated{2022-08-25 10:08:46}
\Level{A}
\SubArea{Expressions with powers and roots}
\LANGen{
\Question{
The average distance of the Moon from  Earth is \( 3.84\cdot10^8\,\mathrm{m} \), and the average distance of  Earth from the Sun is \( 1.5\cdot10^{11}\mathrm{m} \).  How many times further is it from  Earth to the Sun than from Earth to the Moon?}
{about \( 390 \) times}{about \( 39 \) times}{\( 3\,900 \) times}{about \( 4 \) times}{}{}}
\LANGpl{
\Question{
Średnia odległość Księżyca od Ziemi wynosi  \( 3{,}84\cdot10^8\,\mathrm{m} \), a Ziemi od Słońca \( 1{,}5\cdot10^{11}\mathrm{m} \).  Ile razy dalej jest z Ziemi do Słońca niż z Ziemi do Księżyca?}
{około \( 390 \) razy}{około \( 39 \) razy}{\( 3900 \) razy}{około \( 4 \) razy}{}{}}
\LANGcs{
\Question{
Průměrná vzdálenost ze Země k Měsíci je \( 3{,}84\cdot10^8\,\mathrm{m} \) a průměrná vzdálenost Slunce od Země je \( 1{,}5\cdot10^{11}\mathrm{m} \).  Kolikrát dále je ze Země ke Slunci než ze Země na Měsíc?}
{okolo \( 390 \) krát}{okolo \( 39 \) krát}{\( 3900 \) krát}{okolo \( 4 \) krát}{}{}}
\LANGsk{
\Question{
Priemerná vzdialenosť zo Zeme na Mesiac je \( 3{,}84\cdot10^8\,\mathrm{m} \) a priemerná vzdialenosť zo Slnka na Zem je \( 1{,}5\cdot10^{11}\mathrm{m} \).  Koľkokrát ďalej je zo Zeme na Slnko ako zo Zeme na Mesiac?}
{okolo \( 390 \) krát}{okolo \( 39 \) krát}{\( 3900 \) krát}{okolo \( 4 \) krát}{}{}}
\LANGes{
\Question{
La distancia media de la Tierra a la Luna es \( 3{,}84\cdot10^8\,\mathrm{m} \) y la distancia media del Sol a la Tierra es \( 1{,}5\cdot10^{11}\mathrm{m} \). ¿Cuántas veces más lejos está el Sol de la Tierra que la Luna?}
{aproximadamente \( 390 \) veces}{aproximadamente \( 39 \) veces}{\( 3900 \) veces}{aproximadamente \( 4 \) veces}{}{}}
\End

\Data\ID{1003032110}
\Updated{2022-08-25 10:08:46}
\Level{A}
\SubArea{Expressions with powers and roots}
\LANGen{
\Question{
Calculate \( \sqrt[3]{-\frac8{125}}-\sqrt[3]{-2\frac{10}{27}} \).}
{\( \frac{14}{15} \)}{\( -\frac25 \)}{\( 1\frac2{97} \)}{\( 2\frac{18}{152} \)}{}{}}
\LANGpl{
\Question{
Oblicz \( \sqrt[3]{-\frac8{125}}-\sqrt[3]{-2\frac{10}{27}} \).}
{\( \frac{14}{15} \)}{\( -\frac25 \)}{\( 1\frac2{97} \)}{\( 2\frac{18}{152} \)}{}{}}
\LANGcs{
\Question{
Vypočítejte \( \sqrt[3]{-\frac8{125}}-\sqrt[3]{-2\frac{10}{27}} \).}
{\( \frac{14}{15} \)}{\( -\frac25 \)}{\( 1\frac2{97} \)}{\( 2\frac{18}{152} \)}{}{}}
\LANGsk{
\Question{
Vypočítajte \( \sqrt[3]{-\frac8{125}}-\sqrt[3]{-2\frac{10}{27}} \).}
{\( \frac{14}{15} \)}{\( -\frac25 \)}{\( 1\frac2{97} \)}{\( 2\frac{18}{152} \)}{}{}}
\LANGes{
\Question{
Calcula \( \sqrt[3]{-\frac8{125}}-\sqrt[3]{-2\frac{10}{27}} \).}
{\( \frac{14}{15} \)}{\( -\frac25 \)}{\( 1\frac2{97} \)}{\( 2\frac{18}{152} \)}{}{}}
\End

\Data\ID{1003034104}
\Updated{2022-08-25 10:08:57}
\Level{A}
\SubArea{Expressions with powers and roots}
\LANGen{
\Question{
Evaluating the expression \( \frac13\sqrt[3]4\cdot\sqrt[3]2 \) we get:}
{\( \frac23 \)}{\( \frac83 \)}{\( \frac26 \)}{\( \frac33 \)}{}{}}
\LANGpl{
\Question{
Wartością wyrażenia \( \frac13\sqrt[3]4\cdot\sqrt[3]2 \) jest:}
{\( \frac23 \)}{\( \frac83 \)}{\( \frac26 \)}{\( \frac33 \)}{}{}}
\LANGcs{
\Question{
Hodnota výrazu \( \frac13\sqrt[3]4\cdot\sqrt[3]2 \) je:}
{\( \frac23 \)}{\( \frac83 \)}{\( \frac26 \)}{\( \frac33 \)}{}{}}
\LANGsk{
\Question{
Hodnota výrazu \( \frac13\sqrt[3]4\cdot\sqrt[3]2 \) je:}
{\( \frac23 \)}{\( \frac83 \)}{\( \frac26 \)}{\( \frac33 \)}{}{}}
\LANGes{
\Question{
Evaluando la expresión \( \frac13\sqrt[3]4\cdot\sqrt[3]2 \) obtenemos:}
{\( \frac23 \)}{\( \frac83 \)}{\( \frac26 \)}{\( \frac33 \)}{}{}}
\End

\Data\ID{1003034106}
\Updated{2022-08-25 10:08:00}
\Level{A}
\SubArea{Expressions with powers and roots}
\LANGen{
\Question{
The value of the expression \( \frac{\sqrt[3]{81}}{\sqrt[3]3} \)  is:}
{\( 3 \)}{\( 3\sqrt3 \)}{\( \frac9{\sqrt[3]3} \)}{\( 27 \)}{}{}}
\LANGpl{
\Question{
Wartość wyrażenia \( \frac{\sqrt[3]{81}}{\sqrt[3]3} \)  jest równa:}
{\( 3 \)}{\( 3\sqrt3 \)}{\( \frac9{\sqrt[3]3} \)}{\( 27 \)}{}{}}
\LANGcs{
\Question{
Hodnota výrazu \( \frac{\sqrt[3]{81}}{\sqrt[3]3} \)  je:}
{\( 3 \)}{\( 3\sqrt3 \)}{\( \frac9{\sqrt[3]3} \)}{\( 27 \)}{}{}}
\LANGsk{
\Question{
Hodnota výrazu \( \frac{\sqrt[3]{81}}{\sqrt[3]3} \)  je:}
{\( 3 \)}{\( 3\sqrt3 \)}{\( \frac9{\sqrt[3]3} \)}{\( 27 \)}{}{}}
\LANGes{
\Question{
El valor de la expresión \( \frac{\sqrt[3]{81}}{\sqrt[3]3} \)  es:}
{\( 3 \)}{\( 3\sqrt3 \)}{\( \frac9{\sqrt[3]3} \)}{\( 27 \)}{}{}}
\End

\Data\ID{1003099207}
\Updated{2020-08-27 09:08:50}
\Level{A}
\SubArea{Expressions with powers and roots}
\LANGen{
\Question{
Rationalize the denominator  of   \( \frac5{\sqrt 2} \).}
{\( 2.5\sqrt2 \)}{\( 2.5 \)}{\( \frac54\sqrt 2 \)}{\( 5\sqrt2 \)}{}{}}
\LANGpl{
\Question{
Ułamek   \( \frac5{\sqrt 2} \) jest równy:}
{\( 2{,}5\sqrt2 \)}{\( 2{,}5 \)}{\( \frac54\sqrt 2 \)}{\( 5\sqrt2 \)}{}{}}
\LANGcs{
\Question{
Usměrněte zlomek \( \frac5{\sqrt 2} \).}
{\( 2{,}5\sqrt2 \)}{\( 2{,}5 \)}{\( \frac54\sqrt 2 \)}{\( 5\sqrt2 \)}{}{}}
\LANGsk{
\Question{
Usmernite zlomok \( \frac5{\sqrt 2} \).}
{\( 2{,}5\sqrt2 \)}{\( 2{,}5 \)}{\( \frac54\sqrt 2 \)}{\( 5\sqrt2 \)}{}{}}
\LANGes{
\Question{
Racionaliza el denominador  \( \frac5{\sqrt 2} \).}
{\( 2.5\sqrt2 \)}{\( 2.5 \)}{\( \frac54\sqrt 2 \)}{\( 5\sqrt2 \)}{}{}}
\End

\Data\ID{1003099208}
\Updated{2022-04-23 05:04:20}
\Level{A}
\SubArea{Expressions with powers and roots}
\LANGen{
\Question{
Rationalize the denominator   of  \( \frac{2+\sqrt{2}}{\sqrt 2} \).}
{\( \sqrt 2 +1 \)}{\( 2 \)}{\( 2\sqrt2 \)}{\( 2\sqrt2 +1 \)}{}{}}
\LANGpl{
\Question{
Ułamek  \( \frac{2+\sqrt{2}}{\sqrt 2} \) jest równy:}
{\( \sqrt 2 +1 \)}{\( 2 \)}{\( 2\sqrt2 \)}{\( 2\sqrt2 +1 \)}{}{}}
\LANGcs{
\Question{
Usměrněte zlomek \( \frac{2+\sqrt{2}}{\sqrt 2} \).}
{\( \sqrt 2 +1 \)}{\( 2 \)}{\( 2\sqrt2 \)}{\( 2\sqrt2 +1 \)}{}{}}
\LANGsk{
\Question{
Usmernite zlomok \( \frac{2+\sqrt{2}}{\sqrt 2} \).}
{\( \sqrt 2 +1 \)}{\( 2 \)}{\( 2\sqrt2 \)}{\( 2\sqrt2 +1 \)}{}{}}
\LANGes{
\Question{
Racionaliza el denominador de la fracción  \( \frac{2+\sqrt{2}}{\sqrt 2} \).}
{\( \sqrt 2 +1 \)}{\( 2 \)}{\( 2\sqrt2 \)}{\( 2\sqrt2 +1 \)}{}{}}
\End

\Data\ID{1003099209}
\Updated{2022-04-23 05:04:20}
\Level{A}
\SubArea{Expressions with powers and roots}
\LANGen{
\Question{
Rationalize the denominator  of \( \frac1{\sqrt[3]2} \).}
{\( 0.5\sqrt[3]4 \)}{\( 0.5 \)}{\( 0.5\sqrt[3]2 \)}{\( 0.5\sqrt2 \)}{}{}}
\LANGpl{
\Question{
Ułamek  \( \frac1{\sqrt[3]2} \) jest równy:}
{\( 0{,}5\sqrt[3]4 \)}{\( 0{,}5 \)}{\( 0{,}5\sqrt[3]2 \)}{\( 0{,}5\sqrt2 \)}{}{}}
\LANGcs{
\Question{
Usměrněte zlomek \( \frac1{\sqrt[3]2} \).}
{\( 0{,}5\sqrt[3]4 \)}{\( 0{,}5 \)}{\( 0{,}5\sqrt[3]2 \)}{\( 0{,}5\sqrt2 \)}{}{}}
\LANGsk{
\Question{
Usmernite zlomok \( \frac1{\sqrt[3]2} \).}
{\( 0{,}5\sqrt[3]4 \)}{\( 0{,}5 \)}{\( 0{,}5\sqrt[3]2 \)}{\( 0{,}5\sqrt2 \)}{}{}}
\LANGes{
\Question{
Racionaliza el denominador de la fracción \( \frac1{\sqrt[3]2} \).}
{\( 0.5\sqrt[3]4 \)}{\( 0.5 \)}{\( 0.5\sqrt[3]2 \)}{\( 0.5\sqrt2 \)}{}{}}
\End

\Data\ID{1003099210}
\Updated{2020-08-28 01:08:28}
\Level{A}
\SubArea{Expressions with powers and roots}
\LANGen{
\Question{
Write the fraction  \( \frac{\sqrt[3]2}{\sqrt[3]3}\) in an equivalent form which does not contain a radical in the denominator.}
{\( \frac{\sqrt[3]{18}}3 \)}{\( \frac{\sqrt[3]6}3 \)}{\( \frac23 \)}{\( \sqrt[3]6 \)}{}{}}
\LANGpl{
\Question{
Ułamek  \( \frac{\sqrt[3]2}{\sqrt[3]3}\) jest równy:}
{\( \frac{\sqrt[3]{18}}3 \)}{\( \frac{\sqrt[3]6}3 \)}{\( \frac23 \)}{\( \sqrt[3]6 \)}{}{}}
\LANGcs{
\Question{
Usměrněte zlomek \( \frac{\sqrt[3]2}{\sqrt[3]3}\).}
{\( \frac{\sqrt[3]{18}}3 \)}{\( \frac{\sqrt[3]6}3 \)}{\( \frac23 \)}{\( \sqrt[3]6 \)}{}{}}
\LANGsk{
\Question{
Usmernite zlomok \( \frac{\sqrt[3]2}{\sqrt[3]3}\).}
{\( \frac{\sqrt[3]{18}}3 \)}{\( \frac{\sqrt[3]6}3 \)}{\( \frac23 \)}{\( \sqrt[3]6 \)}{}{}}
\LANGes{
\Question{
Escribe la fracción \( \frac{\sqrt[3]2}{\sqrt[3]3}\) en forma equivalente que no contenga raíz en el denominador.}
{\( \frac{\sqrt[3]{18}}3 \)}{\( \frac{\sqrt[3]6}3 \)}{\( \frac23 \)}{\( \sqrt[3]6 \)}{}{}}
\End

\Data\ID{1003099211}
\Updated{2020-08-28 01:08:37}
\Level{A}
\SubArea{Expressions with powers and roots}
\LANGen{
\Question{
Write the fraction \(\frac{2- \sqrt3}{2\sqrt3} \) in an equivalent form which does not contain a radical in the denominator.}
{\( \frac{2\sqrt3-3}6 \)}{\( 0 \)}{\( 1 \)}{\( \frac{2\sqrt3-3}3 \)}{}{}}
\LANGpl{
\Question{
Ułamek \(\frac{2-\sqrt3 }{2\sqrt3} \) równa się:}
{\( \frac{2\sqrt3-3}6 \)}{\( 0 \)}{\( 1 \)}{\( \frac{2\sqrt3-3}3 \)}{}{}}
\LANGcs{
\Question{
Usměrněte zlomek \(\frac{2-\sqrt3 }{2\sqrt3} \).}
{\( \frac{2\sqrt3-3}6 \)}{\( 0 \)}{\( 1 \)}{\( \frac{2\sqrt3-3}3 \)}{}{}}
\LANGsk{
\Question{
Usmernite zlomok \(\frac{2-\sqrt3 }{2\sqrt3} \).}
{\( \frac{2\sqrt3-3}6 \)}{\( 0 \)}{\( 1 \)}{\( \frac{2\sqrt3-3}3 \)}{}{}}
\LANGes{
\Question{
Escribe la fracción \(\frac{2- \sqrt3}{2\sqrt3} \) en forma equivalente que no contenga raíz en el denominador,}
{\( \frac{2\sqrt3-3}6 \)}{\( 0 \)}{\( 1 \)}{\( \frac{2\sqrt3-3}3 \)}{}{}}
\End

\Data\ID{1003099504}
\Updated{2020-08-27 09:08:17}
\Level{A}
\SubArea{Expressions with powers and roots}
\LANGen{
\Question{
Rationalizing the denominator \( \frac1{\sqrt5+\sqrt7} \) we get:}
{\( \frac{\sqrt7-\sqrt5}2 \)}{\( \frac{\sqrt5+\sqrt7}2 \)}{\( \frac{\sqrt5-\sqrt7}2 \)}{\( \frac{-\sqrt7-\sqrt5}2 \)}{}{}}
\LANGpl{
\Question{
Liczba \( \frac1{\sqrt5+\sqrt7} \)  jest równa:}
{\( \frac{\sqrt7-\sqrt5}2 \)}{\( \frac{\sqrt5+\sqrt7}2 \)}{\( \frac{\sqrt5-\sqrt7}2 \)}{\( \frac{-\sqrt7-\sqrt5}2 \)}{}{}}
\LANGcs{
\Question{
Usměrněním zlomku \( \frac1{\sqrt5+\sqrt7} \) dostaneme:}
{\( \frac{\sqrt7-\sqrt5}2 \)}{\( \frac{\sqrt5+\sqrt7}2 \)}{\( \frac{\sqrt5-\sqrt7}2 \)}{\( \frac{-\sqrt7-\sqrt5}2 \)}{}{}}
\LANGsk{
\Question{
Usmernením zlomku \( \frac1{\sqrt5+\sqrt7} \) dostaneme:}
{\( \frac{\sqrt7-\sqrt5}2 \)}{\( \frac{\sqrt5+\sqrt7}2 \)}{\( \frac{\sqrt5-\sqrt7}2 \)}{\( \frac{-\sqrt7-\sqrt5}2 \)}{}{}}
\LANGes{
\Question{
Racionalizando el denominador de la fracción \( \frac1{\sqrt5+\sqrt7} \) obtenemos:}
{\( \frac{\sqrt7-\sqrt5}2 \)}{\( \frac{\sqrt5+\sqrt7}2 \)}{\( \frac{\sqrt5-\sqrt7}2 \)}{\( \frac{-\sqrt7-\sqrt5}2 \)}{}{}}
\End

\Data\ID{1003099505}
\Updated{2020-08-27 09:08:52}
\Level{A}
\SubArea{Expressions with powers and roots}
\LANGen{
\Question{
Rewrite \( \frac{2-\sqrt3}{2+\sqrt3} \) by rationalizing the denominator.}
{\( 7-4\sqrt3 \)}{\( \left(2-\sqrt3\right)\left(2+\sqrt3\right) \)}{\( \frac{7-4\sqrt3}5 \)}{\( \frac{7-4\sqrt3}7 \)}{}{}}
\LANGpl{
\Question{
Po usunięciu niewymierności z mianownika ułamka  \( \frac{2-\sqrt3}{2+\sqrt3} \) otrzymamy liczbę:}
{\( 7-4\sqrt3 \)}{\( \left(2-\sqrt3\right)\left(2+\sqrt3\right) \)}{\( \frac{7-4\sqrt3}5 \)}{\( \frac{7-4\sqrt3}7 \)}{}{}}
\LANGcs{
\Question{
Usměrněním zlomku \( \frac{2-\sqrt3}{2+\sqrt3} \) dostaneme:}
{\( 7-4\sqrt3 \)}{\( \left(2-\sqrt3\right)\left(2+\sqrt3\right) \)}{\( \frac{7-4\sqrt3}5 \)}{\( \frac{7-4\sqrt3}7 \)}{}{}}
\LANGsk{
\Question{
V danom výraze odstráňte odmocninu z menovateľa \( \frac{2-\sqrt3}{2+\sqrt3} \).}
{\( 7-4\sqrt3 \)}{\( \left(2-\sqrt3\right)\left(2+\sqrt3\right) \)}{\( \frac{7-4\sqrt3}5 \)}{\( \frac{7-4\sqrt3}7 \)}{}{}}
\LANGes{
\Question{
Racionalizando el denominador de la fracción \( \frac{2-\sqrt3}{2+\sqrt3} \) obtenemos:}
{\( 7-4\sqrt3 \)}{\( \left(2-\sqrt3\right)\left(2+\sqrt3\right) \)}{\( \frac{7-4\sqrt3}5 \)}{\( \frac{7-4\sqrt3}7 \)}{}{}}
\End

\Data\ID{1003099508}
\Updated{2020-08-27 09:08:38}
\Level{A}
\SubArea{Expressions with powers and roots}
\LANGen{
\Question{
Evaluate the expression \( \frac{2-x}{x-2} \) for \( x=2-\sqrt2 \).}
{\( -1 \)}{\( \sqrt2 - 2 \)}{\( 2 - \sqrt2 \)}{\( 1 \)}{}{}}
\LANGpl{
\Question{
Wartość wyrażenia  \( \frac{2-x}{x-2} \) dla \( x=2-\sqrt2 \) jest równa:}
{\( -1 \)}{\( \sqrt2 - 2 \)}{\( 2 - \sqrt2 \)}{\( 1 \)}{}{}}
\LANGcs{
\Question{
Vypočítejte výraz \( \frac{2-x}{x-2} \) pro \( x=2-\sqrt2 \).}
{\( -1 \)}{\( \sqrt2 - 2 \)}{\( 2 - \sqrt2 \)}{\( 1 \)}{}{}}
\LANGsk{
\Question{
Vypočítajte hodnotu výrazu \( \frac{2-x}{x-2} \) pre \( x=2-\sqrt2 \).}
{\( -1 \)}{\( \sqrt2 - 2 \)}{\( 2 - \sqrt2 \)}{\( 1 \)}{}{}}
\LANGes{
\Question{
Evalúa la expresión \( \frac{2-x}{x-2} \) para \( x=2-\sqrt2 \).}
{\( -1 \)}{\( \sqrt2 - 2 \)}{\( 2 - \sqrt2 \)}{\( 1 \)}{}{}}
\End

\Data\ID{1003099509}
\Updated{2020-10-12 10:10:44}
\Level{A}
\SubArea{Expressions with powers and roots}
\LANGen{
\Question{
Given the numbers \( x = 4+2\sqrt5 \) and \( y=6-2\sqrt5 \), the fraction \( \frac xy \) can be written in the form:}
{\( \frac{11+5\sqrt5}4 \)}{\( \frac{7\sqrt5-9}4 \)}{\( \frac{-5\sqrt5}2 \)}{\( 8\sqrt5 \)}{}{}}
\LANGpl{
\Question{
Dane są liczby \( x = 4+2\sqrt5 \)  i  \( y=6-2\sqrt5 \), iloraz \( \frac xy \) można zapisać w postaci:}
{\( \frac{11+5\sqrt5}4 \)}{\( \frac{7\sqrt5-9}4 \)}{\( \frac{-5\sqrt5}2 \)}{\( 8\sqrt5 \)}{}{}}
\LANGcs{
\Question{
Jsou dána čísla \( x = 4+2\sqrt5 \) a \( y=6-2\sqrt5 \), zlomek \( \frac xy \) může být zapsán ve tvaru:}
{\( \frac{11+5\sqrt5}4 \)}{\( \frac{7\sqrt5-9}4 \)}{\( \frac{-5\sqrt5}2 \)}{\( 8\sqrt5 \)}{}{}}
\LANGsk{
\Question{
Dané sú čísla \( x = 4+2\sqrt5 \) a \( y=6-2\sqrt5 \), zlomok \( \frac xy \) môžme zapísať v tvare:}
{\( \frac{11+5\sqrt5}4 \)}{\( \frac{7\sqrt5-9}4 \)}{\( \frac{-5\sqrt5}2 \)}{\( 8\sqrt5 \)}{}{}}
\LANGes{
\Question{
Sean \( x = 4+2\sqrt5 \) y \( y=6-2\sqrt5 \), la fracción\( \frac xy \) puede escribirse de la forma:}
{\( \frac{11+5\sqrt5}4 \)}{\( \frac{7\sqrt5-9}4 \)}{\( \frac{-5\sqrt5}2 \)}{\( 8\sqrt5 \)}{}{}}
\End

\Data\ID{1003099601}
\Updated{2020-08-27 08:08:28}
\Level{A}
\SubArea{Expressions with powers and roots}
\LANGen{
\Question{
Given the numbers \( x=1+2\sqrt2 \) and \( y=\sqrt2-1 \), calculate \( xy \).}
{\( 3-\sqrt2 \)}{\( 4-\sqrt2 \)}{\( 3 \)}{\( -\sqrt2 \)}{}{}}
\LANGpl{
\Question{
Jeżeli \( x=1+2\sqrt2 \)  i  \( y=\sqrt2-1 \), to \( xy \) jest równe:}
{\( 3-\sqrt2 \)}{\( 4-\sqrt2 \)}{\( 3 \)}{\( -\sqrt2 \)}{}{}}
\LANGcs{
\Question{
Jsou dána čísla \( x=1+2\sqrt2 \) a \( y=\sqrt2-1 \), vypočítejte \( xy \).}
{\( 3-\sqrt2 \)}{\( 4-\sqrt2 \)}{\( 3 \)}{\( -\sqrt2 \)}{}{}}
\LANGsk{
\Question{
Dané sú čísla \( x=1+2\sqrt2 \) a \( y=\sqrt2-1 \), vypočítajte \( xy \).}
{\( 3-\sqrt2 \)}{\( 4-\sqrt2 \)}{\( 3 \)}{\( -\sqrt2 \)}{}{}}
\LANGes{
\Question{
Sean \( x=1+2\sqrt2 \) a \( y=\sqrt2-1 \), calcula \( xy \).}
{\( 3-\sqrt2 \)}{\( 4-\sqrt2 \)}{\( 3 \)}{\( -\sqrt2 \)}{}{}}
\End

\Data\ID{1003099602}
\Updated{2022-04-23 05:04:20}
\Level{A}
\SubArea{Expressions with powers and roots}
\LANGen{
\Question{
Simplifying \( \frac32\sqrt8 + \sqrt{16} + \sqrt{32} - \frac13\sqrt{18} \) you get:}
{\( 4+6\sqrt2 \)}{\( 4+\sqrt{12} \)}{\( 2+\sqrt{56} \)}{\( 4+\sqrt{40} \)}{}{}}
\LANGpl{
\Question{
Liczba \( \frac32\sqrt8 + \sqrt{16} + \sqrt{32} - \frac13\sqrt{18} \)  jest równa:}
{\( 4+6\sqrt2 \)}{\( 4+\sqrt{12} \)}{\( 2+\sqrt{56} \)}{\( 4+\sqrt{40} \)}{}{}}
\LANGcs{
\Question{
Zjednodušením \( \frac32\sqrt8 + \sqrt{16} + \sqrt{32} - \frac13\sqrt{18} \) dostanete:}
{\( 4+6\sqrt2 \)}{\( 4+\sqrt{12} \)}{\( 2+\sqrt{56} \)}{\( 4+\sqrt{40} \)}{}{}}
\LANGsk{
\Question{
Zjednodušením \( \frac32\sqrt8 + \sqrt{16} + \sqrt{32} - \frac13\sqrt{18} \) dostaneme:}
{\( 4+6\sqrt2 \)}{\( 4+\sqrt{12} \)}{\( 2+\sqrt{56} \)}{\( 4+\sqrt{40} \)}{}{}}
\LANGes{
\Question{
Simplificando \( \frac32\sqrt8 + \sqrt{16} + \sqrt{32} - \frac13\sqrt{18} \) obtenemos:}
{\( 4+6\sqrt2 \)}{\( 4+\sqrt{12} \)}{\( 2+\sqrt{56} \)}{\( 4+\sqrt{40} \)}{}{}}
\End

\Data\ID{1003099603}
\Updated{2022-04-23 05:04:20}
\Level{A}
\SubArea{Expressions with powers and roots}
\LANGen{
\Question{
Calculate \( \left(2\sqrt{75}-3\sqrt{48}+2\sqrt{27}\right)^2 \).}
{\( 48 \)}{\( 192 \)}{\( 12 \)}{\( 60 \)}{}{}}
\LANGpl{
\Question{
Wartość wyrażenia \( \left(2\sqrt{75}-3\sqrt{48}+2\sqrt{27}\right)^2 \) jest równa:}
{\( 48 \)}{\( 192 \)}{\( 12 \)}{\( 60 \)}{}{}}
\LANGcs{
\Question{
Vypočítejte \( \left(2\sqrt{75}-3\sqrt{48}+2\sqrt{27}\right)^2 \).}
{\( 48 \)}{\( 192 \)}{\( 12 \)}{\( 60 \)}{}{}}
\LANGsk{
\Question{
Vypočítajte \( \left(2\sqrt{75}-3\sqrt{48}+2\sqrt{27}\right)^2 \).}
{\( 48 \)}{\( 192 \)}{\( 12 \)}{\( 60 \)}{}{}}
\LANGes{
\Question{
Calcula \( \left(2\sqrt{75}-3\sqrt{48}+2\sqrt{27}\right)^2 \).}
{\( 48 \)}{\( 192 \)}{\( 12 \)}{\( 60 \)}{}{}}
\End

\Data\ID{1003099604}
\Updated{2020-08-27 08:08:51}
\Level{A}
\SubArea{Expressions with powers and roots}
\LANGen{
\Question{
Express \( \left(\sqrt2+3\right)^2 \) in the simplest form:}
{\( 11+6\sqrt2 \)}{\( 11 \)}{\( 6\sqrt2 \)}{\( 5 \)}{}{}}
\LANGpl{
\Question{
Najprostszą postacią wyrażenia \( \left(\sqrt2+3\right)^2 \) jest:}
{\( 11+6\sqrt2 \)}{\( 11 \)}{\( 6\sqrt2 \)}{\( 5 \)}{}{}}
\LANGcs{
\Question{
Vyjádřete \( \left(\sqrt2+3\right)^2 \) nejjednodušší formou:}
{\( 11+6\sqrt2 \)}{\( 11 \)}{\( 6\sqrt2 \)}{\( 5 \)}{}{}}
\LANGsk{
\Question{
Zapíšte výraz \( \left(\sqrt2+3\right)^2 \) v najjednoduchšej forme:}
{\( 11+6\sqrt2 \)}{\( 11 \)}{\( 6\sqrt2 \)}{\( 5 \)}{}{}}
\LANGes{
\Question{
Expresa \( \left(\sqrt2+3\right)^2 \) en la forma más simple:}
{\( 11+6\sqrt2 \)}{\( 11 \)}{\( 6\sqrt2 \)}{\( 5 \)}{}{}}
\End

\Data\ID{1003099607}
\Updated{2020-08-27 08:08:51}
\Level{A}
\SubArea{Expressions with powers and roots}
\LANGen{
\Question{
Let \( \frac{m}{6-\sqrt6}=\frac{6+\sqrt6}6 \), determine \( m \).}
{\( m=5 \)}{\( m=6 \)}{\( m=1 \)}{\( m=-5 \)}{}{}}
\LANGpl{
\Question{
Równość  \( \frac{m}{6-\sqrt6}=\frac{6+\sqrt6}6 \) zachodzi dla \( m \):}
{\( m=5 \)}{\( m=6 \)}{\( m=1 \)}{\( m=-5 \)}{}{}}
\LANGcs{
\Question{
Nechť \( \frac{m}{6-\sqrt6}=\frac{6+\sqrt6}6 \), určete \( m \).}
{\( m=5 \)}{\( m=6 \)}{\( m=1 \)}{\( m=-5 \)}{}{}}
\LANGsk{
\Question{
Nech platí \( \frac{m}{6-\sqrt6}=\frac{6+\sqrt6}6 \), určte \( m \).}
{\( m=5 \)}{\( m=6 \)}{\( m=1 \)}{\( m=-5 \)}{}{}}
\LANGes{
\Question{
Sea \( \frac{m}{6-\sqrt6}=\frac{6+\sqrt6}6 \), determina \( m \).}
{\( m=5 \)}{\( m=6 \)}{\( m=1 \)}{\( m=-5 \)}{}{}}
\End

\Data\ID{1003118601}
\Updated{2020-08-27 02:08:08}
\Level{A}
\SubArea{Expressions with powers and roots}
\LANGen{
\Question{
Decide which of the following equations is false.}
{\( \sqrt5-\sqrt2=\sqrt3 \)}{\( \sqrt{15}:\sqrt3=\sqrt5 \)}{\( \sqrt5 \cdot \sqrt2 =\sqrt{10} \)}{\( \sqrt{\sqrt4}=\sqrt2 \)}{}{}}
\LANGpl{
\Question{
Wskaż równość fałszywą.}
{\( \sqrt5-\sqrt2=\sqrt3 \)}{\( \sqrt{15}:\sqrt3=\sqrt5 \)}{\( \sqrt5 \cdot \sqrt2 =\sqrt{10} \)}{\( \sqrt{\sqrt4}=\sqrt2 \)}{}{}}
\LANGcs{
\Question{
Která z uvedených rovnic neplatí?}
{\( \sqrt5-\sqrt2=\sqrt3 \)}{\( \sqrt{15}:\sqrt3=\sqrt5 \)}{\( \sqrt5 \cdot \sqrt2 =\sqrt{10} \)}{\( \sqrt{\sqrt4}=\sqrt2 \)}{}{}}
\LANGsk{
\Question{
Rozhodnite, ktorá z nasledujúcich rovnosti je nepravdivá.}
{\( \sqrt5-\sqrt2=\sqrt3 \)}{\( \sqrt{15}:\sqrt3=\sqrt5 \)}{\( \sqrt5 \cdot \sqrt2 =\sqrt{10} \)}{\( \sqrt{\sqrt4}=\sqrt2 \)}{}{}}
\LANGes{
\Question{
Decide cuál de las ecuaciones no es correcta.}
{\( \sqrt5-\sqrt2=\sqrt3 \)}{\( \sqrt{15}:\sqrt3=\sqrt5 \)}{\( \sqrt5 \cdot \sqrt2 =\sqrt{10} \)}{\( \sqrt{\sqrt4}=\sqrt2 \)}{}{}}
\End

\Data\ID{1003118602}
\Updated{2020-08-28 01:08:46}
\Level{A}
\SubArea{Expressions with powers and roots}
\LANGen{
\Question{
Choose the number equal to \( \sqrt{18}-\sqrt8 \).}
{\( \sqrt2 \)}{\( \sqrt{10} \)}{\( 10 \)}{\( 5\sqrt2 \)}{}{}}
\LANGpl{
\Question{
Wskaż liczbę równą liczbie  \( \sqrt{18}-\sqrt8 \).}
{\( \sqrt2 \)}{\( \sqrt{10} \)}{\( 10 \)}{\( 5\sqrt2 \)}{}{}}
\LANGcs{
\Question{
Které z uvedených čísel je rovno \( \sqrt{18}-\sqrt8 \)?}
{\( \sqrt2 \)}{\( \sqrt{10} \)}{\( 10 \)}{\( 5\sqrt2 \)}{}{}}
\LANGsk{
\Question{
Vyberte číslo rovné \( \sqrt{18}-\sqrt8 \).}
{\( \sqrt2 \)}{\( \sqrt{10} \)}{\( 10 \)}{\( 5\sqrt2 \)}{}{}}
\LANGes{
\Question{
Elije el número igual a \( \sqrt{18}-\sqrt8 \).}
{\( \sqrt2 \)}{\( \sqrt{10} \)}{\( 10 \)}{\( 5\sqrt2 \)}{}{}}
\End

\Data\ID{1003118604}
\Updated{2022-04-23 05:04:20}
\Level{A}
\SubArea{Expressions with powers and roots}
\LANGen{
\Question{
The sum of \( 3^{100}+3^{100}+3^{100} \)  equals:}
{\( 3^{101} \)}{\( 3^{103} \)}{\( 3^{300} \)}{\( 9^{100} \)}{}{}}
\LANGpl{
\Question{
Suma  \( 3^{100}+3^{100}+3^{100} \)  jest równa:}
{\( 3^{101} \)}{\( 3^{103} \)}{\( 3^{300} \)}{\( 9^{100} \)}{}{}}
\LANGcs{
\Question{
Součet \( 3^{100}+3^{100}+3^{100} \)  je roven:}
{\( 3^{101} \)}{\( 3^{103} \)}{\( 3^{300} \)}{\( 9^{100} \)}{}{}}
\LANGsk{
\Question{
Súčet \( 3^{100}+3^{100}+3^{100} \) sa rovná:}
{\( 3^{101} \)}{\( 3^{103} \)}{\( 3^{300} \)}{\( 9^{100} \)}{}{}}
\LANGes{
\Question{
La suma \( 3^{100}+3^{100}+3^{100} \)  es igual a}
{\( 3^{101} \)}{\( 3^{103} \)}{\( 3^{300} \)}{\( 9^{100} \)}{}{}}
\End

\Data\ID{1003118606}
\Updated{2020-08-27 02:08:08}
\Level{A}
\SubArea{Expressions with powers and roots}
\LANGen{
\Question{
The number \(2^{10}+10\cdot2^9 \) is divisible by:}
{\( 3 \)}{\( 5 \)}{\( 10 \)}{\( 11 \)}{}{}}
\LANGpl{
\Question{
Liczba  \(2^{10}+10\cdot2^9 \) jest podzielna przez:}
{\( 3 \)}{\( 5 \)}{\( 10 \)}{\( 11 \)}{}{}}
\LANGcs{
\Question{
Číslo \(2^{10}+10\cdot2^9 \) je dělitelné:}
{\( 3 \)}{\( 5 \)}{\( 10 \)}{\( 11 \)}{}{}}
\LANGsk{
\Question{
Číslo \(2^{10}+10\cdot2^9 \) je deliteľné:}
{\( 3 \)}{\( 5 \)}{\( 10 \)}{\( 11 \)}{}{}}
\LANGes{
\Question{
El número \(2^{10}+10\cdot2^9 \) es divisible por:}
{\( 3 \)}{\( 5 \)}{\( 10 \)}{\( 11 \)}{}{}}
\End

\Data\ID{1003124910}
\Updated{2022-04-23 05:04:14}
\Level{A}
\SubArea{Expressions with powers and roots}
\LANGen{
\Question{
The value of the expression \( \sqrt[3]{16}-\sqrt{50}+5\sqrt{32}-\sqrt[3]{250} \) is equal to:}
{\( -3\sqrt[3]2+15\sqrt2 \)}{\( -3\sqrt[3]2+6\sqrt2 \)}{\( 2\sqrt[3]4-26\sqrt2-5\sqrt[3]{10} \)}{\( 2\sqrt[3]4+6\sqrt2-5\sqrt[3]{10} \)}{}{}}
\LANGpl{
\Question{
Wartość wyrażenia \( \sqrt[3]{16}-\sqrt{50}+5\sqrt{32}-\sqrt[3]{250} \) jest równa:}
{\( -3\sqrt[3]2+15\sqrt2 \)}{\( -3\sqrt[3]2+6\sqrt2 \)}{\( 2\sqrt[3]4-26\sqrt2-5\sqrt[3]{10} \)}{\( 2\sqrt[3]4+6\sqrt2-5\sqrt[3]{10} \)}{}{}}
\LANGcs{
\Question{
Hodnota výrazu \( \sqrt[3]{16}-\sqrt{50}+5\sqrt{32}-\sqrt[3]{250} \) je rovna:}
{\( -3\sqrt[3]2+15\sqrt2 \)}{\( -3\sqrt[3]2+6\sqrt2 \)}{\( 2\sqrt[3]4-26\sqrt2-5\sqrt[3]{10} \)}{\( 2\sqrt[3]4+6\sqrt2-5\sqrt[3]{10} \)}{}{}}
\LANGsk{
\Question{
Hodnota výrazu \( \sqrt[3]{16}-\sqrt{50}+5\sqrt{32}-\sqrt[3]{250} \) sa rovná:}
{\( -3\sqrt[3]2+15\sqrt2 \)}{\( -3\sqrt[3]2+6\sqrt2 \)}{\( 2\sqrt[3]4-26\sqrt2-5\sqrt[3]{10} \)}{\( 2\sqrt[3]4+6\sqrt2-5\sqrt[3]{10} \)}{}{}}
\LANGes{
\Question{
El valor de la expresión \( \sqrt[3]{16}-\sqrt{50}+5\sqrt{32}-\sqrt[3]{250} \) es igual a:}
{\( -3\sqrt[3]2+15\sqrt2 \)}{\( -3\sqrt[3]2+6\sqrt2 \)}{\( 2\sqrt[3]4-26\sqrt2-5\sqrt[3]{10} \)}{\( 2\sqrt[3]4+6\sqrt2-5\sqrt[3]{10} \)}{}{}}
\End

\Data\ID{2000009401}
\Updated{2022-11-02 09:11:37}
\Level{A}
\SubArea{Expressions with powers and roots}
\LANGen{
\Question{
The expression \( 5^{12} +5^{11}+5^{10}-6\cdot 5^{10}\) equals:}
{\(5^{12}\)}{\(5^{-27}\)}{\(19\cdot 5^{10}\)}{\(-15^{23}\)}{}{}}
\LANGpl{
\Question{
Wyrażenie \( 5^{12} +5^{11}+5^{10}-6\cdot 5^{10}\) jest równe:}
{\(5^{12}\)}{\(5^{-27}\)}{\(19\cdot 5^{10}\)}{\(-15^{23}\)}{}{}}
\LANGcs{
\Question{
Výraz \( 5^{12} +5^{11}+5^{10}-6\cdot 5^{10}\) je roven:}
{\(5^{12}\)}{\(5^{-27}\)}{\(19\cdot 5^{10}\)}{\(-15^{23}\)}{}{}}
\LANGsk{
\Question{
Výraz \( 5^{12} +5^{11}+5^{10}-6\cdot 5^{10}\) sa rovná:}
{\(5^{12}\)}{\(5^{-27}\)}{\(19\cdot 5^{10}\)}{\(-15^{23}\)}{}{}}
\LANGes{
\Question{
La expresión \( 5^{12} +5^{11}+5^{10}-6\cdot 5^{10}\) es igual a:}
{\(5^{12}\)}{\(5^{-27}\)}{\(19\cdot 5^{10}\)}{\(-15^{23}\)}{}{}}
\End

\Data\ID{2000009402}
\Updated{2022-11-02 09:11:37}
\Level{A}
\SubArea{Expressions with powers and roots}
\LANGen{
\Question{
The expression \( \left(2\left(4\left(6\cdot 8^0\right)^1\right)^{-1}\right)^2\) equals:}
{\(\frac1{144}\)}{\(\frac1{1152}\)}{\(6\)}{\(0\)}{}{}}
\LANGpl{
\Question{
Wyrażenie \( \left(2\left(4\left(6\cdot 8^0\right)^1\right)^{-1}\right)^2\) jest równe:}
{\(\frac1{144}\)}{\(\frac1{1152}\)}{\(6\)}{\(0\)}{}{}}
\LANGcs{
\Question{
Výraz \( \left(2\left(4\left(6\cdot 8^0\right)^1\right)^{-1}\right)^2\) je roven:}
{\(\frac1{144}\)}{\(\frac1{1152}\)}{\(6\)}{\(0\)}{}{}}
\LANGsk{
\Question{
Výraz \( \left(2\left(4\left(6\cdot 8^0\right)^1\right)^{-1}\right)^2\) sa rovná:}
{\(\frac1{144}\)}{\(\frac1{1152}\)}{\(6\)}{\(0\)}{}{}}
\LANGes{
\Question{
La expresión \( \left(2\left(4\left(6\cdot 8^0\right)^1\right)^{-1}\right)^2\) es igual a:}
{\(\frac1{144}\)}{\(\frac1{1152}\)}{\(6\)}{\(0\)}{}{}}
\End

\Data\ID{2000009404}
\Updated{2022-11-02 09:11:37}
\Level{A}
\SubArea{Expressions with powers and roots}
\LANGen{
\Question{
The expression \( \frac{31 \cdot 10^3 \cdot 0.001}{10^4 \cdot 10^{-2}}\) equals:}
{\(0.31\)}{\(3.1\)}{\(3100\)}{\(310\)}{}{}}
\LANGpl{
\Question{
Wyrażenie \( \frac{31 \cdot 10^3 \cdot 0{,}001}{10^4 \cdot 10^{-2}}\) jest równe:}
{\(0{,}31\)}{\(3{,}1\)}{\(3100\)}{\(310\)}{}{}}
\LANGcs{
\Question{
Výraz \( \frac{31 \cdot 10^3 \cdot 0{,}001}{10^4 \cdot 10^{-2}}\) je roven:}
{\(0{,}31\)}{\(3{,}1\)}{\(3100\)}{\(310\)}{}{}}
\LANGsk{
\Question{
Výraz \( \frac{31 \cdot 10^3 \cdot 0{,}001}{10^4 \cdot 10^{-2}}\) sa rovná:}
{\(0{,}31\)}{\(3{,}1\)}{\(3100\)}{\(310\)}{}{}}
\LANGes{
\Question{
La expresión \( \frac{31 \cdot 10^3 \cdot 0.001}{10^4 \cdot 10^{-2}}\) es igual a:}
{\(0.31\)}{\(3.1\)}{\(3100\)}{\(310\)}{}{}}
\End

\Data\ID{2000009405}
\Updated{2022-11-02 09:11:37}
\Level{A}
\SubArea{Expressions with powers and roots}
\LANGen{
\Question{
The expression \( \frac{6^3\cdot 50^2}{2^3 \cdot 3^3 \cdot 10^2}\) equals:}
{\(25\)}{\(5\)}{\(1.25\)}{\(\frac1{125}\)}{}{}}
\LANGpl{
\Question{
Wyrażenie \( \frac{6^3\cdot 50^2}{2^3 \cdot 3^3 \cdot 10^2}\) jest równe:}
{\(25\)}{\(5\)}{\(1{,}25\)}{\(\frac1{125}\)}{}{}}
\LANGcs{
\Question{
Výraz  \( \frac{6^3\cdot 50^2}{2^3 \cdot 3^3 \cdot 10^2}\) je roven:}
{\(25\)}{\(5\)}{\(1{,}25\)}{\(\frac1{125}\)}{}{}}
\LANGsk{
\Question{
Výraz \( \frac{6^3\cdot 50^2}{2^3 \cdot 3^3 \cdot 10^2}\) sa rovná:}
{\(25\)}{\(5\)}{\(1{,}25\)}{\(\frac1{125}\)}{}{}}
\LANGes{
\Question{
La expresión \( \frac{6^3\cdot 50^2}{2^3 \cdot 3^3 \cdot 10^2}\) es igual a:}
{\(25\)}{\(5\)}{\(1.25\)}{\(\frac1{125}\)}{}{}}
\End

\Data\ID{2000009406}
\Updated{2022-11-02 09:11:37}
\Level{A}
\SubArea{Expressions with powers and roots}
\LANGen{
\Question{
For \(x\), \(a\), \(b \in \mathbb{R}\), \(x>0\), simplify the expression \( \sqrt{\frac{x^{a-b}}{x^{b-a}}}\).}
{\(x^{a-b}\)}{\(x^{-\frac12}\)}{\(1\)}{\(-1\)}{}{}}
\LANGpl{
\Question{
Dla \(x\), \(a\), \(b \in \mathbb{R}\), \(x>0\), uprość wyrażenie \( \sqrt{\frac{x^{a-b}}{x^{b-a}}}\).}
{\(x^{a-b}\)}{\(x^{-\frac12}\)}{\(1\)}{\(-1\)}{}{}}
\LANGcs{
\Question{
Pro dané \(x\), \(a\), \(b \in \mathbb{R}\), \(x>0\), zjednodušte výraz \( \sqrt{\frac{x^{a-b}}{x^{b-a}}}\).}
{\(x^{a-b}\)}{\(x^{-\frac12}\)}{\(1\)}{\(-1\)}{}{}}
\LANGsk{
\Question{
Zjednodušte výraz \( \sqrt{\frac{x^{a-b}}{x^{b-a}}}\) pre \(x\), \(a\), \(b \in \mathbb{R}\), \(x>0\).}
{\(x^{a-b}\)}{\(x^{-\frac12}\)}{\(1\)}{\(-1\)}{}{}}
\LANGes{
\Question{
Para \(x\), \(a\), \(b \in \mathbb{R}\), \(x>0\), simplifica la expresión \( \sqrt{\frac{x^{a-b}}{x^{b-a}}}\).}
{\(x^{a-b}\)}{\(x^{-\frac12}\)}{\(1\)}{\(-1\)}{}{}}
\End

\Data\ID{2000009407}
\Updated{2022-11-02 09:11:37}
\Level{A}
\SubArea{Expressions with powers and roots}
\LANGen{
\Question{
For \(x \in \mathbb{R}\), \(x \neq 0\), simplify the expression \( \frac{x^{-3}x^4}{(x^{-2})^3}\).}
{\(x^7\)}{\(x^2\)}{\(1\)}{\(x^{-5}\)}{}{}}
\LANGpl{
\Question{
Dla \(x \in \mathbb{R}\), \(x \neq 0\), uprość wyrażenie \( \frac{x^{-3}x^4}{(x^{-2})^3}\).}
{\(x^7\)}{\(x^2\)}{\(1\)}{\(x^{-5}\)}{}{}}
\LANGcs{
\Question{
Pro dané \(x \in \mathbb{R}\), \(x \neq 0\), zjednodušte výraz \( \frac{x^{-3}x^4}{(x^{-2})^3}\).}
{\(x^7\)}{\(x^2\)}{\(1\)}{\(x^{-5}\)}{}{}}
\LANGsk{
\Question{
Zjednodušte výraz \( \frac{x^{-3}x^4}{(x^{-2})^3}\) pre \(x \in \mathbb{R}\), \(x \neq 0\).}
{\(x^7\)}{\(x^2\)}{\(1\)}{\(x^{-5}\)}{}{}}
\LANGes{
\Question{
Para \(x \in \mathbb{R}\), \(x \neq 0\), simplifica la expresión \( \frac{x^{-3}x^4}{(x^{-2})^3}\).}
{\(x^7\)}{\(x^2\)}{\(1\)}{\(x^{-5}\)}{}{}}
\End

\Data\ID{2000009408}
\Updated{2022-11-02 09:11:37}
\Level{A}
\SubArea{Expressions with powers and roots}
\LANGen{
\Question{
Choose a false statement.}
{\( 2^4\cdot 4^2 > 2^3\cdot 4^3\)}{\(3^8=9^4\)}{\( \sqrt{3} + \sqrt{6} > \sqrt{3+6}\)}{\(\sqrt{2}\cdot \sqrt{2} = \sqrt{2+2}\)}{}{}}
\LANGpl{
\Question{
Wybierz  błędne stwierdzenie.}
{\( 2^4\cdot 4^2 > 2^3\cdot 4^3\)}{\(3^8=9^4\)}{\( \sqrt{3} + \sqrt{6} > \sqrt{3+6}\)}{\(\sqrt{2}\cdot \sqrt{2} = \sqrt{2+2}\)}{}{}}
\LANGcs{
\Question{
Vyberte nepravdivé tvrzení.}
{\( 2^4\cdot 4^2 > 2^3\cdot 4^3\)}{\(3^8=9^4\)}{\( \sqrt{3} + \sqrt{6} > \sqrt{3+6}\)}{\(\sqrt{2}\cdot \sqrt{2} = \sqrt{2+2}\)}{}{}}
\LANGsk{
\Question{
Vyberte nepravdivé tvrdenie.}
{\( 2^4\cdot 4^2 > 2^3\cdot 4^3\)}{\(3^8=9^4\)}{\( \sqrt{3} + \sqrt{6} > \sqrt{3+6}\)}{\(\sqrt{2}\cdot \sqrt{2} = \sqrt{2+2}\)}{}{}}
\LANGes{
\Question{
Elige la afirmación falsa.}
{\( 2^4\cdot 4^2 > 2^3\cdot 4^3\)}{\(3^8=9^4\)}{\( \sqrt{3} + \sqrt{6} > \sqrt{3+6}\)}{\(\sqrt{2}\cdot \sqrt{2} = \sqrt{2+2}\)}{}{}}
\End

\Data\ID{2010004801}
\Updated{2022-11-02 09:11:13}
\Level{A}
\SubArea{Expressions with powers and roots}
\LANGen{
\Question{
The value of the expression \( \sqrt[3]{81}-\sqrt{48}+5\sqrt{27}-\sqrt[3]{375} \) is equal to:}
{\( -2\sqrt[3]3+11\sqrt3 \)}{\( 2\sqrt[3]9+11\sqrt3 - 5\sqrt[3]{15}\)}{\(- 2\sqrt[3]3-\sqrt{3} \)}{\( 2\sqrt[3]9-\sqrt3-5\sqrt[3]{15} \)}{}{}}
\LANGpl{
\Question{
Wartość wyrażenia \( \sqrt[3]{81}-\sqrt{48}+5\sqrt{27}-\sqrt[3]{375} \) jest równa:}
{\( -2\sqrt[3]3+11\sqrt3 \)}{\( 2\sqrt[3]9+11\sqrt3 - 5\sqrt[3]{15}\)}{\(- 2\sqrt[3]3-\sqrt{3} \)}{\( 2\sqrt[3]9-\sqrt3-5\sqrt[3]{15} \)}{}{}}
\LANGcs{
\Question{
Hodnota výrazu \( \sqrt[3]{81}-\sqrt{48}+5\sqrt{27}-\sqrt[3]{375} \) je:}
{\( -2\sqrt[3]3+11\sqrt3 \)}{\( 2\sqrt[3]9+11\sqrt3 - 5\sqrt[3]{15}\)}{\(- 2\sqrt[3]3-\sqrt{3} \)}{\( 2\sqrt[3]9-\sqrt3-5\sqrt[3]{15} \)}{}{}}
\LANGsk{
\Question{
Aká je hodnota výrazu \( \sqrt[3]{81}-\sqrt{48}+5\sqrt{27}-\sqrt[3]{375} \)?}
{\( -2\sqrt[3]3+11\sqrt3 \)}{\( 2\sqrt[3]9+11\sqrt3 - 5\sqrt[3]{15}\)}{\(- 2\sqrt[3]3-\sqrt{3} \)}{\( 2\sqrt[3]9-\sqrt3-5\sqrt[3]{15} \)}{}{}}
\LANGes{
\Question{
El valor de la expresión \( \sqrt[3]{81}-\sqrt{48}+5\sqrt{27}-\sqrt[3]{375} \) es igual a:}
{\( -2\sqrt[3]3+11\sqrt3 \)}{\( 2\sqrt[3]9+11\sqrt3 - 5\sqrt[3]{15}\)}{\(- 2\sqrt[3]3-\sqrt{3} \)}{\( 2\sqrt[3]9-\sqrt3-5\sqrt[3]{15} \)}{}{}}
\End

\Data\ID{2010004805}
\Updated{2022-11-02 09:11:34}
\Level{A}
\SubArea{Expressions with powers and roots}
\LANGen{
\Question{
For \(x\in \mathbb{R}\),
\(x > 0\),
simplify the following expression. 
\[ 
 x\cdot \root{3}\of{x^{7}}
\]}
{\(x^{3}\root{3}\of{x}\)}{\(x^{7}\root{3}\of{x}\)}{\(x^{8}\root{3}\of{x}\)}{\(x^2\root{3}\of{x^2}\)}{}{}}
\LANGpl{
\Question{
Uprość wyrażenie dla \(x\in \mathbb{R}\),
\(x > 0\). 
\[ 
 x\cdot \root{3}\of{x^{7}}
\]}
{\(x^{3}\root{3}\of{x}\)}{\(x^{7}\root{3}\of{x}\)}{\(x^{8}\root{3}\of{x}\)}{\(x^2\root{3}\of{x^2}\)}{}{}}
\LANGcs{
\Question{
Je-li \(x\)
kladné reálné číslo, zjednodušte výraz.  
\[ 
 x\cdot \root{3}\of{x^{7}}
\]}
{\(x^{3}\root{3}\of{x}\)}{\(x^{7}\root{3}\of{x}\)}{\(x^{8}\root{3}\of{x}\)}{\(x^2\root{3}\of{x^2}\)}{}{}}
\LANGsk{
\Question{
Ak je \(x\) kladné reálne číslo, zjednodušte výraz.
 \[
 x\cdot \root{3}\of{x^{7}}
\]}
{\(x^{3}\root{3}\of{x}\)}{\(x^{7}\root{3}\of{x}\)}{\(x^{8}\root{3}\of{x}\)}{\(x^2\root{3}\of{x^2}\)}{}{}}
\LANGes{
\Question{
Para \(x\in \mathbb{R}\),
\(x > 0\),
simplifica la siguiente expresión. 
\[ 
 x\cdot \root{3}\of{x^{7}}
\]}
{\(x^{3}\root{3}\of{x}\)}{\(x^{7}\root{3}\of{x}\)}{\(x^{8}\root{3}\of{x}\)}{\(x^2\root{3}\of{x^2}\)}{}{}}
\End

\Data\ID{2010005601}
\Updated{2022-11-02 09:11:36}
\Level{A}
\SubArea{Expressions with powers and roots}
\LANGen{
\Question{
The mean distance of Uranus from the Sun is \( 4.53\cdot10^{12}\,\mathrm{m} \), and the mean distance of Mercury from the Sun is \( 5.79\cdot10^{10}\,\mathrm{m} \).  How many times farther Uranus from the Sun than Mercury is?}
{about \( 78 \) times}{about \( 780 \) times}{\( 130\) times}{about \( 8\) times}{}{}}
\LANGpl{
\Question{
Średnia odległość Urana od Słońca wynosi \( 4{,}53\cdot10^{12}\,\mathrm{m} \), a średnia odległość Merkurego od Słońca to \( 5{,}79\cdot10^{10}\,\mathrm{m} \).  Ile razy dalej od Słońca znajduje się Uran niż Markury?}
{około \( 78 \) razy}{około \( 780 \) razy}{\( 130\) razy}{około \( 8\) razy}{}{}}
\LANGcs{
\Question{
Střední vzdálenost Uranu od Slunce je \( 4{,}53\cdot10^{12}\,\mathrm{m} \) a střední vzdálenost Merkuru od Slunce je  \( 5{,}79\cdot10^{10}\,\mathrm{m} \). Kolikrát je Uran dále od Slunce než Merkur?}
{přibližně \( 78 \) krát}{přibližně \( 780 \) krát}{\( 130\) krát}{přibližně \( 8\) krát}{}{}}
\LANGsk{
\Question{
Stredná vzdialenosť Uránu od Slnka je \( 4{,}53\cdot10^{12}\,\mathrm{m} \) a stredná vzdialenosť Merkuru od Slnka je \( 5{,}79\cdot10^{10}\,\mathrm{m} \). Koľkokrát je Uran ďalej od Slnka ako Merkur?}
{okolo \( 78 \) krát}{okolo \( 780 \) krát}{\( 130\) krát}{okolo \( 8\) krát}{}{}}
\LANGes{
\Question{
La distancia media de Urano al Sol es \( 4.53\cdot10^{12}\,\mathrm{m} \), y la distancia media de Mercurio al Sol es \( 5.79\cdot10^{10}\,\mathrm{m} \).  ¿Cuántas veces está Urano más lejos del Sol que Mercurio?}
{alrededor de \( 78 \) veces}{alrededor de \( 780 \) veces}{\( 130\) veces}{alrededor de \( 8\) veces}{}{}}
\End

\Data\ID{2010005602}
\Updated{2022-11-02 09:11:36}
\Level{A}
\SubArea{Expressions with powers and roots}
\LANGen{
\Question{
Rationalize the denominator   of  \( \frac{\sqrt{3}-3}{\sqrt 3} \).}
{\(1- \sqrt 3 \)}{\( -3 \)}{\( 1+\sqrt3 \)}{\( 1-\sqrt3\)}{}{}}
\LANGpl{
\Question{
Usuń niewymierność z mianownika  \( \frac{\sqrt{3}-3}{\sqrt 3} \).}
{\(1- \sqrt 3 \)}{\( -3 \)}{\( 1+\sqrt3 \)}{\( 1-\sqrt3\)}{}{}}
\LANGcs{
\Question{
Usměrněte zlomek  \( \frac{\sqrt{3}-3}{\sqrt 3} \).}
{\(1- \sqrt 3 \)}{\( -3 \)}{\( 1+\sqrt3 \)}{\( 1-\sqrt3\)}{}{}}
\LANGsk{
\Question{
Usmernite zlomok \( \frac{\sqrt{3}-3}{\sqrt 3} \).}
{\(1- \sqrt 3 \)}{\( -3 \)}{\( 1+\sqrt3 \)}{\( 1-\sqrt3\)}{}{}}
\LANGes{
\Question{
Racionaliza el denominador de  \( \frac{\sqrt{3}-3}{\sqrt 3} \).}
{\(1- \sqrt 3 \)}{\( -3 \)}{\( 1+\sqrt3 \)}{\( 1-\sqrt3\)}{}{}}
\End

\Data\ID{2010005603}
\Updated{2022-11-02 09:11:36}
\Level{A}
\SubArea{Expressions with powers and roots}
\LANGen{
\Question{
Rationalize the denominator  of \( \frac1{\sqrt[3]5} \).}
{\( \frac{\sqrt[3]{25}}{5} \)}{\( \frac15 \)}{\( \frac{\sqrt[3]{5}}{5} \)}{\( \frac{\sqrt5}5 \)}{}{}}
\LANGpl{
\Question{
Usuń niewymierność z mianownika \( \frac1{\sqrt[3]5} \).}
{\( \frac{\sqrt[3]{25}}{5} \)}{\( \frac15 \)}{\( \frac{\sqrt[3]{5}}{5} \)}{\( \frac{\sqrt5}5 \)}{}{}}
\LANGcs{
\Question{
Usměrněte zlomek \( \frac1{\sqrt[3]5} \).}
{\( \frac{\sqrt[3]{25}}{5} \)}{\( \frac15 \)}{\( \frac{\sqrt[3]{5}}{5} \)}{\( \frac{\sqrt5}5 \)}{}{}}
\LANGsk{
\Question{
Usmernite zlomok \( \frac1{\sqrt[3]5} \).}
{\( \frac{\sqrt[3]{25}}{5} \)}{\( \frac15 \)}{\( \frac{\sqrt[3]{5}}{5} \)}{\( \frac{\sqrt5}5 \)}{}{}}
\LANGes{
\Question{
Racionaliza el denominador de \( \frac1{\sqrt[3]5} \).}
{\( \frac{\sqrt[3]{25}}{5} \)}{\( \frac15 \)}{\( \frac{\sqrt[3]{5}}{5} \)}{\( \frac{\sqrt5}5 \)}{}{}}
\End

\Data\ID{2010005604}
\Updated{2022-11-02 09:11:36}
\Level{A}
\SubArea{Expressions with powers and roots}
\LANGen{
\Question{
Simplifying \( \frac23\sqrt{27} + \sqrt{81} - \sqrt{12} + \frac14\sqrt{48} \) you get:}
{\( 9+\sqrt3 \)}{\( 3+3\sqrt{3} \)}{\( 9+5\sqrt{3} \)}{\( 15-\sqrt{3} \)}{}{}}
\LANGpl{
\Question{
Upraszczając \( \frac23\sqrt{27} + \sqrt{81} - \sqrt{12} + \frac14\sqrt{48} \) otrzymamy:}
{\( 9+\sqrt3 \)}{\( 3+3\sqrt{3} \)}{\( 9+5\sqrt{3} \)}{\( 15-\sqrt{3} \)}{}{}}
\LANGcs{
\Question{
Zjednodušením výrazu \( \frac23\sqrt{27} + \sqrt{81} - \sqrt{12} + \frac14\sqrt{48} \) dostaneme:}
{\( 9+\sqrt3 \)}{\( 3+3\sqrt{3} \)}{\( 9+5\sqrt{3} \)}{\( 15-\sqrt{3} \)}{}{}}
\LANGsk{
\Question{
Zjednodušením výrazu \( \frac23\sqrt{27} + \sqrt{81} - \sqrt{12} + \frac14\sqrt{48} \) dostaneme:}
{\( 9+\sqrt3 \)}{\( 3+3\sqrt{3} \)}{\( 9+5\sqrt{3} \)}{\( 15-\sqrt{3} \)}{}{}}
\LANGes{
\Question{
Simplificando \( \frac23\sqrt{27} + \sqrt{81} - \sqrt{12} + \frac14\sqrt{48} \) se obtiene:}
{\( 9+\sqrt3 \)}{\( 3+3\sqrt{3} \)}{\( 9+5\sqrt{3} \)}{\( 15-\sqrt{3} \)}{}{}}
\End

\Data\ID{2010005605}
\Updated{2022-11-02 09:11:36}
\Level{A}
\SubArea{Expressions with powers and roots}
\LANGen{
\Question{
Calculate \( \left(2\sqrt{72}-3\sqrt{50}+2\sqrt{32}\right)^2 \).}
{\( 50 \)}{\( 100 \)}{\( 5 \)}{\( 25 \)}{}{}}
\LANGpl{
\Question{
Oblicz \( \left(2\sqrt{72}-3\sqrt{50}+2\sqrt{32}\right)^2 \).}
{\( 50 \)}{\( 100 \)}{\( 5 \)}{\( 25 \)}{}{}}
\LANGcs{
\Question{
Vypočítejte \( \left(2\sqrt{72}-3\sqrt{50}+2\sqrt{32}\right)^2 \).}
{\( 50 \)}{\( 100 \)}{\( 5 \)}{\( 25 \)}{}{}}
\LANGsk{
\Question{
Vypočítajte \( \left(2\sqrt{72}-3\sqrt{50}+2\sqrt{32}\right)^2 \).}
{\( 50 \)}{\( 100 \)}{\( 5 \)}{\( 25 \)}{}{}}
\LANGes{
\Question{
Calcula \( \left(2\sqrt{72}-3\sqrt{50}+2\sqrt{32}\right)^2 \).}
{\( 50 \)}{\( 100 \)}{\( 5 \)}{\( 25 \)}{}{}}
\End

\Data\ID{2010005606}
\Updated{2022-11-02 09:11:36}
\Level{A}
\SubArea{Expressions with powers and roots}
\LANGen{
\Question{
The sum of \( 4^{50}+4^{50}+4^{50}+4^{50} \)  equals:}
{\( 4^{51} \)}{\( 4^{54} \)}{\( 4^{200} \)}{\( 16^{50} \)}{}{}}
\LANGpl{
\Question{
Oblicz \( 4^{50}+4^{50}+4^{50}+4^{50} \).}
{\( 4^{51} \)}{\( 4^{54} \)}{\( 4^{200} \)}{\( 16^{50} \)}{}{}}
\LANGcs{
\Question{
Součet  \( 4^{50}+4^{50}+4^{50}+4^{50} \)  je roven:}
{\( 4^{51} \)}{\( 4^{54} \)}{\( 4^{200} \)}{\( 16^{50} \)}{}{}}
\LANGsk{
\Question{
Súčet \( 4^{50}+4^{50}+4^{50}+4^{50} \) sa rovná:}
{\( 4^{51} \)}{\( 4^{54} \)}{\( 4^{200} \)}{\( 16^{50} \)}{}{}}
\LANGes{
\Question{
La suma \( 4^{50}+4^{50}+4^{50}+4^{50} \) es igual a:}
{\( 4^{51} \)}{\( 4^{54} \)}{\( 4^{200} \)}{\( 16^{50} \)}{}{}}
\End

\Data\ID{9000010501}
\Updated{2020-08-28 01:08:02}
\Level{A}
\SubArea{Expressions with powers and roots}
\LANGen{
\Question{
For \(x\in \mathbb{R}\),
\(x > 0\),
simplify the following expression. 
\[ 
 \root{3}\of{x^{5}}
\]}
{\(x\root{3}\of{x^{2}}\)}{\(x^{2}\root{3}\of{x^{2}}\)}{\(x^{3}\root{3}\of{x^{2}}\)}{\(x^{2}\root{5}\of{x^{3}}\)}{}{}}
\LANGpl{
\Question{
Dla \(x\in \mathbb{R}\),
\(x > 0\),
uprość następujące wyrażenie. 
\[ 
 \root{3}\of{x^{5}}
\]}
{\(x\root{3}\of{x^{2}}\)}{\(x^{2}\root{3}\of{x^{2}}\)}{\(x^{3}\root{3}\of{x^{2}}\)}{\(x^{2}\root{5}\of{x^{3}}\)}{}{}}
\LANGcs{
\Question{
Je-li \(x\) kladné reálné
číslo, pak je výraz \(\root{3}\of{x^{5}}\) 
roven:}
{\(x\root{3}\of{x^{2}}\)}{\(x^{2}\root{3}\of{x^{2}}\)}{\(x^{3}\root{3}\of{x^{2}}\)}{\(x^{2}\root{5}\of{x^{3}}\)}{}{}}
\LANGsk{
\Question{
Ak je \(x\) kladné reálne
číslo, tak potom je výraz \(\root{3}\of{x^{5}}\) 
rovný:}
{\(x\root{3}\of{x^{2}}\)}{\(x^{2}\root{3}\of{x^{2}}\)}{\(x^{3}\root{3}\of{x^{2}}\)}{\(x^{2}\root{5}\of{x^{3}}\)}{}{}}
\LANGes{
\Question{
Para \(x\in \mathbb{R}\),
\(x > 0\),
simplifica la siguiente expresión 
\[ 
 \root{3}\of{x^{5}}
\]}
{\(x\root{3}\of{x^{2}}\)}{\(x^{2}\root{3}\of{x^{2}}\)}{\(x^{3}\root{3}\of{x^{2}}\)}{\(x^{2}\root{5}\of{x^{3}}\)}{}{}}
\End

\Data\ID{9000010502}
\Updated{2020-08-28 01:08:42}
\Level{A}
\SubArea{Expressions with powers and roots}
\LANGen{
\Question{
For \(x\in \mathbb{R}\),
\(x > 0\),
simplify the following expression. 
\[ 
 \root{3}\of{x^{5}}\cdot \root{3}\of{x^{4}}
\]}
{\(x^{3}\)}{\(\root{3}\of{x^{12}}\)}{\(\root{3}\of{x}\)}{\(x\root{3}\of{x^{4}}\)}{}{}}
\LANGpl{
\Question{
Dla \(x\in \mathbb{R}\),
\(x > 0\),
uprość następujące wyrażenie. 
\[ 
 \root{3}\of{x^{5}}\cdot \root{3}\of{x^{4}}
\]}
{\(x^{3}\)}{\(\root{3}\of{x^{12}}\)}{\(\root{3}\of{x}\)}{\(x\root{3}\of{x^{4}}\)}{}{}}
\LANGcs{
\Question{
Je-li \(x\) kladné reálné
číslo, pak je výraz \(\root{3}\of{x^{5}}\cdot \root{3}\of{x^{4}}\) 
roven:}
{\(x^{3}\)}{\(\root{3}\of{x^{12}}\)}{\(\root{3}\of{x}\)}{\(x\root{3}\of{x^{4}}\)}{}{}}
\LANGsk{
\Question{
Ak je \(x\) kladné reálne
číslo, potom je výraz \(\root{3}\of{x^{5}}\cdot \root{3}\of{x^{4}}\) 
rovný:}
{\(x^{3}\)}{\(\root{3}\of{x^{12}}\)}{\(\root{3}\of{x}\)}{\(x\root{3}\of{x^{4}}\)}{}{}}
\LANGes{
\Question{
Para \(x\in \mathbb{R}\),
\(x > 0\),
simplifica la siguiente expresión:
\[ 
 \root{3}\of{x^{5}}\cdot \root{3}\of{x^{4}}
\]}
{\(x^{3}\)}{\(\root{3}\of{x^{12}}\)}{\(\root{3}\of{x}\)}{\(x\root{3}\of{x^{4}}\)}{}{}}
\End

\Data\ID{9000010504}
\Updated{2020-08-28 01:08:49}
\Level{A}
\SubArea{Expressions with powers and roots}
\LANGen{
\Question{
For \(x\in \mathbb{R}\),
\(x > 0\),
simplify the following expression. 
\[ 
 \root{3}\of{x^{2}} : \root{3}\of{x}
\]}
{\(\root{3}\of{x}\)}{\(x\)}{\(1\)}{\(\root{9}\of{x}\)}{}{}}
\LANGpl{
\Question{
Dla \(x\in \mathbb{R}\),
\(x > 0\),
uprość następujące wyrażenie. 
\[ 
 \root{3}\of{x^{2}} : \root{3}\of{x}
\]}
{\(\root{3}\of{x}\)}{\(x\)}{\(1\)}{\(\root{9}\of{x}\)}{}{}}
\LANGcs{
\Question{
Je-li \(x\) kladné reálné
číslo, pak je výraz \(\root{3}\of{x^{2}} : \root{3}\of{x}\) 
roven:}
{\(\root{3}\of{x}\)}{\(x\)}{\(1\)}{\(\root{9}\of{x}\)}{}{}}
\LANGsk{
\Question{
Ak je \(x\) kladné reálne
číslo, potom je výraz \(\root{3}\of{x^{2}} : \root{3}\of{x}\) 
rovný:}
{\(\root{3}\of{x}\)}{\(x\)}{\(1\)}{\(\root{9}\of{x}\)}{}{}}
\LANGes{
\Question{
Para \(x\in \mathbb{R}\),
\(x > 0\),
simplifica la siguiente expresión.
\[ 
 \root{3}\of{x^{2}} : \root{3}\of{x}
\]}
{\(\root{3}\of{x}\)}{\(x\)}{\(1\)}{\(\root{9}\of{x}\)}{}{}}
\End

\Data\ID{9000010507}
\Updated{2020-08-28 01:08:27}
\Level{A}
\SubArea{Expressions with powers and roots}
\LANGen{
\Question{
For \(x\in \mathbb{R}\),
\(x > 0\),
simplify the following expression. 
\[ 
 x^{3} : \root{}\of{x}
\]}
{\(x^{2}\root{}\of{x}\)}{\(x^{3}\root{}\of{x}\)}{\(\root{}\of{x^{3}}\)}{\(\root{6}\of{x}\)}{}{}}
\LANGpl{
\Question{
Dla\(x\in \mathbb{R}\),
\(x > 0\),
uprość następujące wyrażenie. 
\[ 
 x^{3} : \root{}\of{x}
\]}
{\(x^{2}\root{}\of{x}\)}{\(x^{3}\root{}\of{x}\)}{\(\root{}\of{x^{3}}\)}{\(\root{6}\of{x}\)}{}{}}
\LANGcs{
\Question{
Je-li \(x\) kladné reálné
číslo, pak je výraz \(x^{3} : \root{}\of{x}\) 
roven:}
{\(x^{2}\root{}\of{x}\)}{\(x^{3}\root{}\of{x}\)}{\(\root{}\of{x^{3}}\)}{\(\root{6}\of{x}\)}{}{}}
\LANGsk{
\Question{
Ak je \(x\) kladné reálne
číslo, potom je výraz \(x^{3} : \root{}\of{x}\) 
rovný:}
{\(x^{2}\root{}\of{x}\)}{\(x^{3}\root{}\of{x}\)}{\(\root{}\of{x^{3}}\)}{\(\root{6}\of{x}\)}{}{}}
\LANGes{
\Question{
Para \(x\in \mathbb{R}\),
\(x > 0\),
simplifica la siguiente expresión. 
\[ 
 x^{3} : \root{}\of{x}
\]}
{\(x^{2}\root{}\of{x}\)}{\(x^{3}\root{}\of{x}\)}{\(\root{}\of{x^{3}}\)}{\(\root{6}\of{x}\)}{}{}}
\End

\Data\ID{9000010509}
\Updated{2022-04-23 05:04:14}
\Level{A}
\SubArea{Expressions with powers and roots}
\LANGen{
\Question{
For \(x\in \mathbb{R}\),
\(x > 0\),
simplify the following expression. 
\[ 
 x\cdot \root{3}\of{x^{11}}
\]}
{\(x^{4}\root{3}\of{x^{2}}\)}{\(x^{11}\root{3}\of{x}\)}{\(x^{12}\root{3}\of{x}\)}{\(x\root{3}\of{x}\)}{}{}}
\LANGpl{
\Question{
Dla \(x\in \mathbb{R}\),
\(x > 0\),
uprość następujące wyrażenie. 
\[ 
 x\cdot \root{3}\of{x^{11}}
\]}
{\(x^{4}\root{3}\of{x^{2}}\)}{\(x^{11}\root{3}\of{x}\)}{\(x^{12}\root{3}\of{x}\)}{\(x\root{3}\of{x}\)}{}{}}
\LANGcs{
\Question{
Je-li \(x\) kladné reálné
číslo, pak je výraz \(x\cdot \root{3}\of{x^{11}}\) 
roven:}
{\(x^{4}\root{3}\of{x^{2}}\)}{\(x^{11}\root{3}\of{x}\)}{\(x^{12}\root{3}\of{x}\)}{\(x\root{3}\of{x}\)}{}{}}
\LANGsk{
\Question{
Ak je \(x\) kladné reálne
číslo, potom je výraz \(x\cdot \root{3}\of{x^{11}}\) 
rovný:}
{\(x^{4}\root{3}\of{x^{2}}\)}{\(x^{11}\root{3}\of{x}\)}{\(x^{12}\root{3}\of{x}\)}{\(x\root{3}\of{x}\)}{}{}}
\LANGes{
\Question{
Para \(x\in \mathbb{R}\),
\(x > 0\),
simplifica la siguiente expresión. 
\[ 
 x\cdot \root{3}\of{x^{11}}
\]}
{\(x^{4}\root{3}\of{x^{2}}\)}{\(x^{11}\root{3}\of{x}\)}{\(x^{12}\root{3}\of{x}\)}{\(x\root{3}\of{x}\)}{}{}}
\End

\Data\ID{9000013505}
\Updated{2020-08-28 01:08:01}
\Level{A}
\SubArea{Expressions with powers and roots}
\LANGen{
\Question{
Write the number \(3\root{3}\of{3}\) 
in an equivalent form involving just radical of a single number or radical of a single
power.}
{\(\root{3}\of{81}\)}{\(\sqrt{9}\)}{\(\root{3}\of{27}\)}{\(\root{3}\of{3^{2}}\)}{}{}}
\LANGpl{
\Question{
Zapisz liczbę  \(3\root{3}\of{3}\) 
w odpowiedniej formie używając tylko pierwiastka pojedynczej liczby lub pierwiastka pojedynczej potęgi.}
{\(\root{3}\of{81}\)}{\(\sqrt{9}\)}{\(\root{3}\of{27}\)}{\(\root{3}\of{3^{2}}\)}{}{}}
\LANGcs{
\Question{
Číslo \(3\root{3}\of{3}\) 
zapište ve tvaru jediné odmocniny.}
{\(\root{3}\of{81}\)}{\(\sqrt{9}\)}{\(\root{3}\of{27}\)}{\(\root{3}\of{3^{2}}\)}{}{}}
\LANGsk{
\Question{
Číslo \(3\root{3}\of{3}\) 
zapíšte v tvare jedinej odmocniny.}
{\(\root{3}\of{81}\)}{\(\sqrt{9}\)}{\(\root{3}\of{27}\)}{\(\root{3}\of{3^{2}}\)}{}{}}
\LANGes{
\Question{
Escribe el número \(3\root{3}\of{3}\) 
en forma de una potencia.}
{\(\root{3}\of{81}\)}{\(\sqrt{9}\)}{\(\root{3}\of{27}\)}{\(\root{3}\of{3^{2}}\)}{}{}}
\End

\Data\ID{9000013506}
\Updated{2020-08-28 01:08:45}
\Level{A}
\SubArea{Expressions with powers and roots}
\LANGen{
\Question{
Write the number \(\root{3}\of{16000}\) 
in an equivalent form with smallest possible number under the radical.}
{\(20\root{3}\of{2}\)}{\(10\root{3}\of{4}\)}{\(100\root{3}\of{2}\)}{\(4\root{3}\of{10}\)}{}{}}
\LANGpl{
\Question{
Zapisz liczbę \(\root{3}\of{16000}\) 
 w odpowiedniej formie z możliwie najmniejszą liczbą pod pierwiastkiem.}
{\(20\root{3}\of{2}\)}{\(10\root{3}\of{4}\)}{\(100\root{3}\of{2}\)}{\(4\root{3}\of{10}\)}{}{}}
\LANGcs{
\Question{
Číslo \(\root{3}\of{16000}\) 
upravte na součin racionálního čísla a odmocniny z co nejmenšího
přirozeného čísla.}
{\(20\root{3}\of{2}\)}{\(10\root{3}\of{4}\)}{\(100\root{3}\of{2}\)}{\(4\root{3}\of{10}\)}{}{}}
\LANGsk{
\Question{
Číslo \(\root{3}\of{16000}\) 
upravte na súčin racionálneho čísla a odmocniny z čo najmenšieho
prirodzeného čísla.}
{\(20\root{3}\of{2}\)}{\(10\root{3}\of{4}\)}{\(100\root{3}\of{2}\)}{\(4\root{3}\of{10}\)}{}{}}
\LANGes{
\Question{
Escribe el número \(\root{3}\of{16000}\) 
como producto de un número racional y una raíz con el radicando menor posible número natural.}
{\(20\root{3}\of{2}\)}{\(10\root{3}\of{4}\)}{\(100\root{3}\of{2}\)}{\(4\root{3}\of{10}\)}{}{}}
\End

\Data\ID{9000013508}
\Updated{2020-10-12 10:10:47}
\Level{A}
\SubArea{Expressions with powers and roots}
\LANGen{
\Question{
Write the fraction \(\frac{\sqrt{7}}
{\sqrt{3}}\) 
in an equivalent form which does not contain radical in the denominator.}
{\(\frac{\sqrt{21}}
 {3} \)}{\(\frac{\sqrt{10}}
 {3} \)}{\(\frac{\sqrt{7}+\sqrt{3}}
 {3} \)}{\(\frac{7}
{3}\)}{}{}}
\LANGpl{
\Question{
Usuń niewymierność z mianownika \(\frac{\sqrt{7}}
{\sqrt{3}}\).}
{\(\frac{\sqrt{21}}
 {3} \)}{\(\frac{\sqrt{10}}
 {3} \)}{\(\frac{\sqrt{7}+\sqrt{3}}
 {3} \)}{\(\frac{7}
{3}\)}{}{}}
\LANGcs{
\Question{
Usměrněte zlomek \(\frac{\sqrt{7}}
{\sqrt{3}}\).}
{\(\frac{\sqrt{21}}
 {3} \)}{\(\frac{\sqrt{10}}
 {3} \)}{\(\frac{\sqrt{7}+\sqrt{3}}
 {3} \)}{\(\frac{7}
{3}\)}{}{}}
\LANGsk{
\Question{
Zjednodušte zlomok \(\frac{\sqrt{7}}
{\sqrt{3}}\).}
{\(\frac{\sqrt{21}}
 {3} \)}{\(\frac{\sqrt{10}}
 {3} \)}{\(\frac{\sqrt{7}+\sqrt{3}}
 {3} \)}{\(\frac{7}
{3}\)}{}{}}
\LANGes{
\Question{
Escribe la fracción \(\frac{\sqrt{7}}
{\sqrt{3}}\) 
de forma que no contenga una raíz en el denominador,}
{\(\frac{\sqrt{21}}
 {3} \)}{\(\frac{\sqrt{10}}
 {3} \)}{\(\frac{\sqrt{7}+\sqrt{3}}
 {3} \)}{\(\frac{7}
{3}\)}{}{}}
\End

\Data\ID{9000013509}
\Updated{2020-08-28 01:08:13}
\Level{A}
\SubArea{Expressions with powers and roots}
\LANGen{
\Question{
Simplify the expression \((1 + \sqrt{2})^{2}\).}
{\(3 + 2\sqrt{2}\)}{\(3\)}{\(3 - 2\sqrt{2}\)}{\(3 + \sqrt{2}\)}{}{}}
\LANGpl{
\Question{
Uprość wyrażenie \((1 + \sqrt{2})^{2}\).}
{\(3 + 2\sqrt{2}\)}{\(3\)}{\(3 - 2\sqrt{2}\)}{\(3 + \sqrt{2}\)}{}{}}
\LANGcs{
\Question{
Upravte \((1 + \sqrt{2})^{2}\).}
{\(3 + 2\sqrt{2}\)}{\(3\)}{\(3 - 2\sqrt{2}\)}{\(3 + \sqrt{2}\)}{}{}}
\LANGsk{
\Question{
Upravte \((1 + \sqrt{2})^{2}\).}
{\(3 + 2\sqrt{2}\)}{\(3\)}{\(3 - 2\sqrt{2}\)}{\(3 + \sqrt{2}\)}{}{}}
\LANGes{
\Question{
Simplifica la expresión \((1 + \sqrt{2})^{2}\).}
{\(3 + 2\sqrt{2}\)}{\(3\)}{\(3 - 2\sqrt{2}\)}{\(3 + \sqrt{2}\)}{}{}}
\End

\Data\ID{9000013510}
\Updated{2020-10-12 10:10:09}
\Level{A}
\SubArea{Expressions with powers and roots}
\LANGen{
\Question{
Write the fraction \(\frac{1}
{1+\sqrt{2}}\) 
in an equivalent form which does not contain radical in the denominator.}
{\(\sqrt{2} - 1\)}{\(\sqrt{2}\)}{\(\frac{1}
{\sqrt{2}}\)}{\(1 -\sqrt{2}\)}{}{}}
\LANGpl{
\Question{
Usuń niewymierność z mianownika \(\frac{1}
{1+\sqrt{2}}\).}
{\(\sqrt{2} - 1\)}{\(\sqrt{2}\)}{\(\frac{1}
{\sqrt{2}}\)}{\(1 -\sqrt{2}\)}{}{}}
\LANGcs{
\Question{
Usměrněte zlomek \(\frac{1}
{1+\sqrt{2}}\).}
{\(\sqrt{2} - 1\)}{\(\sqrt{2}\)}{\(\frac{1}
{\sqrt{2}}\)}{\(1 -\sqrt{2}\)}{}{}}
\LANGsk{
\Question{
Zjednodušte zlomok \(\frac{1}
{1+\sqrt{2}}\).}
{\(\sqrt{2} - 1\)}{\(\sqrt{2}\)}{\(\frac{1}
{\sqrt{2}}\)}{\(1 -\sqrt{2}\)}{}{}}
\LANGes{
\Question{
Escribe la fracción \(\frac{1}
{1+\sqrt{2}}\) 
de forma que no contenga una raíz en el denominador.}
{\(\sqrt{2} - 1\)}{\(\sqrt{2}\)}{\(\frac{1}
{\sqrt{2}}\)}{\(1 -\sqrt{2}\)}{}{}}
\End

\Data\ID{1003032101}
\Updated{2022-08-25 10:08:42}
\Level{B}
\SubArea{Expressions with powers and roots}
\LANGen{
\Question{
The product of the numbers \( \left(\sqrt[3]{25}\cdot\sqrt5\right)^{-1} \)   and  \( \sqrt[6]5\cdot625^{\frac14} \) equals:}
{\( 1 \)}{\( \frac15 \)}{\( \sqrt5 \)}{\( 5 \)}{}{}}
\LANGpl{
\Question{
Iloczyn liczb   \( \left(\sqrt[3]{25}\cdot\sqrt5\right)^{-1} \)   i \( \sqrt[6]5\cdot625^{\frac14} \) jest równy:}
{\( 1 \)}{\( \frac15 \)}{\( \sqrt5 \)}{\( 5 \)}{}{}}
\LANGcs{
\Question{
Součin čísel \( \left(\sqrt[3]{25}\cdot\sqrt5\right)^{-1} \)   a  \( \sqrt[6]5\cdot625^{\frac14} \) je roven:}
{\( 1 \)}{\( \frac15 \)}{\( \sqrt5 \)}{\( 5 \)}{}{}}
\LANGsk{
\Question{
Súčin čísel \( \left(\sqrt[3]{25}\cdot\sqrt5\right)^{-1} \)   a  \( \sqrt[6]5\cdot625^{\frac14} \) je rovný:}
{\( 1 \)}{\( \frac15 \)}{\( \sqrt5 \)}{\( 5 \)}{}{}}
\LANGes{
\Question{
El producto de los números \( \left(\sqrt[3]{25}\cdot\sqrt5\right)^{-1} \)   y  \( \sqrt[6]5\cdot625^{\frac14} \) es igual a:}
{\( 1 \)}{\( \frac15 \)}{\( \sqrt5 \)}{\( 5 \)}{}{}}
\End

\Data\ID{1003032102}
\Updated{2022-08-25 10:08:42}
\Level{B}
\SubArea{Expressions with powers and roots}
\LANGen{
\Question{
The number \( \frac{\left(1.4\cdot10^{6}\right)\cdot\left(5.4\cdot10^{-8}\right)}{\left(3.6\cdot10^{-3}\right)\left(3.5\cdot10^{-4}\right)} \) is \( k \)-times bigger than the number \( 3000 \) for:}
{\( k=20 \)}{\( k=2 \)}{\( k=6 \)}{\( k=10 \)}{}{}}
\LANGpl{
\Question{
Liczba \( \frac{\left(1{,}4\cdot10^{6}\right)\cdot\left(5{,}4\cdot10^{-8}\right)}{\left(3{,}6\cdot10^{-3}\right)\left(3{,}5\cdot10^{-4}\right)} \) jest \( k \)-krotnie większa od liczby \( 3000 \) dla:}
{\( k=20 \)}{\( k=2 \)}{\( k=6 \)}{\( k=10 \)}{}{}}
\LANGcs{
\Question{
Číslo \( \frac{\left(1{,}4\cdot10^{6}\right)\cdot\left(5{,}4\cdot10^{-8}\right)}{\left(3{,}6\cdot10^{-3}\right)\left(3{,}5\cdot10^{-4}\right)} \) je \( k \)-krát větší než číslo \( 3000 \) pro:}
{\( k=20 \)}{\( k=2 \)}{\( k=6 \)}{\( k=10 \)}{}{}}
\LANGsk{
\Question{
Číslo \( \frac{\left(1{,}4\cdot10^{6}\right)\cdot\left(5{,}4\cdot10^{-8}\right)}{\left(3{,}6\cdot10^{-3}\right)\left(3{,}5\cdot10^{-4}\right)} \) je \( k \)-krát väčšie ako číslo \( 3000 \) pre:}
{\( k=20 \)}{\( k=2 \)}{\( k=6 \)}{\( k=10 \)}{}{}}
\LANGes{
\Question{
El número \( \frac{\left(1{,}4\cdot10^{6}\right)\cdot\left(5{,}4\cdot10^{-8}\right)}{\left(3{,}6\cdot10^{-3}\right)\left(3{,}5\cdot10^{-4}\right)} \) es \( k \) veces mayor que \( 3000 \) para:}
{\( k=20 \)}{\( k=2 \)}{\( k=6 \)}{\( k=10 \)}{}{}}
\End

\Data\ID{1003032103}
\Updated{2022-08-25 10:08:42}
\Level{B}
\SubArea{Expressions with powers and roots}
\LANGen{
\Question{
The number \( \sqrt[3]{\left(-27\right)^{-1}}\cdot81^{\frac34} \)  equals:}
{\( -9 \)}{\( -27 \)}{\( 3 \)}{\( 9 \)}{}{}}
\LANGpl{
\Question{
Liczba \( \sqrt[3]{\left(-27\right)^{-1}}\cdot81^{\frac34} \)  jest równa:}
{\( -9 \)}{\( -27 \)}{\( 3 \)}{\( 9 \)}{}{}}
\LANGcs{
\Question{
Číslo \( \sqrt[3]{\left(-27\right)^{-1}}\cdot81^{\frac34} \)  je rovno:}
{\( -9 \)}{\( -27 \)}{\( 3 \)}{\( 9 \)}{}{}}
\LANGsk{
\Question{
Číslo \( \sqrt[3]{\left(-27\right)^{-1}}\cdot81^{\frac34} \) je rovné:}
{\( -9 \)}{\( -27 \)}{\( 3 \)}{\( 9 \)}{}{}}
\LANGes{
\Question{
El número \( \sqrt[3]{\left(-27\right)^{-1}}\cdot81^{\frac34} \)  es igual a:}
{\( -9 \)}{\( -27 \)}{\( 3 \)}{\( 9 \)}{}{}}
\End

\Data\ID{1003032104}
\Updated{2022-08-25 10:08:42}
\Level{B}
\SubArea{Expressions with powers and roots}
\LANGen{
\Question{
The number \(\sqrt[3]{(-2)^4}\cdot\left(\frac1{16}\right)^{-\frac23} \) equals:}
{\( 16 \)}{\( -16 \)}{\( -4 \)}{\( 4 \)}{}{}}
\LANGpl{
\Question{
Liczba \(\sqrt[3]{(-2)^4}\cdot\left(\frac1{16}\right)^{-\frac23} \) jest równa:}
{\( 16 \)}{\( -16 \)}{\( -4 \)}{\( 4 \)}{}{}}
\LANGcs{
\Question{
Číslo \(\sqrt[3]{(-2)^4}\cdot\left(\frac1{16}\right)^{-\frac23} \) je rovno:}
{\( 16 \)}{\( -16 \)}{\( -4 \)}{\( 4 \)}{}{}}
\LANGsk{
\Question{
Číslo \(\sqrt[3]{(-2)^4}\cdot\left(\frac1{16}\right)^{-\frac23} \) je rovné:}
{\( 16 \)}{\( -16 \)}{\( -4 \)}{\( 4 \)}{}{}}
\LANGes{
\Question{
El número \(\sqrt[3]{(-2)^4}\cdot\left(\frac1{16}\right)^{-\frac23} \) es igual a:}
{\( 16 \)}{\( -16 \)}{\( -4 \)}{\( 4 \)}{}{}}
\End

\Data\ID{1003032105}
\Updated{2022-08-25 10:08:42}
\Level{B}
\SubArea{Expressions with powers and roots}
\LANGen{
\Question{
Let \( a=5^{-1}\cdot\sqrt5 \), \( b=5^{-\frac32}\cdot25 \), \( c=125^{\frac14}:5^{-3} \), \( d=5^{\frac13}\cdot25^{-\frac13} \). Which of these numbers is the biggest?}
{\( c \)}{\( a \)}{\( b \)}{\( d \)}{}{}}
\LANGpl{
\Question{
Spośród podanych liczb: \( a=5^{-1}\cdot\sqrt5 \), \( b=5^{-\frac32}\cdot25 \), \( c=125^{\frac14}:5^{-3} \), \( d=5^{\frac13}\cdot25^{-\frac13} \), największą jest:}
{\( c \)}{\( a \)}{\( b \)}{\( d \)}{}{}}
\LANGcs{
\Question{
Platí, že \( a=5^{-1}\cdot\sqrt5 \), \( b=5^{-\frac32}\cdot25 \), \( c=125^{\frac14}:5^{-3} \), \( d=5^{\frac13}\cdot25^{-\frac13} \). Které z těchto čísel je největší?}
{\( c \)}{\( a \)}{\( b \)}{\( d \)}{}{}}
\LANGsk{
\Question{
Ktoré z daných čísel \( a=5^{-1}\cdot\sqrt5 \), \( b=5^{-\frac32}\cdot25 \), \( c=125^{\frac14}:5^{-3} \), \( d=5^{\frac13}\cdot25^{-\frac13} \) je najväčšie?}
{\( c \)}{\( a \)}{\( b \)}{\( d \)}{}{}}
\LANGes{
\Question{
Sean \( a=5^{-1}\cdot\sqrt5 \), \( b=5^{-\frac32}\cdot25 \), \( c=125^{\frac14}:5^{-3} \), \( d=5^{\frac13}\cdot25^{-\frac13} \). ¿Cuál de los números es el mayor?}
{\( c \)}{\( a \)}{\( b \)}{\( d \)}{}{}}
\End

\Data\ID{1003032106}
\Updated{2022-08-25 10:08:46}
\Level{B}
\SubArea{Expressions with powers and roots}
\LANGen{
\Question{
The value of the expression  \( \frac{16^{\frac32}-16^{\frac54}}{40} \) is:}
{\( 0.8 \)}{\( 0.4 \)}{\( \frac1{20} \)}{\( \frac1{40} \)}{}{}}
\LANGpl{
\Question{
Wartość wyrażenia  \( \frac{16^{\frac32}-16^{\frac54}}{40} \) to:}
{\( 0{,}8 \)}{\( 0{,}4 \)}{\( \frac1{20} \)}{\( \frac1{40} \)}{}{}}
\LANGcs{
\Question{
Hodnota výrazu  \( \frac{16^{\frac32}-16^{\frac54}}{40} \) je:}
{\( 0{,}8 \)}{\( 0{,}4 \)}{\( \frac1{20} \)}{\( \frac1{40} \)}{}{}}
\LANGsk{
\Question{
Aká je hodnota výrazu \( \frac{16^{\frac32}-16^{\frac54}}{40} \)?}
{\( 0{,}8 \)}{\( 0{,}4 \)}{\( \frac1{20} \)}{\( \frac1{40} \)}{}{}}
\LANGes{
\Question{
El valor de la expresión \( \frac{16^{\frac32}-16^{\frac54}}{40} \) es:}
{\( 0{,}8 \)}{\( 0{,}4 \)}{\( \frac1{20} \)}{\( \frac1{40} \)}{}{}}
\End

\Data\ID{1003032108}
\Updated{2022-08-25 10:08:46}
\Level{B}
\SubArea{Expressions with powers and roots}
\LANGen{
\Question{
Write the number \( 4^{-5}\cdot8^2 \) as \( 2^m \), where \( m \) is an integer.}
{\( 2^{-4} \)}{\( 2^{-3} \)}{\( 2^{-16} \)}{\( 2^{-30} \)}{}{}}
\LANGpl{
\Question{
Zapisz liczbę  \( 4^{-5}\cdot8^2 \) w postaci \( 2^m \), gdzie \( m \) jest liczbą całkowitą.}
{\( 2^{-4} \)}{\( 2^{-3} \)}{\( 2^{-16} \)}{\( 2^{-30} \)}{}{}}
\LANGcs{
\Question{
Napište číslo \( 4^{-5}\cdot8^2 \) v tvaru \( 2^m \), kde \( m \) je celé číslo.}
{\( 2^{-4} \)}{\( 2^{-3} \)}{\( 2^{-16} \)}{\( 2^{-30} \)}{}{}}
\LANGsk{
\Question{
Napíšte číslo \( 4^{-5}\cdot8^2 \) v tvare \( 2^m \), kde \( m \) je celé číslo.}
{\( 2^{-4} \)}{\( 2^{-3} \)}{\( 2^{-16} \)}{\( 2^{-30} \)}{}{}}
\LANGes{
\Question{
Escribe el número \( 4^{-5}\cdot8^2 \) en la forma  \( 2^m \), donde \( m \) es un número entero.}
{\( 2^{-4} \)}{\( 2^{-3} \)}{\( 2^{-16} \)}{\( 2^{-30} \)}{}{}}
\End

\Data\ID{1003032109}
\Updated{2022-08-25 10:08:46}
\Level{B}
\SubArea{Expressions with powers and roots}
\LANGen{
\Question{
Calculate \( \left(3.4\cdot10^7\right)\cdot\left(4\cdot10^{-5}\right) \)  and give the result in the scientific   notation.}
{\( 1\,360 = 1.36\cdot10^3\)}{\(1\,360 = 13.6\cdot10^2 \)}{\( 1\,360 = 136\cdot10^1\)}{\( 1\,360\,000\,000\,000 = 1.36\cdot10^{12}\)}{}{}}
\LANGpl{
\Question{
Oblicz \( \left(3{,}4\cdot10^7\right)\cdot\left(4\cdot10^{-5}\right) \)  i odpowiedź podaj w notacji wykładniczej i dziesiętnej.}
{\( 1\,360 = 1.36\cdot10^3\)}{\(1\,360 = 13.6\cdot10^2 \)}{\( 1\,360 = 136\cdot10^1\)}{\( 1\,360\,000\,000\,000 = 1.36\cdot10^{12}\)}{}{}}
\LANGcs{
\Question{
Vypočítejte \( \left(3{,}4\cdot10^7\right)\cdot\left(4\cdot10^{-5}\right) \)  a výsledek zapište vědeckým zápisem čísla.}
{\( 1\,360 = 1.36\cdot10^3\)}{\(1\,360 = 13.6\cdot10^2 \)}{\( 1\,360 = 136\cdot10^1\)}{\( 1\,360\,000\,000\,000 = 1.36\cdot10^{12}\)}{}{}}
\LANGsk{
\Question{
Vypočítajte \( \left(3{,}4\cdot10^7\right)\cdot\left(4\cdot10^{-5}\right) \)  a výsledok zapíšte v tvare vedeckého zápisu čísel.}
{\( 1\,360 = 1.36\cdot10^3\)}{\(1\,360 = 13.6\cdot10^2 \)}{\( 1\,360 = 136\cdot10^1\)}{\( 1\,360\,000\,000\,000 = 1.36\cdot10^{12}\)}{}{}}
\LANGes{
\Question{
Calcula \( \left(3.4\cdot10^7\right)\cdot\left(4\cdot10^{-5}\right) \)  y escríbe el resultado en notación científica.}
{\( 1\,360 = 1.36\cdot10^3\)}{\(1\,360 = 13.6\cdot10^2 \)}{\( 1\,360 = 136\cdot10^1\)}{\( 1\,360\,000\,000\,000 = 1.36\cdot10^{12}\)}{}{}}
\End

\Data\ID{1003032201}
\Updated{2022-08-25 10:08:46}
\Level{B}
\SubArea{Expressions with powers and roots}
\LANGen{
\Question{
Let \( a=\left(\frac23\right)^{\sqrt7-\sqrt3} \)  and \( b=\left(\frac23\right)^{\sqrt3+2} \). 
Which of the following relations is true for \( a \) and \( b \)?}
{\( a>b \)}{\( b>a \)}{\( a=b \)}{\( a\cdot b=\left(\frac23\right)^{2\sqrt7} \)}{}{}}
\LANGpl{
\Question{
Dane są liczby  \( a=\left(\frac23\right)^{\sqrt7-\sqrt3} \)  i \( b=\left(\frac23\right)^{\sqrt3+2} \). 
Prawdą jest, że:}
{\( a>b \)}{\( b>a \)}{\( a=b \)}{\( a\cdot b=\left(\frac23\right)^{2\sqrt7} \)}{}{}}
\LANGcs{
\Question{
Jsou dána čísla \( a=\left(\frac23\right)^{\sqrt7-\sqrt3} \)  a \( b=\left(\frac23\right)^{\sqrt3+2} \). 
Které tvrzení o vztahu \( a \) a \( b \) je pravdivé?}
{\( a>b \)}{\( b>a \)}{\( a=b \)}{\( a\cdot b=\left(\frac23\right)^{2\sqrt7} \)}{}{}}
\LANGsk{
\Question{
Dané sú čísla \( a=\left(\frac23\right)^{\sqrt7-\sqrt3} \)  a \( b=\left(\frac23\right)^{\sqrt3+2} \). 
Ktoré tvrdenie o vzťahu \( a \) a \( b \) je pravdivé?}
{\( a>b \)}{\( b>a \)}{\( a=b \)}{\( a\cdot b=\left(\frac23\right)^{2\sqrt7} \)}{}{}}
\LANGes{
\Question{
Sean \( a=\left(\frac23\right)^{\sqrt7-\sqrt3} \)  y \( b=\left(\frac23\right)^{\sqrt3+2} \). 
¿Cuál de las (in)ecuaciones es correcta para \( a \) y \( b \)?}
{\( a>b \)}{\( b>a \)}{\( a=b \)}{\( a\cdot b=\left(\frac23\right)^{2\sqrt7} \)}{}{}}
\End

\Data\ID{1003032202}
\Updated{2022-08-25 10:08:46}
\Level{B}
\SubArea{Expressions with powers and roots}
\LANGen{
\Question{
Writing the expression \( \frac{(4x)^2\cdot(x:y)^{-2}}{(0.5)^{-4}\cdot y^{-1}}\), \( x\neq0\), \( y\neq0 \) in a  simplified   form, we  will get:}
{\( y^3 \)}{\( \frac12y^{-1} \)}{\( 8y^3 \)}{\( 256y^{-1} \)}{}{}}
\LANGpl{
\Question{
Zapisując wyrażenie  \( \frac{(4x)^2\cdot(x:y)^{-2}}{(0{,}5)^{-4}\cdot y^{-1}}\), \( x\neq0\), \( y\neq0 \) w prostszej postaci, otrzymamy:}
{\( y^3 \)}{\( \frac12y^{-1} \)}{\( 8y^3 \)}{\( 256y^{-1} \)}{}{}}
\LANGcs{
\Question{
Když výraz \( \frac{(4x)^2\cdot(x:y)^{-2}}{(0{,}5)^{-4}\cdot y^{-1}}\), \( x\neq0\), \( y\neq0 \) zjednodušíme, dostaneme:}
{\( y^3 \)}{\( \frac12y^{-1} \)}{\( 8y^3 \)}{\( 256y^{-1} \)}{}{}}
\LANGsk{
\Question{
Zjednodušte daný výraz \( \frac{(4x)^2\cdot(x:y)^{-2}}{(0{,}5)^{-4}\cdot y^{-1}}\), ak \( x\neq0\), \( y\neq0 \).}
{\( y^3 \)}{\( \frac12y^{-1} \)}{\( 8y^3 \)}{\( 256y^{-1} \)}{}{}}
\LANGes{
\Question{
Al simplificar la expresión \( \frac{(4x)^2\cdot(x:y)^{-2}}{(0{,}5)^{-4}\cdot y^{-1}}\), \( x\neq0\), \( y\neq0 \) obtenemos:}
{\( y^3 \)}{\( \frac12y^{-1} \)}{\( 8y^3 \)}{\( 256y^{-1} \)}{}{}}
\End

\Data\ID{1003032203}
\Updated{2022-08-25 10:08:46}
\Level{B}
\SubArea{Expressions with powers and roots}
\LANGen{
\Question{
Writing the expression \( \frac{\sqrt[3]3\cdot9\cdot\sqrt{27}\cdot\sqrt[6]{81}}{81\cdot\sqrt[3]{\frac13}\cdot\sqrt[4]9} \) in the form of a  power of \( 3 \), we will get:}
{\( 3^{\frac13} \)}{\( 3^{-\frac56} \)}{\( 3^{-\frac12} \)}{\( 3^{-\frac23} \)}{}{}}
\LANGpl{
\Question{
Zapisując wyrażenie  \( \frac{\sqrt[3]3\cdot9\cdot\sqrt{27}\cdot\sqrt[6]{81}}{81\cdot\sqrt[3]{\frac13}\cdot\sqrt[4]9} \)  w postaci potęgi liczby  \( 3 \), otrzymamy:}
{\( 3^{\frac13} \)}{\( 3^{-\frac56} \)}{\( 3^{-\frac12} \)}{\( 3^{-\frac23} \)}{}{}}
\LANGcs{
\Question{
Zjednodušte daný výraz \( \frac{\sqrt[3]3\cdot9\cdot\sqrt{27}\cdot\sqrt[6]{81}}{81\cdot\sqrt[3]{\frac13}\cdot\sqrt[4]9} \) a výsledek zapište ve tvaru mocniny čísla \( 3 \).}
{\( 3^{\frac13} \)}{\( 3^{-\frac56} \)}{\( 3^{-\frac12} \)}{\( 3^{-\frac23} \)}{}{}}
\LANGsk{
\Question{
Zjednodušte daný výraz \( \frac{\sqrt[3]3\cdot9\cdot\sqrt{27}\cdot\sqrt[6]{81}}{81\cdot\sqrt[3]{\frac13}\cdot\sqrt[4]9} \) a výsledok zapíšte v tvare mocniny čísla \( 3 \).}
{\( 3^{\frac13} \)}{\( 3^{-\frac56} \)}{\( 3^{-\frac12} \)}{\( 3^{-\frac23} \)}{}{}}
\LANGes{
\Question{
Simplifica \( \frac{\sqrt[3]3\cdot9\cdot\sqrt{27}\cdot\sqrt[6]{81}}{81\cdot\sqrt[3]{\frac13}\cdot\sqrt[4]9} \) y escríbe el resultado en forma de potencia de \( 3 \).}
{\( 3^{\frac13} \)}{\( 3^{-\frac56} \)}{\( 3^{-\frac12} \)}{\( 3^{-\frac23} \)}{}{}}
\End

\Data\ID{1003032204}
\Updated{2022-08-25 10:08:46}
\Level{B}
\SubArea{Expressions with powers and roots}
\LANGen{
\Question{
Order the  numbers \( a=5^{\sqrt5} \), \( b=25\sqrt5 \), \( c=125^{\frac45} \), \( d=25^{1.1} \)  from the smallest to the largest.}
{\( d < a < c < b \)}{\( a < b < c < d \)}{\( d < a < b < c \)}{\( a < b < d < c \)}{}{}}
\LANGpl{
\Question{
Uporządkuj liczby  \( a=5^{\sqrt5} \), \( b=25\sqrt5 \), \( c=125^{\frac45} \), \( d=25^{1{,}1} \)  od najmniejszej do największej.}
{\( d < a < c < b \)}{\( a < b < c < d \)}{\( d < a < b < c \)}{\( a < b < d < c \)}{}{}}
\LANGcs{
\Question{
Seřaďte čísla \( a=5^{\sqrt5} \), \( b=25\sqrt5 \), \( c=125^{\frac45} \), \( d=25^{1{,}1} \) od nejmenšího po největší.}
{\( d < a < c < b \)}{\( a < b < c < d \)}{\( d < a < b < c \)}{\( a < b < d < c \)}{}{}}
\LANGsk{
\Question{
Usporiadajte čísla \( a=5^{\sqrt5} \), \( b=25\sqrt5 \), \( c=125^{\frac45} \), \( d=25^{1{,}1} \) od najmenšieho po najväčšie.}
{\( d < a < c < b \)}{\( a < b < c < d \)}{\( d < a < b < c \)}{\( a < b < d < c \)}{}{}}
\LANGes{
\Question{
Ordena los números \( a=5^{\sqrt5} \), \( b=25\sqrt5 \), \( c=125^{\frac45} \), \( d=25^{1{,}1} \) de menor a mayor.}
{\( d < a < c < b \)}{\( a < b < c < d \)}{\( d < a < b < c \)}{\( a < b < d < c \)}{}{}}
\End

\Data\ID{1003032205}
\Updated{2022-08-25 10:08:46}
\Level{B}
\SubArea{Expressions with powers and roots}
\LANGen{
\Question{
Order the  numbers  \( a=4.5^{-2} \), \( b=\sqrt{0.4}^{\sqrt{12}} \), \( c=2.5^{-2.4} \), \( d=0.064 \) from the smallest to the largest.}
{\( a < d < c < b \)}{\( d < a < c < b \)}{\( a < d < b < c \)}{\( c < a < b < d \)}{}{}}
\LANGpl{
\Question{
Uporządkuj liczby   \( a=4{,}5^{-2} \), \( b=\sqrt{0{,}4}^{\sqrt{12}} \), \( c=2{,}5^{-2{,}4} \), \( d=0{,}064 \) od najmniejszej do największej.}
{\( a < d < c < b \)}{\( d < a < c < b \)}{\( a < d < b < c \)}{\( c < a < b < d \)}{}{}}
\LANGcs{
\Question{
Seřaďte čísla \( a=4{,}5^{-2} \), \( b=\sqrt{0{,}4}^{\sqrt{12}} \), \( c=2{,}5^{-2{,}4} \), \( d=0{,}064 \) od nejmenšího po největší.}
{\( a < d < c < b \)}{\( d < a < c < b \)}{\( a < d < b < c \)}{\( c < a < b < d \)}{}{}}
\LANGsk{
\Question{
Usporiadajte čísla \( a=4{,}5^{-2} \), \( b=\sqrt{0{,}4}^{\sqrt{12}} \), \( c=2{,}5^{-2{,}4} \), \( d=0{,}064 \) od najmenšieho po najväčšie.}
{\( a < d < c < b \)}{\( d < a < c < b \)}{\( a < d < b < c \)}{\( c < a < b < d \)}{}{}}
\LANGes{
\Question{
Ordena los números  \( a=4.5^{-2} \), \( b=\sqrt{0.4}^{\sqrt{12}} \), \( c=2.5^{-2.4} \), \( d=0.064 \) de menor a mayor.}
{\( a < d < c < b \)}{\( d < a < c < b \)}{\( a < d < b < c \)}{\( c < a < b < d \)}{}{}}
\End

\Data\ID{1003032206}
\Updated{2022-08-25 10:08:50}
\Level{B}
\SubArea{Expressions with powers and roots}
\LANGen{
\Question{
Simplifying the expression \( \frac{\left(\frac13\right)^{-3}\cdot27^{-2}}{9^{-5}\cdot\sqrt{3^4}} \) we will get:}
{\( 3^5 \)}{\( 3^{-11} \)}{\( 3^2 \)}{\( 3^{-5} \)}{}{}}
\LANGpl{
\Question{
Wyrażenie  \( \frac{\left(\frac13\right)^{-3}\cdot27^{-2}}{9^{-5}\cdot\sqrt{3^4}} \) zapisane w postaci potęgi o podstawie 3 ma postać:}
{\( 3^5 \)}{\( 3^{-11} \)}{\( 3^2 \)}{\( 3^{-5} \)}{}{}}
\LANGcs{
\Question{
Zjednodušením výrazu \( \frac{\left(\frac13\right)^{-3}\cdot27^{-2}}{9^{-5}\cdot\sqrt{3^4}} \) dostaneme:}
{\( 3^5 \)}{\( 3^{-11} \)}{\( 3^2 \)}{\( 3^{-5} \)}{}{}}
\LANGsk{
\Question{
Zjednodušením výrazu \( \frac{\left(\frac13\right)^{-3}\cdot27^{-2}}{9^{-5}\cdot\sqrt{3^4}} \) dostaneme:}
{\( 3^5 \)}{\( 3^{-11} \)}{\( 3^2 \)}{\( 3^{-5} \)}{}{}}
\LANGes{
\Question{
Simplificando la expresión \( \frac{\left(\frac13\right)^{-3}\cdot27^{-2}}{9^{-5}\cdot\sqrt{3^4}} \) obtenemos:}
{\( 3^5 \)}{\( 3^{-11} \)}{\( 3^2 \)}{\( 3^{-5} \)}{}{}}
\End

\Data\ID{1003032207}
\Updated{2022-08-25 10:08:50}
\Level{B}
\SubArea{Expressions with powers and roots}
\LANGen{
\Question{
Let \( x=2^{4\sqrt2} \), \( y=4^{\frac{\sqrt2}2} \) and \( z=8^{-\frac{\sqrt2}3} \). Which two numbers are multiplicative inverses of each other?}
{\( y\), \( z \)}{\( x\), \( y \)}{\( x\), \( z \)}{no two numbers}{}{}}
\LANGpl{
\Question{
Spośród  liczb \( x=2^{4\sqrt2} \), \( y=4^{\frac{\sqrt2}2} \) i \( z=8^{-\frac{\sqrt2}3} \) wybierz dwie takie liczby, z których jedna jest odwrotnością drugiej.}
{\( y\), \( z \)}{\( x\), \( y \)}{\( x\), \( z \)}{nie ma takich liczb}{}{}}
\LANGcs{
\Question{
Která dvě z daných čísel \( x=2^{4\sqrt2} \), \( y=4^{\frac{\sqrt2}2} \) a \( z=8^{-\frac{\sqrt2}3} \) jsou navzájem inverzní?}
{\( y\), \( z \)}{\( x\), \( y \)}{\( x\), \( z \)}{žádné dvě čísla}{}{}}
\LANGsk{
\Question{
Ktoré dve z daných čísle \( x=2^{4\sqrt2} \), \( y=4^{\frac{\sqrt2}2} \) a \( z=8^{-\frac{\sqrt2}3} \) sú si navzájom inverzné?}
{\( y\), \( z \)}{\( x\), \( y \)}{\( x\), \( z \)}{žiadne dve čísla}{}{}}
\LANGes{
\Question{
¿Cuáles de los números \( x=2^{4\sqrt2} \), \( y=4^{\frac{\sqrt2}2} \) y \( z=8^{-\frac{\sqrt2}3} \) son inversos?}
{\( y\), \( z \)}{\( x\), \( y \)}{\( x\), \( z \)}{ningunos}{}{}}
\End

\Data\ID{1003032208}
\Updated{2022-08-25 10:08:50}
\Level{B}
\SubArea{Expressions with powers and roots}
\LANGen{
\Question{
Write the number \( \frac13\cdot9^{\pi+\frac12}:81^{2\pi} \)  in the form of \( a^x \), where \( a \) is a natural number.}
{\( 3^{-6\pi} \)}{\( 3^{6\pi} \)}{\( 3^{-\frac12+3\pi} \)}{\( 9^{-\frac12+3\pi} \)}{}{}}
\LANGpl{
\Question{
Przedstaw liczbę  \( \frac13\cdot9^{\pi+\frac12}:81^{2\pi} \)  w postaci  \( a^x \), gdzie \( a \) jest liczbą naturalną.}
{\( 3^{-6\pi} \)}{\( 3^{6\pi} \)}{\( 3^{-\frac12+3\pi} \)}{\( 9^{-\frac12+3\pi} \)}{}{}}
\LANGcs{
\Question{
Zapište číslo \( \frac13\cdot9^{\pi+\frac12}:81^{2\pi} \) v tvaru \( a^x \), kde \( a \) je přirozené číslo.}
{\( 3^{-6\pi} \)}{\( 3^{6\pi} \)}{\( 3^{-\frac12+3\pi} \)}{\( 9^{-\frac12+3\pi} \)}{}{}}
\LANGsk{
\Question{
Zapíšte číslo \( \frac13\cdot9^{\pi+\frac12}:81^{2\pi} \) v tvare \( a^x \), kde \( a \) je prirodzené číslo.}
{\( 3^{-6\pi} \)}{\( 3^{6\pi} \)}{\( 3^{-\frac12+3\pi} \)}{\( 9^{-\frac12+3\pi} \)}{}{}}
\LANGes{
\Question{
Escriba el número \( \frac13\cdot9^{\pi+\frac12}:81^{2\pi} \) en la forma \( a^x \), donde \( a \) es un número natural.}
{\( 3^{-6\pi} \)}{\( 3^{6\pi} \)}{\( 3^{-\frac12+3\pi} \)}{\( 9^{-\frac12+3\pi} \)}{}{}}
\End

\Data\ID{1003032209}
\Updated{2022-08-25 10:08:50}
\Level{B}
\SubArea{Expressions with powers and roots}
\LANGen{
\Question{
What is the value of \( \left(2^{\sqrt7-\sqrt2}\right)^{\sqrt7+\sqrt2} \)?}
{\( 32 \)}{\( 2^{\sqrt{45}} \)}{\( 2^9 \)}{\( 1024 \)}{}{}}
\LANGpl{
\Question{
Oblicz  \( \left(2^{\sqrt7-\sqrt2}\right)^{\sqrt7+\sqrt2} \)?}
{\( 32 \)}{\( 2^{\sqrt{45}} \)}{\( 2^9 \)}{\( 1024 \)}{}{}}
\LANGcs{
\Question{
Jaká je hodnota výrazu \( \left(2^{\sqrt7-\sqrt2}\right)^{\sqrt7+\sqrt2} \)?}
{\( 32 \)}{\( 2^{\sqrt{45}} \)}{\( 2^9 \)}{\( 1024 \)}{}{}}
\LANGsk{
\Question{
Aká je hodnota výrazu \( \left(2^{\sqrt7-\sqrt2}\right)^{\sqrt7+\sqrt2} \)?}
{\( 32 \)}{\( 2^{\sqrt{45}} \)}{\( 2^9 \)}{\( 1024 \)}{}{}}
\LANGes{
\Question{
¿Cuál es el valor de la expresión \( \left(2^{\sqrt7-\sqrt2}\right)^{\sqrt7+\sqrt2} \)?}
{\( 32 \)}{\( 2^{\sqrt{45}} \)}{\( 2^9 \)}{\( 1024 \)}{}{}}
\End

\Data\ID{1003032210}
\Updated{2022-08-25 10:08:50}
\Level{B}
\SubArea{Expressions with powers and roots}
\LANGen{
\Question{
Which of the given numbers belongs to the interval  \( (-5;5) \)?}
{\( 3\left(\sqrt{0.1}\right)^4\cdot\left(\sqrt3\right)^8 \)}{\( \left(3\sqrt{5}\right)^2-\left(\sqrt2\right)^6 \)}{\( \left(\sqrt3\right)^4-\left(\sqrt2\right)^4 \)}{\( 3\left(\sqrt{0.1}\right)^4+\left(\sqrt3\right)^8 \)}{}{}}
\LANGpl{
\Question{
Która z podanych liczb należy do przedziału   \( (-5;5) \)?}
{\( 3\left(\sqrt{0{,}1}\right)^4\cdot\left(\sqrt3\right)^8 \)}{\( \left(3\sqrt{5}\right)^2-\left(\sqrt2\right)^6 \)}{\( \left(\sqrt3\right)^4-\left(\sqrt2\right)^4 \)}{\( 3\left(\sqrt{0{,}1}\right)^4+\left(\sqrt3\right)^8 \)}{}{}}
\LANGcs{
\Question{
Které z daných čísel náleží intervalu \( (-5;5) \)?}
{\( 3\left(\sqrt{0{,}1}\right)^4\cdot\left(\sqrt3\right)^8 \)}{\( \left(3\sqrt{5}\right)^2-\left(\sqrt2\right)^6 \)}{\( \left(\sqrt3\right)^4-\left(\sqrt2\right)^4 \)}{\( 3\left(\sqrt{0{,}1}\right)^4+\left(\sqrt3\right)^8 \)}{}{}}
\LANGsk{
\Question{
Ktoré z daných čísel patrí do intervalu \( (-5;5) \)?}
{\( 3\left(\sqrt{0{,}1}\right)^4\cdot\left(\sqrt3\right)^8 \)}{\( \left(3\sqrt{5}\right)^2-\left(\sqrt2\right)^6 \)}{\( \left(\sqrt3\right)^4-\left(\sqrt2\right)^4 \)}{\( 3\left(\sqrt{0{,}1}\right)^4+\left(\sqrt3\right)^8 \)}{}{}}
\LANGes{
\Question{
¿Cuál de los números pertenece al intervalo  \( (-5;5) \)?}
{\( 3\left(\sqrt{0.1}\right)^4\cdot\left(\sqrt3\right)^8 \)}{\( \left(3\sqrt{5}\right)^2-\left(\sqrt2\right)^6 \)}{\( \left(\sqrt3\right)^4-\left(\sqrt2\right)^4 \)}{\( 3\left(\sqrt{0.1}\right)^4+\left(\sqrt3\right)^8 \)}{}{}}
\End

\Data\ID{1003034107}
\Updated{2022-08-25 10:08:00}
\Level{B}
\SubArea{Expressions with powers and roots}
\LANGen{
\Question{
By computing \( 4^{11}\cdot4^{-11}\) we get:}
{\( 1 \)}{\( 0 \)}{\( 4^{22} \)}{\( 16^{-121} \)}{}{}}
\LANGpl{
\Question{
Wartością wyrażenia \( 4^{11}\cdot4^{-11}\) jest:}
{\( 1 \)}{\( 0 \)}{\( 4^{22} \)}{\( 16^{-121} \)}{}{}}
\LANGcs{
\Question{
Zjednodušením výrazu \( 4^{11}\cdot4^{-11}\) dostaneme:}
{\( 1 \)}{\( 0 \)}{\( 4^{22} \)}{\( 16^{-121} \)}{}{}}
\LANGsk{
\Question{
Zjednodušením výrazu \( 4^{11}\cdot4^{-11}\) dostaneme:}
{\( 1 \)}{\( 0 \)}{\( 4^{22} \)}{\( 16^{-121} \)}{}{}}
\LANGes{
\Question{
Calculando \( 4^{11}\cdot4^{-11}\) obtenemos:}
{\( 1 \)}{\( 0 \)}{\( 4^{22} \)}{\( 16^{-121} \)}{}{}}
\End

\Data\ID{1003099408}
\Updated{2020-08-27 09:08:04}
\Level{B}
\SubArea{Expressions with powers and roots}
\LANGen{
\Question{
The value of the expression \( \frac12\cdot\left[\frac{5\cdot\left(0.2+\frac35\right)^2}{3.2}\right]+\frac13 \) is:}
{\( \frac56 \)}{\( \frac32 \)}{\( \frac43 \)}{\( \frac52 \)}{}{}}
\LANGpl{
\Question{
Wartość wyrażenia  \( \frac12\cdot\left[\frac{5\cdot\left(0{,}2+\frac35\right)^2}{3{,}2}\right]+\frac13 \) wynosi:}
{\( \frac56 \)}{\( \frac32 \)}{\( \frac43 \)}{\( \frac52 \)}{}{}}
\LANGcs{
\Question{
Hodnota výrazu \( \frac12\cdot\left[\frac{5\cdot\left(0{,}2+\frac35\right)^2}{3{,}2}\right]+\frac13 \) je:}
{\( \frac56 \)}{\( \frac32 \)}{\( \frac43 \)}{\( \frac52 \)}{}{}}
\LANGsk{
\Question{
Hodnota výrazu \( \frac12\cdot\left[\frac{5\cdot\left(0{,}2+\frac35\right)^2}{3{,}2}\right]+\frac13 \) je:}
{\( \frac56 \)}{\( \frac32 \)}{\( \frac43 \)}{\( \frac52 \)}{}{}}
\LANGes{
\Question{
El valor de la expresión \( \frac12\cdot\left[\frac{5\cdot\left(0.2+\frac35\right)^2}{3.2}\right]+\frac13 \) es:}
{\( \frac56 \)}{\( \frac32 \)}{\( \frac43 \)}{\( \frac52 \)}{}{}}
\End

\Data\ID{1003099409}
\Updated{2020-08-27 09:08:26}
\Level{B}
\SubArea{Expressions with powers and roots}
\LANGen{
\Question{
Simplifying \( \left( \frac1{\left( \sqrt[3]{729}+\sqrt[4]{256}+2 \right)^0} \right)^{-2} \) we get:}
{\( 1 \)}{\( \frac1{15} \)}{\( \frac1{225} \)}{\( 15 \)}{}{}}
\LANGpl{
\Question{
Liczba \( \left( \frac1{\left( \sqrt[3]{729}+\sqrt[4]{256}+2 \right)^0} \right)^{-2} \) jest równa:}
{\( 1 \)}{\( \frac1{15} \)}{\( \frac1{225} \)}{\( 15 \)}{}{}}
\LANGcs{
\Question{
Zjednodušením výrazu \( \left( \frac1{\left( \sqrt[3]{729}+\sqrt[4]{256}+2 \right)^0} \right)^{-2} \) dostaneme:}
{\( 1 \)}{\( \frac1{15} \)}{\( \frac1{225} \)}{\( 15 \)}{}{}}
\LANGsk{
\Question{
Zjednodušením výrazu \( \left( \frac1{\left( \sqrt[3]{729}+\sqrt[4]{256}+2 \right)^0} \right)^{-2} \) dostaneme:}
{\( 1 \)}{\( \frac1{15} \)}{\( \frac1{225} \)}{\( 15 \)}{}{}}
\LANGes{
\Question{
Simplificando \( \left( \frac1{\left( \sqrt[3]{729}+\sqrt[4]{256}+2 \right)^0} \right)^{-2} \) obtenemos:}
{\( 1 \)}{\( \frac1{15} \)}{\( \frac1{225} \)}{\( 15 \)}{}{}}
\End

\Data\ID{1003099410}
\Updated{2022-04-23 05:04:14}
\Level{B}
\SubArea{Expressions with powers and roots}
\LANGen{
\Question{
The value of the multiplicative inverse of \( \left[ 2^{-2}+\left( \frac16 \right)^{-1} \right]^{\frac12} \) is:}
{\( \frac25 \)}{\( \frac12+\sqrt6 \)}{\( \frac4{25} \)}{\( \frac52 \)}{}{}}
\LANGpl{
\Question{
Wartość odwrotna wyrażenia\( \left[ 2^{-2}+\left( \frac16 \right)^{-1} \right]^{\frac12} \) jest równa:}
{\( \frac25 \)}{\( \frac12+\sqrt6 \)}{\( \frac4{25} \)}{\( \frac52 \)}{}{}}
\LANGcs{
\Question{
Převrácená hodnota výrazu \( \left[ 2^{-2}+\left( \frac16 \right)^{-1} \right]^{\frac12} \) je:}
{\( \frac25 \)}{\( \frac12+\sqrt6 \)}{\( \frac4{25} \)}{\( \frac52 \)}{}{}}
\LANGsk{
\Question{
Prevrátená hodnota výrazu \( \left[ 2^{-2}+\left( \frac16 \right)^{-1} \right]^{\frac12} \) je:}
{\( \frac25 \)}{\( \frac12+\sqrt6 \)}{\( \frac4{25} \)}{\( \frac52 \)}{}{}}
\LANGes{
\Question{
Halla el inverso multiplicativo de \( \left[ 2^{-2}+\left( \frac16 \right)^{-1} \right]^{\frac12} \).}
{\( \frac25 \)}{\( \frac12+\sqrt6 \)}{\( \frac4{25} \)}{\( \frac52 \)}{}{}}
\End

\Data\ID{1003099501}
\Updated{2020-08-27 09:08:21}
\Level{B}
\SubArea{Expressions with powers and roots}
\LANGen{
\Question{
Let \( x=4^{-1}+4^{-\frac12}-\left(\frac{\sqrt2}2\right)^2 \). Which of the following inequalities is true?}
{\( x \geq 2^{-2} \)}{\( x < 4^{-1} \)}{\( x > 2 \)}{\( x \leq 4^{-3} \)}{}{}}
\LANGpl{
\Question{
Wiadomo, że \( x=4^{-1}+4^{-\frac12}-\left(\frac{\sqrt2}2\right)^2 \). Która z poniższych nierówności jest prawdziwa?}
{\( x \geq 2^{-2} \)}{\( x < 4^{-1} \)}{\( x > 2 \)}{\( x \leq 4^{-3} \)}{}{}}
\LANGcs{
\Question{
Nechť \( x=4^{-1}+4^{-\frac12}-\left(\frac{\sqrt2}2\right)^2 \). Která z následujících nerovnic je pravdivá?}
{\( x \geq 2^{-2} \)}{\( x < 4^{-1} \)}{\( x > 2 \)}{\( x \leq 4^{-3} \)}{}{}}
\LANGsk{
\Question{
Nech \( x=4^{-1}+4^{-\frac12}-\left(\frac{\sqrt2}2\right)^2 \). Ktorá z nasledujúcich nerovníc je pravdivá?}
{\( x \geq 2^{-2} \)}{\( x < 4^{-1} \)}{\( x > 2 \)}{\( x \leq 4^{-3} \)}{}{}}
\LANGes{
\Question{
Sea \( x=4^{-1}+4^{-\frac12}-\left(\frac{\sqrt2}2\right)^2 \).  ¿Cuál de las inecuaciones siguientes es correcta?}
{\( x \geq 2^{-2} \)}{\( x < 4^{-1} \)}{\( x > 2 \)}{\( x \leq 4^{-3} \)}{}{}}
\End

\Data\ID{1003099502}
\Updated{2022-04-23 05:04:14}
\Level{B}
\SubArea{Expressions with powers and roots}
\LANGen{
\Question{
Simplify the fraction \( \frac{\sqrt[8]9\cdot\sqrt[12]{27}\cdot\sqrt[4]{14}}{\sqrt[4]{42}} \).}
{\( \sqrt[4]3 \)}{\( \frac1{\sqrt[4]3} \)}{\( 1 \)}{\( 3 \)}{}{}}
\LANGpl{
\Question{
Liczba \( \frac{\sqrt[8]9\cdot\sqrt[12]{27}\cdot\sqrt[4]{14}}{\sqrt[4]{42}} \) jest równa:}
{\( \sqrt[4]3 \)}{\( \frac1{\sqrt[4]3} \)}{\( 1 \)}{\( 3 \)}{}{}}
\LANGcs{
\Question{
Zjednodušte zlomek \( \frac{\sqrt[8]9\cdot\sqrt[12]{27}\cdot\sqrt[4]{14}}{\sqrt[4]{42}} \).}
{\( \sqrt[4]3 \)}{\( \frac1{\sqrt[4]3} \)}{\( 1 \)}{\( 3 \)}{}{}}
\LANGsk{
\Question{
Zjednodušte zlomok \( \frac{\sqrt[8]9\cdot\sqrt[12]{27}\cdot\sqrt[4]{14}}{\sqrt[4]{42}} \).}
{\( \sqrt[4]3 \)}{\( \frac1{\sqrt[4]3} \)}{\( 1 \)}{\( 3 \)}{}{}}
\LANGes{
\Question{
Simplifica la fracción \( \frac{\sqrt[8]9\cdot\sqrt[12]{27}\cdot\sqrt[4]{14}}{\sqrt[4]{42}} \).}
{\( \sqrt[4]3 \)}{\( \frac1{\sqrt[4]3} \)}{\( 1 \)}{\( 3 \)}{}{}}
\End

\Data\ID{1003099503}
\Updated{2020-10-12 10:10:34}
\Level{B}
\SubArea{Expressions with powers and roots}
\LANGen{
\Question{
Let \( a = 2\sqrt7 + \sqrt5 \) and \( b=\frac1{\sqrt7-\sqrt5} \).  Choose the right  relation between \( a \) and \( b \).}
{\( a > b \)}{\( a = b \)}{\( a < b \)}{\( a + b = 0 \)}{}{}}
\LANGpl{
\Question{
Dane są liczby \( a = 2\sqrt7 + \sqrt5 \)   i  \( b=\frac1{\sqrt7-\sqrt5} \).  Wybierz właściwą relację między \( a \)  i  \( b \).}
{\( a > b \)}{\( a = b \)}{\( a < b \)}{\( a + b = 0 \)}{}{}}
\LANGcs{
\Question{
Nechť \( a = 2\sqrt7 + \sqrt5 \) a \( b=\frac1{\sqrt7-\sqrt5} \). Vyberte správný vztah mezi \( a \) a \( b \).}
{\( a > b \)}{\( a = b \)}{\( a < b \)}{\( a + b = 0 \)}{}{}}
\LANGsk{
\Question{
Dané sú čísla \( a = 2\sqrt7 + \sqrt5 \) a \( b=\frac1{\sqrt7-\sqrt5} \). Vyberte správny vzťah medzi číslami \( a \) a \( b \).}
{\( a > b \)}{\( a = b \)}{\( a < b \)}{\( a + b = 0 \)}{}{}}
\LANGes{
\Question{
Sean \( a = 2\sqrt7 + \sqrt5 \) y \( b=\frac1{\sqrt7-\sqrt5} \).  Elije la relación correcta entre \( a \) y \( b \).}
{\( a > b \)}{\( a = b \)}{\( a < b \)}{\( a + b = 0 \)}{}{}}
\End

\Data\ID{1003099506}
\Updated{2020-08-27 09:08:30}
\Level{B}
\SubArea{Expressions with powers and roots}
\LANGen{
\Question{
The additive inverse of \( \frac1{5-2\sqrt5} \) is:}
{\( \frac{-5-2\sqrt5}5 \)}{\( \frac{-1}{2\sqrt5-5} \)}{\( \frac{-1}{2\sqrt5+5} \)}{\( 5-2\sqrt5 \)}{}{}}
\LANGpl{
\Question{
Liczbą przeciwną do liczby \( \frac1{5-2\sqrt5} \) jest:}
{\( \frac{-5-2\sqrt5}5 \)}{\( \frac{-1}{2\sqrt5-5} \)}{\( \frac{-1}{2\sqrt5+5} \)}{\( 5-2\sqrt5 \)}{}{}}
\LANGcs{
\Question{
Opačný výraz k výrazu \( \frac1{5-2\sqrt5} \) je:}
{\( \frac{-5-2\sqrt5}5 \)}{\( \frac{-1}{2\sqrt5-5} \)}{\( \frac{-1}{2\sqrt5+5} \)}{\( 5-2\sqrt5 \)}{}{}}
\LANGsk{
\Question{
Opačný výraz k výrazu \( \frac1{5-2\sqrt5} \) je:}
{\( \frac{-5-2\sqrt5}5 \)}{\( \frac{-1}{2\sqrt5-5} \)}{\( \frac{-1}{2\sqrt5+5} \)}{\( 5-2\sqrt5 \)}{}{}}
\LANGes{
\Question{
El inverso aditivo de \( \frac1{5-2\sqrt5} \) es:}
{\( \frac{-5-2\sqrt5}5 \)}{\( \frac{-1}{2\sqrt5-5} \)}{\( \frac{-1}{2\sqrt5+5} \)}{\( 5-2\sqrt5 \)}{}{}}
\End

\Data\ID{1003099507}
\Updated{2020-08-27 09:08:59}
\Level{B}
\SubArea{Expressions with powers and roots}
\LANGen{
\Question{
The multiplicative inverse of \( \frac2{\sqrt3-1} \) is:}
{\( \frac1{\sqrt3+1} \)}{\( \frac2{\sqrt3+1} \)}{\( \frac{-2}{\sqrt3-1} \)}{\( \frac{1-\sqrt3}2 \)}{}{}}
\LANGpl{
\Question{
Odwrotnością liczby \( \frac2{\sqrt3-1} \) jest:}
{\( \frac1{\sqrt3+1} \)}{\( \frac2{\sqrt3+1} \)}{\( \frac{-2}{\sqrt3-1} \)}{\( \frac{1-\sqrt3}2 \)}{}{}}
\LANGcs{
\Question{
Převrácená hodnota  výrazu \( \frac2{\sqrt3-1} \) je:}
{\( \frac1{\sqrt3+1} \)}{\( \frac2{\sqrt3+1} \)}{\( \frac{-2}{\sqrt3-1} \)}{\( \frac{1-\sqrt3}2 \)}{}{}}
\LANGsk{
\Question{
Prevrátená hodnota daného výrazu \( \frac2{\sqrt3-1} \) je:}
{\( \frac1{\sqrt3+1} \)}{\( \frac2{\sqrt3+1} \)}{\( \frac{-2}{\sqrt3-1} \)}{\( \frac{1-\sqrt3}2 \)}{}{}}
\LANGes{
\Question{
El inverso multiplicativo de \( \frac2{\sqrt3-1} \) es:}
{\( \frac1{\sqrt3+1} \)}{\( \frac2{\sqrt3+1} \)}{\( \frac{-2}{\sqrt3-1} \)}{\( \frac{1-\sqrt3}2 \)}{}{}}
\End

\Data\ID{1003099510}
\Updated{2020-10-12 10:10:08}
\Level{B}
\SubArea{Expressions with powers and roots}
\LANGen{
\Question{
What is the value of \( \sqrt[4]{2\sqrt2}\sqrt[8]{32} \)?}
{\( 2 \)}{\( 2^{\frac12} \)}{\( 2^{\frac34} \)}{\( 2^0 \)}{}{}}
\LANGpl{
\Question{
Liczba \( \sqrt[4]{2\sqrt2}\sqrt[8]{32} \) jest równa:}
{\( 2 \)}{\( 2^{\frac12} \)}{\( 2^{\frac34} \)}{\( 2^0 \)}{}{}}
\LANGcs{
\Question{
Jaká je hodnota výrazu \( \sqrt[4]{2\sqrt2}\sqrt[8]{32} \)?}
{\( 2 \)}{\( 2^{\frac12} \)}{\( 2^{\frac34} \)}{\( 2^0 \)}{}{}}
\LANGsk{
\Question{
Určte hodnotu výrazu \( \sqrt[4]{2\sqrt2}\sqrt[8]{32} \).}
{\( 2 \)}{\( 2^{\frac12} \)}{\( 2^{\frac34} \)}{\( 2^0 \)}{}{}}
\LANGes{
\Question{
¿Cuál es el valor de \( \sqrt[4]{2\sqrt2}\sqrt[8]{32} \)?}
{\( 2 \)}{\( 2^{\frac12} \)}{\( 2^{\frac34} \)}{\( 2^0 \)}{}{}}
\End

\Data\ID{1003099605}
\Updated{2020-08-27 08:08:20}
\Level{B}
\SubArea{Expressions with powers and roots}
\LANGen{
\Question{
Simplifying \( \left( \sqrt[3]{3\sqrt9} \right)^{\frac32} \sqrt{9^{-1}} \) we get:}
{\( 1 \)}{\( 3\sqrt[6]3 \)}{\( 3\sqrt[3]3 \)}{\( 3 \)}{}{}}
\LANGpl{
\Question{
Liczba \( \left( \sqrt[3]{3\sqrt9} \right)^{\frac32} \sqrt{9^{-1}} \)  jest równa:}
{\( 1 \)}{\( 3\sqrt[6]3 \)}{\( 3\sqrt[3]3 \)}{\( 3 \)}{}{}}
\LANGcs{
\Question{
Zjednodušením \( \left( \sqrt[3]{3\sqrt9} \right)^{\frac32} \sqrt{9^{-1}} \) dostaneme:}
{\( 1 \)}{\( 3\sqrt[6]3 \)}{\( 3\sqrt[3]3 \)}{\( 3 \)}{}{}}
\LANGsk{
\Question{
Zjednodušením výrazu \( \left( \sqrt[3]{3\sqrt9} \right)^{\frac32} \sqrt{9^{-1}} \) dostaneme:}
{\( 1 \)}{\( 3\sqrt[6]3 \)}{\( 3\sqrt[3]3 \)}{\( 3 \)}{}{}}
\LANGes{
\Question{
Simplificando \( \left( \sqrt[3]{3\sqrt9} \right)^{\frac32} \sqrt{9^{-1}} \) obtenemos:}
{\( 1 \)}{\( 3\sqrt[6]3 \)}{\( 3\sqrt[3]3 \)}{\( 3 \)}{}{}}
\End

\Data\ID{1003099606}
\Updated{2020-10-12 10:10:26}
\Level{B}
\SubArea{Expressions with powers and roots}
\LANGen{
\Question{
Calculate the value of the expression \( \frac{2a+12}{-a^2} \)  for \( a=-2\sqrt3 \).}
{\( \frac{\sqrt3-3}3 \)}{\( 4\sqrt3 -1 \)}{\( \frac{-\sqrt3+3}3 \)}{\( -4\sqrt3+1 \)}{}{}}
\LANGpl{
\Question{
Wartość wyrażenia  \( \frac{2a+12}{-a^2} \)  dla \( a=-2\sqrt3 \)   jest równa:}
{\( \frac{\sqrt3-3}3 \)}{\( 4\sqrt3 -1 \)}{\( \frac{-\sqrt3+3}3 \)}{\( -4\sqrt3+1 \)}{}{}}
\LANGcs{
\Question{
Vypočítejte hodnotu výrazu \( \frac{2a+12}{-a^2} \)  pro \( a=-2\sqrt3 \).}
{\( \frac{\sqrt3-3}3 \)}{\( 4\sqrt3 -1 \)}{\( \frac{-\sqrt3+3}3 \)}{\( -4\sqrt3+1 \)}{}{}}
\LANGsk{
\Question{
Vypočítajte hodnotu výrazu \( \frac{2a+12}{-a^2} \) pre \( a=-2\sqrt3 \).}
{\( \frac{\sqrt3-3}3 \)}{\( 4\sqrt3 -1 \)}{\( \frac{-\sqrt3+3}3 \)}{\( -4\sqrt3+1 \)}{}{}}
\LANGes{
\Question{
Calcula el valor de la expresión \( \frac{2a+12}{-a^2} \)  para \( a=-2\sqrt3 \).}
{\( \frac{\sqrt3-3}3 \)}{\( 4\sqrt3 -1 \)}{\( \frac{-\sqrt3+3}3 \)}{\( -4\sqrt3+1 \)}{}{}}
\End

\Data\ID{1003099608}
\Updated{2020-08-27 08:08:22}
\Level{B}
\SubArea{Expressions with powers and roots}
\LANGen{
\Question{
Choose the correct statement about the number \( 4\sqrt2-\frac{2\sqrt2+2}{\sqrt2-1} \).}
{It is a rational number.}{It is an irrational number.}{It is greater than \( \sqrt2 \).}{It is a positive integer.}{}{}}
\LANGpl{
\Question{
Liczba \( 4\sqrt2-\frac{2\sqrt2+2}{\sqrt2-1} \) jest liczbą:}
{wymierną.}{niewymierną.}{większą niż \( \sqrt2 \).}{dodatnią całkowitą.}{}{}}
\LANGcs{
\Question{
Vyber správné tvrzení o číslu \( 4\sqrt2-\frac{2\sqrt2+2}{\sqrt2-1} \).}
{Je racionální číslo.}{Je iracionální číslo.}{Je větší než \( \sqrt2 \).}{Je kladné celé číslo.}{}{}}
\LANGsk{
\Question{
Vyberte pravdivé tvrdenie o čísle \( 4\sqrt2-\frac{2\sqrt2+2}{\sqrt2-1} \).}
{Je racionálne číslo.}{Je iracionálne číslo.}{Je väčšie ako \( \sqrt2 \).}{Je kladné celé číslo.}{}{}}
\LANGes{
\Question{
Elije la expresión correcta sobre el número \( 4\sqrt2-\frac{2\sqrt2+2}{\sqrt2-1} \).}
{Es un número racional.}{Es un número irracional.}{Es mayor que \( \sqrt2 \).}{Es un número natural.}{}{}}
\End

\Data\ID{1003099609}
\Updated{2020-10-12 10:10:43}
\Level{B}
\SubArea{Expressions with powers and roots}
\LANGen{
\Question{
Complete the sentence to get a true statement: The numbers \( -\frac{\sqrt3}6-\frac12 \) and \( \sqrt3-3 \) are ...}
{multiplicative inverses of each other.}{equal.}{rational numbers.}{additive inverses of each other.}{}{}}
\LANGpl{
\Question{
Uzupełnij zdanie, aby otrzymać zdanie prawdziwe. Liczby \( -\frac{\sqrt3}6-\frac12 \)  i  \( \sqrt3-3 \)  są liczbami...}
{odwrotnymi.}{równymi.}{wymiernymi.}{przeciwnymi.}{}{}}
\LANGcs{
\Question{
Doplňte větu tak, aby tvrzení bylo pravdivé: Čísla \( -\frac{\sqrt3}6-\frac12 \) a \( \sqrt3-3 \) jsou ...}
{převrácené výrazy jeden druhého.}{si rovna.}{racionální čísla.}{opačné výrazy jeden druhého.}{}{}}
\LANGsk{
\Question{
Doplňte vetu tak, aby bolo tvrdenie pravdivé: Čísla \( -\frac{\sqrt3}6-\frac12 \) a \( \sqrt3-3 \) sú ...}
{prevrátené výrazy jeden druhého.}{rovné.}{racionálne čísla.}{opačné výrazy jeden druhého.}{}{}}
\LANGes{
\Question{
Completa la frase de manera que sea verdadera: Los números \( -\frac{\sqrt3}6-\frac12 \) y \( \sqrt3-3 \) son ...}
{inversos multiplicativos.}{iguales.}{números racionales.}{inversos aditivos.}{}{}}
\End

\Data\ID{1003118005}
\Updated{2020-10-12 10:10:04}
\Level{B}
\SubArea{Expressions with powers and roots}
\LANGen{
\Question{
Determine the numbers \( a \) and \( b \) if you know that \( \left(\sqrt5-3\right)\left(a\sqrt5+b\right)=-9\sqrt5+5\sqrt5 \).}
{\( a=3\text{, }b=5 \)}{\( a=\sqrt5\text{, }b=3 \)}{\( a=-3\text{, }b=1 \)}{\( a=5\text{, }b=\sqrt5 \)}{}{}}
\LANGpl{
\Question{
Wiadomo, że  \( a \)  i  \( b \) są liczbami wymiernymi oraz  \( \left(\sqrt5-3\right)\left(a\sqrt5+b\right)=-9\sqrt5+5\sqrt5 \) zatem:}
{\( a=3\text{, }b=5 \)}{\( a=\sqrt5\text{, }b=3 \)}{\( a=-3\text{, }b=1 \)}{\( a=5\text{, }b=\sqrt5 \)}{}{}}
\LANGcs{
\Question{
Určete čísla \( a \) a \( b \) jestliže víte, že \( \left(\sqrt5-3\right)\left(a\sqrt5+b\right)=-9\sqrt5+5\sqrt5 \).}
{\( a=3\text{, }b=5 \)}{\( a=\sqrt5\text{, }b=3 \)}{\( a=-3\text{, }b=1 \)}{\( a=5\text{, }b=\sqrt5 \)}{}{}}
\LANGsk{
\Question{
Určte čísla \( a \) a \( b \) ak viete, že platí \( \left(\sqrt5-3\right)\left(a\sqrt5+b\right)=-9\sqrt5+5\sqrt5 \).}
{\( a=3\text{, }b=5 \)}{\( a=\sqrt5\text{, }b=3 \)}{\( a=-3\text{, }b=1 \)}{\( a=5\text{, }b=\sqrt5 \)}{}{}}
\LANGes{
\Question{
Determina los números  \( a \) y \( b \) si sabes que  \( \left(\sqrt5-3\right)\left(a\sqrt5+b\right)=-9\sqrt5+5\sqrt5 \).}
{\( a=3\text{, }b=5 \)}{\( a=\sqrt5\text{, }b=3 \)}{\( a=-3\text{, }b=1 \)}{\( a=5\text{, }b=\sqrt5 \)}{}{}}
\End

\Data\ID{1003118009}
\Updated{2020-08-27 09:08:41}
\Level{B}
\SubArea{Expressions with powers and roots}
\LANGen{
\Question{
The square of the number  \( \sqrt2-\sqrt[4]2 \)   is equal to:}
{\( 2-2\sqrt[4]8+\sqrt2 \)}{\( 2-2\sqrt[4]2+\sqrt2 \)}{\( 2-2\sqrt[16]8+\sqrt[8]2 \)}{\( 2-\sqrt2 \)}{}{}}
\LANGpl{
\Question{
Kwadrat liczby \( \sqrt2-\sqrt[4]2 \)   jest równy:}
{\( 2-2\sqrt[4]8+\sqrt2 \)}{\( 2-2\sqrt[4]2+\sqrt2 \)}{\( 2-2\sqrt[16]8+\sqrt[8]2 \)}{\( 2-\sqrt2 \)}{}{}}
\LANGcs{
\Question{
Druhá mocnina čísla \( \sqrt2-\sqrt[4]2 \)   se rovná:}
{\( 2-2\sqrt[4]8+\sqrt2 \)}{\( 2-2\sqrt[4]2+\sqrt2 \)}{\( 2-2\sqrt[16]8+\sqrt[8]2 \)}{\( 2-\sqrt2 \)}{}{}}
\LANGsk{
\Question{
Druhá mocnina čísla  \( \sqrt2-\sqrt[4]2 \)  sa rovná:}
{\( 2-2\sqrt[4]8+\sqrt2 \)}{\( 2-2\sqrt[4]2+\sqrt2 \)}{\( 2-2\sqrt[16]8+\sqrt[8]2 \)}{\( 2-\sqrt2 \)}{}{}}
\LANGes{
\Question{
El cuadrado del número  \( \sqrt2-\sqrt[4]2 \)   es igual a:}
{\( 2-2\sqrt[4]8+\sqrt2 \)}{\( 2-2\sqrt[4]2+\sqrt2 \)}{\( 2-2\sqrt[16]8+\sqrt[8]2 \)}{\( 2-\sqrt2 \)}{}{}}
\End

\Data\ID{1003118103}
\Updated{2020-08-27 09:08:09}
\Level{B}
\SubArea{Expressions with powers and roots}
\LANGen{
\Question{
Express \( \frac{\sqrt[3]5\sqrt[6]5}{3\cdot25^3+2\cdot125^2} \) as a power of \( 5 \).}
{\( 5^{-\frac{13}2} \)}{\( 5^{-5} \)}{\( 5^{-6} \)}{\( 5^{-\frac{11}2} \)}{}{}}
\LANGpl{
\Question{
Liczba \( \frac{\sqrt[3]5\sqrt[6]5}{3\cdot25^3+2\cdot125^2} \) jako potęga  \( 5 \) jest równa:}
{\( 5^{-\frac{13}2} \)}{\( 5^{-5} \)}{\( 5^{-6} \)}{\( 5^{-\frac{11}2} \)}{}{}}
\LANGcs{
\Question{
Vyjádřete \( \frac{\sqrt[3]5\sqrt[6]5}{3\cdot25^3+2\cdot125^2} \) jako mocninu \( 5 \).}
{\( 5^{-\frac{13}2} \)}{\( 5^{-5} \)}{\( 5^{-6} \)}{\( 5^{-\frac{11}2} \)}{}{}}
\LANGsk{
\Question{
Vyjadrite \( \frac{\sqrt[3]5\sqrt[6]5}{3\cdot25^3+2\cdot125^2} \) ako mocninu \( 5 \).}
{\( 5^{-\frac{13}2} \)}{\( 5^{-5} \)}{\( 5^{-6} \)}{\( 5^{-\frac{11}2} \)}{}{}}
\LANGes{
\Question{
Expresa \( \frac{\sqrt[3]5\sqrt[6]5}{3\cdot25^3+2\cdot125^2} \) como potencia de \( 5 \).}
{\( 5^{-\frac{13}2} \)}{\( 5^{-5} \)}{\( 5^{-6} \)}{\( 5^{-\frac{11}2} \)}{}{}}
\End

\Data\ID{1003118107}
\Updated{2022-04-23 05:04:14}
\Level{B}
\SubArea{Expressions with powers and roots}
\LANGen{
\Question{
The number \( \left(\frac{27^{-4}\cdot8^{-4}}{16^{-2}\cdot9^{-5}}\right)^{-3} \) equals:}
{\( 12^{6} \)}{\( 6^6 \)}{\( 6^{12} \)}{\( \frac1{3^6\cdot2^{12}} \)}{}{}}
\LANGpl{
\Question{
Liczba  \( \left(\frac{27^{-4}\cdot8^{-4}}{16^{-2}\cdot9^{-5}}\right)^{-3} \) jest równa:}
{\( 12^{6} \)}{\( 6^6 \)}{\( 6^{12} \)}{\( \frac1{3^6\cdot2^{12}} \)}{}{}}
\LANGcs{
\Question{
Číslo \( \left(\frac{27^{-4}\cdot8^{-4}}{16^{-2}\cdot9^{-5}}\right)^{-3} \) se rovná:}
{\( 12^{6} \)}{\( 6^6 \)}{\( 6^{12} \)}{\( \frac1{3^6\cdot2^{12}} \)}{}{}}
\LANGsk{
\Question{
Číslo \( \left(\frac{27^{-4}\cdot8^{-4}}{16^{-2}\cdot9^{-5}}\right)^{-3} \) sa rovná:}
{\( 12^{6} \)}{\( 6^6 \)}{\( 6^{12} \)}{\( \frac1{3^6\cdot2^{12}} \)}{}{}}
\LANGes{
\Question{
El número \( \left(\frac{27^{-4}\cdot8^{-4}}{16^{-2}\cdot9^{-5}}\right)^{-3} \) es igual a:}
{\( 12^{6} \)}{\( 6^6 \)}{\( 6^{12} \)}{\( \frac1{3^6\cdot2^{12}} \)}{}{}}
\End

\Data\ID{1003118603}
\Updated{2020-08-27 02:08:58}
\Level{B}
\SubArea{Expressions with powers and roots}
\LANGen{
\Question{
Select the number that equals \( \sqrt[5]{64} \).}
{\( 2\sqrt[5]2 \)}{\( \frac{\sqrt[5]2}2 \)}{\( \sqrt2 \)}{\( 2 \)}{}{}}
\LANGpl{
\Question{
Wskaż liczbę równą liczbie  \( \sqrt[5]{64} \).}
{\( 2\sqrt[5]2 \)}{\( \frac{\sqrt[5]2}2 \)}{\( \sqrt2 \)}{\( 2 \)}{}{}}
\LANGcs{
\Question{
Které z uvedených čísel je rovno \( \sqrt[5]{64} \)?}
{\( 2\sqrt[5]2 \)}{\( \frac{\sqrt[5]2}2 \)}{\( \sqrt2 \)}{\( 2 \)}{}{}}
\LANGsk{
\Question{
Vyberte číslo rovné \( \sqrt[5]{64} \).}
{\( 2\sqrt[5]2 \)}{\( \frac{\sqrt[5]2}2 \)}{\( \sqrt2 \)}{\( 2 \)}{}{}}
\LANGes{
\Question{
¿Cuál de los números es igual a \( \sqrt[5]{64} \)?}
{\( 2\sqrt[5]2 \)}{\( \frac{\sqrt[5]2}2 \)}{\( \sqrt2 \)}{\( 2 \)}{}{}}
\End

\Data\ID{1003118605}
\Updated{2020-10-12 10:10:16}
\Level{B}
\SubArea{Expressions with powers and roots}
\LANGen{
\Question{
If \( \frac{2\cdot\sqrt[3]2}{\sqrt8} = 2^x \), then}
{\( x=-\frac16 \).}{\( x=0 \).}{\( x=\frac13 \).}{\( x=-4 \).}{}{}}
\LANGpl{
\Question{
Jeżeli  \( \frac{2\cdot\sqrt[3]2}{\sqrt8} = 2^x \), to}
{\( x=-\frac16 \).}{\( x=0 \).}{\( x=\frac13 \).}{\( x=-4 \).}{}{}}
\LANGcs{
\Question{
Jestliže \( \frac{2\cdot\sqrt[3]2}{\sqrt8} = 2^x \), pak}
{\( x=-\frac16 \).}{\( x=0 \).}{\( x=\frac13 \).}{\( x=-4 \).}{}{}}
\LANGsk{
\Question{
Ak \( \frac{2\cdot\sqrt[3]2}{\sqrt8} = 2^x \), tak}
{\( x=-\frac16 \).}{\( x=0 \).}{\( x=\frac13 \).}{\( x=-4 \).}{}{}}
\LANGes{
\Question{
Si \( \frac{2\cdot\sqrt[3]2}{\sqrt8} = 2^x \), entonces}
{\( x=-\frac16 \).}{\( x=0 \).}{\( x=\frac13 \).}{\( x=-4 \).}{}{}}
\End

\Data\ID{1003118607}
\Updated{2020-08-27 02:08:46}
\Level{B}
\SubArea{Expressions with powers and roots}
\LANGen{
\Question{
Which of the following numbers are ordered from least to greatest?}
{\( (0.3)^4 \),  \( 0.027 \),  \( (0.3)^{\sqrt2}  \)}{\( 81^{\frac34} \), \( 16^{\frac14} \), \( 7^{-2} \)}{\( \left(\frac23 \right)^{1.4} \), \( \left(\frac23 \right)^{\pi} \), \( \left(\frac32 \right)^{-1} \)}{\( 7^0 \), \( 7^{-1} \), \( 7^{-2} \)}{}{}}
\LANGpl{
\Question{
W którym z podanych zestawów liczby zapisano w kolejności od najmniejszej do największej?}
{\( (0{,}3)^4 \),  \( 0{,}027 \),  \( (0{,}3)^{\sqrt2}  \)}{\( 81^{\frac34} \), \( 16^{\frac14} \), \( 7^{-2} \)}{\( \left(\frac23 \right)^{1{,}4} \), \( \left(\frac23 \right)^{\pi} \), \( \left(\frac32 \right)^{-1} \)}{\( 7^0 \), \( 7^{-1} \), \( 7^{-2} \)}{}{}}
\LANGcs{
\Question{
Která z uvedených čísel jsou seřazená od nejmenšího k největšímu?}
{\( (0{,}3)^4 \),  \( 0{,}027 \),  \( (0{,}3)^{\sqrt2}  \)}{\( 81^{\frac34} \), \( 16^{\frac14} \), \( 7^{-2} \)}{\( \left(\frac23 \right)^{1{,}4} \), \( \left(\frac23 \right)^{\pi} \), \( \left(\frac32 \right)^{-1} \)}{\( 7^0 \), \( 7^{-1} \), \( 7^{-2} \)}{}{}}
\LANGsk{
\Question{
Ktoré z nasledujúcich čísel sú usporiadané od najmenšieho po najväčšie?}
{\( (0{,}3)^4 \),  \( 0{,}027 \),  \( (0{,}3)^{\sqrt2}  \)}{\( 81^{\frac34} \), \( 16^{\frac14} \), \( 7^{-2} \)}{\( \left(\frac23 \right)^{1{,}4} \), \( \left(\frac23 \right)^{\pi} \), \( \left(\frac32 \right)^{-1} \)}{\( 7^0 \), \( 7^{-1} \), \( 7^{-2} \)}{}{}}
\LANGes{
\Question{
¿Cuáles de los siguientes números están ordenados de menor a mayor?}
{\( (0.3)^4 \),  \( 0.027 \),  \( (0.3)^{\sqrt2}  \)}{\( 81^{\frac34} \), \( 16^{\frac14} \), \( 7^{-2} \)}{\( \left(\frac23 \right)^{1.4} \), \( \left(\frac23 \right)^{\pi} \), \( \left(\frac32 \right)^{-1} \)}{\( 7^0 \), \( 7^{-1} \), \( 7^{-2} \)}{}{}}
\End

\Data\ID{1003118608}
\Updated{2022-04-23 05:04:14}
\Level{B}
\SubArea{Expressions with powers and roots}
\LANGen{
\Question{
Express the value of the expression \( \left(\frac23-2^{-2}\right)^{-1} \) as a decimal number.}
{\( 2.4 \)}{\( 0.41\overline6\dots\)}{\( \frac{12}5 \)}{\( -1.\overline3 \)}{}{}}
\LANGpl{
\Question{
Oblicz wartość wyrażenia  \( \left(\frac23-2^{-2}\right)^{-1} \).  Wynik zapisz w postaci liczby dziesiętnej.}
{\( 2{,}4 \)}{\( 0{,}41\overline6\dots\)}{\( \frac{12}5 \)}{\( -1{,}\overline3 \)}{}{}}
\LANGcs{
\Question{
Vyjádřete hodnotu výrazu \( \left(\frac23-2^{-2}\right)^{-1} \) jako desetinné číslo.}
{\( 2{,}4 \)}{\( 0{,}41\overline6\dots\)}{\( \frac{12}5 \)}{\( -1{,}\overline3 \)}{}{}}
\LANGsk{
\Question{
Vyjadrite hodnotu výrazu \( \left(\frac23-2^{-2}\right)^{-1} \) ako desatinné číslo.}
{\( 2{,}4 \)}{\( 0{,}41\overline6\dots\)}{\( \frac{12}5 \)}{\( -1{,}\overline3 \)}{}{}}
\LANGes{
\Question{
Escribe el valor de la expresión \( \left(\frac23-2^{-2}\right)^{-1} \) como un número decimal.}
{\( 2.4 \)}{\( 0.41\overline6\dots\)}{\( \frac{12}5 \)}{\( -1.\overline3 \)}{}{}}
\End

\Data\ID{1003124901}
\Updated{2020-08-27 02:08:16}
\Level{B}
\SubArea{Expressions with powers and roots}
\LANGen{
\Question{
The fraction \( \frac{1+\sqrt3}{3+\sqrt{11}} \) is equal to:}
{\( \frac{\sqrt{11}-3}{\sqrt3-1}\)}{\( 9 \)}{\( \frac{\sqrt{11}-3}{1-\sqrt3} \)}{\( \frac{\sqrt{11}+2\sqrt3}2 \)}{}{}}
\LANGpl{
\Question{
Liczba \( \frac{1+\sqrt3}{3+\sqrt{11}} \) jest równa liczbie:}
{\( \frac{\sqrt{11}-3}{\sqrt3-1}\)}{\( 9 \)}{\( \frac{\sqrt{11}-3}{1-\sqrt3} \)}{\( \frac{\sqrt{11}+2\sqrt3}2 \)}{}{}}
\LANGcs{
\Question{
Zlomek \( \frac{1+\sqrt3}{3+\sqrt{11}} \) je roven:}
{\( \frac{\sqrt{11}-3}{\sqrt3-1}\)}{\( 9 \)}{\( \frac{\sqrt{11}-3}{1-\sqrt3} \)}{\( \frac{\sqrt{11}+2\sqrt3}2 \)}{}{}}
\LANGsk{
\Question{
Zlomok \( \frac{1+\sqrt3}{3+\sqrt{11}} \) je rovný:}
{\( \frac{\sqrt{11}-3}{\sqrt3-1}\)}{\( 9 \)}{\( \frac{\sqrt{11}-3}{1-\sqrt3} \)}{\( \frac{\sqrt{11}+2\sqrt3}2 \)}{}{}}
\LANGes{
\Question{
La fracción \( \frac{1+\sqrt3}{3+\sqrt{11}} \) es igual a:}
{\( \frac{\sqrt{11}-3}{\sqrt3-1}\)}{\( 9 \)}{\( \frac{\sqrt{11}-3}{1-\sqrt3} \)}{\( \frac{\sqrt{11}+2\sqrt3}2 \)}{}{}}
\End

\Data\ID{1003124907}
\Updated{2022-04-23 05:04:14}
\Level{B}
\SubArea{Expressions with powers and roots}
\LANGen{
\Question{
The multiplicative inverse of the number \( \frac{\sqrt[3]{27^2}:9^{\frac12}}{\sqrt[3]9} \)  is:}
{\( 3^{-\frac13} \)}{\( 3^{\frac23} \)}{\( 3^{\frac13} \)}{\( 3^{-\frac23} \)}{}{}}
\LANGpl{
\Question{
Liczbą odwrotną do liczby  \(\frac{\sqrt[3]{27^2}:9^{\frac12}}{\sqrt[3]9} \)  jest:}
{\( 3^{-\frac13} \)}{\( 3^{\frac23} \)}{\( 3^{\frac13} \)}{\( 3^{-\frac23} \)}{}{}}
\LANGcs{
\Question{
Převrácená hodnota čísla \(\frac{\sqrt[3]{27^2}:9^{\frac12}}{\sqrt[3]9} \)  je:}
{\( 3^{-\frac13} \)}{\( 3^{\frac23} \)}{\( 3^{\frac13} \)}{\( 3^{-\frac23} \)}{}{}}
\LANGsk{
\Question{
Prevrátená hodnota čísla \( \frac{\sqrt[3]{27^2}:9^{\frac12}}{\sqrt[3]9} \)  je:}
{\( 3^{-\frac13} \)}{\( 3^{\frac23} \)}{\( 3^{\frac13} \)}{\( 3^{-\frac23} \)}{}{}}
\LANGes{
\Question{
El inverso multiplicativo del número \(\frac{\sqrt[3]{27^2}:9^{\frac12}}{\sqrt[3]9} \)  es:}
{\( 3^{-\frac13} \)}{\( 3^{\frac23} \)}{\( 3^{\frac13} \)}{\( 3^{-\frac23} \)}{}{}}
\End

\Data\ID{1003124908}
\Updated{2020-08-28 07:08:18}
\Level{B}
\SubArea{Expressions with powers and roots}
\LANGen{
\Question{
Half the inverse of the cube of the number  \( 8^{19} \) is:}
{\( 4^{-86} \)}{\( 2^{170} \)}{\( \frac1{8^{57}} \)}{\( \frac1{2^{170}} \)}{}{}}
\LANGpl{
\Question{
Połową odwrotności sześcianu liczby   \( 8^{19} \) jest:}
{\( 4^{-86} \)}{\( 2^{170} \)}{\( \frac1{8^{57}} \)}{\( \frac1{2^{170}} \)}{}{}}
\LANGcs{
\Question{
Polovina převrácené hodnoty třetí mocniny čísla \( 8^{19} \) je:}
{\( 4^{-86} \)}{\( 2^{170} \)}{\( \frac1{8^{57}} \)}{\( \frac1{2^{170}} \)}{}{}}
\LANGsk{
\Question{
Polovica prevrátenej hodnoty tretej mocniny čísla \( 8^{19} \) je:}
{\( 4^{-86} \)}{\( 2^{170} \)}{\( \frac1{8^{57}} \)}{\( \frac1{2^{170}} \)}{}{}}
\LANGes{
\Question{
La mitad del inverso de la potencia cúbica del número \( 8^{19} \) es:}
{\( 4^{-86} \)}{\( 2^{170} \)}{\( \frac1{8^{57}} \)}{\( \frac1{2^{170}} \)}{}{}}
\End

\Data\ID{1003124909}
\Updated{2022-04-23 05:04:14}
\Level{B}
\SubArea{Expressions with powers and roots}
\LANGen{
\Question{
The number \( \frac1{2^{2015}}\cdot(0.0005)^{2015} \) equals:}
{\( (0.00025)^{2015} \)}{\( \frac1{2000^{2015}} \)}{\( (0.001)^{2015} \)}{\( (0.0025)^{2015} \)}{}{}}
\LANGpl{
\Question{
Liczba \( \frac1{2^{2015}}\cdot(0{,}0005)^{2015} \) jest równa:}
{\( (0{,}00025)^{2015} \)}{\( \frac1{2000^{2015}} \)}{\( (0{,}001)^{2015} \)}{\( (0{,}0025)^{2015} \)}{}{}}
\LANGcs{
\Question{
Číslo \( \frac1{2^{2015}}\cdot(0{,}0005)^{2015} \) je rovno:}
{\( (0{,}00025)^{2015} \)}{\( \frac1{2000^{2015}} \)}{\( (0{,}001)^{2015} \)}{\( (0{,}0025)^{2015} \)}{}{}}
\LANGsk{
\Question{
Číslo \( \frac1{2^{2015}}\cdot(0{,}0005)^{2015} \) sa rovná:}
{\( (0{,}00025)^{2015} \)}{\( \frac1{2000^{2015}} \)}{\( (0{,}001)^{2015} \)}{\( (0{,}0025)^{2015} \)}{}{}}
\LANGes{
\Question{
El número \( \frac1{2^{2015}}\cdot(0.0005)^{2015} \) es igual a:}
{\( (0.00025)^{2015} \)}{\( \frac1{2000^{2015}} \)}{\( (0.001)^{2015} \)}{\( (0.0025)^{2015} \)}{}{}}
\End

\Data\ID{1003134501}
\Updated{2020-10-12 10:10:11}
\Level{B}
\SubArea{Expressions with powers and roots}
\LANGen{
\Question{
The number \( 4^{100}\cdot\left(\frac12\right)^{-500}\cdot\sqrt[3]{8^{-300}} \) is equal to the number \( 2^p \). It means that:}
{\( p=400 \)}{\( p=300 \)}{\( p=800 \)}{\( p = 700 \)}{}{}}
\LANGpl{
\Question{
Liczba \( 4^{100}\cdot\left(\frac12\right)^{-500}\cdot\sqrt[3]{8^{-300}} \) jest równa liczbie \( 2^p \). Wynika z tego, że:}
{\( p=400 \)}{\( p=300 \)}{\( p=800 \)}{\( p = 700 \)}{}{}}
\LANGcs{
\Question{
Číslo \( 4^{100}\cdot\left(\frac12\right)^{-500}\cdot\sqrt[3]{8^{-300}} \) je rovno číslu \( 2^p \). Platí tedy, že:}
{\( p=400 \)}{\( p=300 \)}{\( p=800 \)}{\( p = 700 \)}{}{}}
\LANGsk{
\Question{
Číslo \( 4^{100}\cdot\left(\frac12\right)^{-500}\cdot\sqrt[3]{8^{-300}} \) je rovné číslu \( 2^p \) pre:}
{\( p=400 \)}{\( p=300 \)}{\( p=800 \)}{\( p = 700 \)}{}{}}
\LANGes{
\Question{
El número \( 4^{100}\cdot\left(\frac12\right)^{-500}\cdot\sqrt[3]{8^{-300}} \) es igual a \( 2^p \). Luego la p vale:}
{\( p=400 \)}{\( p=300 \)}{\( p=800 \)}{\( p = 700 \)}{}{}}
\End

\Data\ID{1003134502}
\Updated{2020-08-28 08:08:32}
\Level{B}
\SubArea{Expressions with powers and roots}
\LANGen{
\Question{
The last digit of the number \( 7^{2015} \) is:}
{\( 3 \)}{\( 7 \)}{\( 9 \)}{\( 1 \)}{}{}}
\LANGpl{
\Question{
Ostatnią cyfrą liczby  \( 7^{2015} \) jest:}
{\( 3 \)}{\( 7 \)}{\( 9 \)}{\( 1 \)}{}{}}
\LANGcs{
\Question{
Poslední číslice čísla \( 7^{2015} \) je:}
{\( 3 \)}{\( 7 \)}{\( 9 \)}{\( 1 \)}{}{}}
\LANGsk{
\Question{
Posledná číslica čísla \( 7^{2015} \) je:}
{\( 3 \)}{\( 7 \)}{\( 9 \)}{\( 1 \)}{}{}}
\LANGes{
\Question{
La última cifra del número \( 7^{2015} \) es:}
{\( 3 \)}{\( 7 \)}{\( 9 \)}{\( 1 \)}{}{}}
\End

\Data\ID{1003134503}
\Updated{2020-08-28 08:08:57}
\Level{B}
\SubArea{Expressions with powers and roots}
\LANGen{
\Question{
The cube of the number \( 4+3\sqrt2 \) equals:}
{\( 280+198\sqrt2 \)}{\( 64+27\sqrt2 \)}{\( 280 + 171\sqrt2 \)}{\( 64+52\sqrt2 \)}{}{}}
\LANGpl{
\Question{
Sześcian liczby  \( 4+3\sqrt2 \) jest równy:}
{\( 280+198\sqrt2 \)}{\( 64+27\sqrt2 \)}{\( 280 + 171\sqrt2 \)}{\( 64+52\sqrt2 \)}{}{}}
\LANGcs{
\Question{
Třetí mocnina čísla \( 4+3\sqrt2 \) je rovná:}
{\( 280+198\sqrt2 \)}{\( 64+27\sqrt2 \)}{\( 280 + 171\sqrt2 \)}{\( 64+52\sqrt2 \)}{}{}}
\LANGsk{
\Question{
Tretia mocnina čísla \( 4+3\sqrt2 \) sa rovná:}
{\( 280+198\sqrt2 \)}{\( 64+27\sqrt2 \)}{\( 280 + 171\sqrt2 \)}{\( 64+52\sqrt2 \)}{}{}}
\LANGes{
\Question{
La potencia cúbica del número \( 4+3\sqrt2 \) es igual a:}
{\( 280+198\sqrt2 \)}{\( 64+27\sqrt2 \)}{\( 280 + 171\sqrt2 \)}{\( 64+52\sqrt2 \)}{}{}}
\End

\Data\ID{1003134504}
\Updated{2020-10-12 10:10:43}
\Level{B}
\SubArea{Expressions with powers and roots}
\LANGen{
\Question{
The equality \( \left( \sqrt2-a\right)^3 = 2\sqrt2+18+3\sqrt2a^2-a^3 \) is valid for:}
{\( a=-3 \)}{\( a=9 \)}{\( a=3 \)}{\( a=-9 \)}{}{}}
\LANGpl{
\Question{
Równość \( \left( \sqrt2-a\right)^3 = 2\sqrt2+18+3\sqrt2a^2-a^3 \) zachodzi dla:}
{\( a=-3 \)}{\( a=9 \)}{\( a=3 \)}{\( a=-9 \)}{}{}}
\LANGcs{
\Question{
Rovnost \( \left( \sqrt2-a\right)^3 = 2\sqrt2+18+3\sqrt2a^2-a^3 \) je pravdivá pro:}
{\( a=-3 \)}{\( a=9 \)}{\( a=3 \)}{\( a=-9 \)}{}{}}
\LANGsk{
\Question{
Rovnosť \( \left( \sqrt2-a\right)^3 = 2\sqrt2+18+3\sqrt2a^2-a^3 \) je pravdivá pre:}
{\( a=-3 \)}{\( a=9 \)}{\( a=3 \)}{\( a=-9 \)}{}{}}
\LANGes{
\Question{
La ecuacuón \( \left( \sqrt2-a\right)^3 = 2\sqrt2+18+3\sqrt2a^2-a^3 \) es valida para:}
{\( a=-3 \)}{\( a=9 \)}{\( a=3 \)}{\( a=-9 \)}{}{}}
\End

\Data\ID{1003134506}
\Updated{2022-04-23 05:04:14}
\Level{B}
\SubArea{Expressions with powers and roots}
\LANGen{
\Question{
The number \( \left( \sqrt{2-\sqrt3}-\sqrt{2+\sqrt3} \right)^2 \) equals:}
{\( 2 \)}{\( 4 \)}{\( \sqrt3 \)}{\( 2\sqrt3 \)}{}{}}
\LANGpl{
\Question{
Liczba \( \left( \sqrt{2-\sqrt3}-\sqrt{2+\sqrt3} \right)^2 \) jest równa:}
{\( 2 \)}{\( 4 \)}{\( \sqrt3 \)}{\( 2\sqrt3 \)}{}{}}
\LANGcs{
\Question{
Číslo \( \left( \sqrt{2-\sqrt3}-\sqrt{2+\sqrt3} \right)^2 \) je rovno:}
{\( 2 \)}{\( 4 \)}{\( \sqrt3 \)}{\( 2\sqrt3 \)}{}{}}
\LANGsk{
\Question{
Číslo \( \left( \sqrt{2-\sqrt3}-\sqrt{2+\sqrt3} \right)^2 \) sa rovná:}
{\( 2 \)}{\( 4 \)}{\( \sqrt3 \)}{\( 2\sqrt3 \)}{}{}}
\LANGes{
\Question{
El número \( \left( \sqrt{2-\sqrt3}-\sqrt{2+\sqrt3} \right)^2 \) es igual a:}
{\( 2 \)}{\( 4 \)}{\( \sqrt3 \)}{\( 2\sqrt3 \)}{}{}}
\End

\Data\ID{1003134507}
\Updated{2020-08-28 08:08:08}
\Level{B}
\SubArea{Expressions with powers and roots}
\LANGen{
\Question{
In the expansion of  \( \left( 2\sqrt3x+4y\right)^3 \) the coefficient of \( xy^2 \) is:}
{\( 96\sqrt3 \)}{\( 32\sqrt3 \)}{\( 48 \)}{\( 144 \)}{}{}}
\LANGpl{
\Question{
W rozwinięciu wyrażenia  \( \left( 2\sqrt3x+4y\right)^3 \) współczynnik przy iloczynie \( xy^2 \) jest równy:}
{\( 96\sqrt3 \)}{\( 32\sqrt3 \)}{\( 48 \)}{\( 144 \)}{}{}}
\LANGcs{
\Question{
V rozvoji \( \left( 2\sqrt3x+4y\right)^3 \) je koeficient členu \( xy^2 \) roven:}
{\( 96\sqrt3 \)}{\( 32\sqrt3 \)}{\( 48 \)}{\( 144 \)}{}{}}
\LANGsk{
\Question{
V rozvoji \( \left( 2\sqrt3x+4y\right)^3 \) je koeficient člena \( xy^2 \) rovný:}
{\( 96\sqrt3 \)}{\( 32\sqrt3 \)}{\( 48 \)}{\( 144 \)}{}{}}
\LANGes{
\Question{
En el desarollo de  \( \left( 2\sqrt3x+4y\right)^3 \) el coeficiente de \( xy^2 \) es:}
{\( 96\sqrt3 \)}{\( 32\sqrt3 \)}{\( 48 \)}{\( 144 \)}{}{}}
\End

\Data\ID{1003134509}
\Updated{2020-10-12 10:10:17}
\Level{B}
\SubArea{Expressions with powers and roots}
\LANGen{
\Question{
Let \( a=4^{1.5} \) and \( b=0.125^{-\frac13} \). Choose the true statement.}
{\( a=4b \)}{\( a=\frac12b \)}{\( a=2b \)}{\( a < b \)}{}{}}
\LANGpl{
\Question{
Dane są liczby \( a=4^{1{,}5} \) oraz \( b=0{,}125^{-\frac13} \). Zatem:}
{\( a=4b \)}{\( a=\frac12b \)}{\( a=2b \)}{\( a < b \)}{}{}}
\LANGcs{
\Question{
Nechť \( a=4^{1{,}5} \) a \( b=0{,}125^{-\frac13} \). Vyberte pravdivé tvrzení.}
{\( a=4b \)}{\( a=\frac12b \)}{\( a=2b \)}{\( a < b \)}{}{}}
\LANGsk{
\Question{
Nech \( a=4^{1{,}5} \) a \( b=0{,}125^{-\frac13} \). Vyberte pravdivé tvrdenie.}
{\( a=4b \)}{\( a=\frac12b \)}{\( a=2b \)}{\( a < b \)}{}{}}
\LANGes{
\Question{
Sean \( a=4^{1.5} \) y \( b=0.125^{-\frac13} \). Elija la opción correcta:}
{\( a=4b \)}{\( a=\frac12b \)}{\( a=2b \)}{\( a < b \)}{}{}}
\End

\Data\ID{1003164006}
\Updated{2020-08-28 01:08:01}
\Level{B}
\SubArea{Expressions with powers and roots}
\LANGen{
\Question{
The value of the expression
\[ \frac{x^4-16}{\left(x^2+4\right)\left(x+2\right) }\]
for \( x=2-\sqrt2 \) is equal to:}
{\( -\sqrt2 \)}{\( \sqrt2 \)}{\( 2 \)}{\( -2 \)}{}{}}
\LANGpl{
\Question{
Wartość wyrażenia
\[ \frac{x^4-16}{\left(x^2+4\right)\left(x+2\right) }\]
dla \( x=2-\sqrt2 \) jest równa:}
{\( -\sqrt2 \)}{\( \sqrt2 \)}{\( 2 \)}{\( -2 \)}{}{}}
\LANGcs{
\Question{
Hodnota výrazu
\[ \frac{x^4-16}{\left(x^2+4\right)\left(x+2\right) }\]
pro \( x=2-\sqrt2 \) je rovna:}
{\( -\sqrt2 \)}{\( \sqrt2 \)}{\( 2 \)}{\( -2 \)}{}{}}
\LANGsk{
\Question{
Hodnota výrazu
\[ \frac{x^4-16}{\left(x^2+4\right)\left(x+2\right) }\]
pre \( x=2-\sqrt2 \) je rovná:}
{\( -\sqrt2 \)}{\( \sqrt2 \)}{\( 2 \)}{\( -2 \)}{}{}}
\LANGes{
\Question{
El valor de la expresión
\[ \frac{x^4-16}{\left(x^2+4\right)\left(x+2\right) }\]
para \( x=2-\sqrt2 \) es igual a:}
{\( -\sqrt2 \)}{\( \sqrt2 \)}{\( 2 \)}{\( -2 \)}{}{}}
\End

\Data\ID{2000009403}
\Updated{2022-11-02 09:11:37}
\Level{B}
\SubArea{Expressions with powers and roots}
\LANGen{
\Question{
The expression \( \sqrt[5]{1024}\) equals:}
{\(4\)}{\(32\)}{\(\sqrt{2}\)}{\(25\)}{}{}}
\LANGpl{
\Question{
Wyrażenie \( \sqrt[5]{1024}\) jest równe:}
{\(4\)}{\(32\)}{\(\sqrt{2}\)}{\(25\)}{}{}}
\LANGcs{
\Question{
Výraz \( \sqrt[5]{1024}\) je roven:}
{\(4\)}{\(32\)}{\(\sqrt{2}\)}{\(25\)}{}{}}
\LANGsk{
\Question{
Výraz \( \sqrt[5]{1024}\) sa rovná:}
{\(4\)}{\(32\)}{\(\sqrt{2}\)}{\(25\)}{}{}}
\LANGes{
\Question{
La expresión \( \sqrt[5]{1024}\) es igual a:}
{\(4\)}{\(32\)}{\(\sqrt{2}\)}{\(25\)}{}{}}
\End

\Data\ID{2010004701}
\Updated{2022-11-02 09:11:13}
\Level{B}
\SubArea{Expressions with powers and roots}
\LANGen{
\Question{
The value of the multiplicative inverse of \( \left[ \left( \frac12 \right)^{-1} -5^{-2}\right]^{\frac12} \) is:}
{\( \frac57 \)}{\( \sqrt2-\frac15 \)}{\( \frac{25}{49} \)}{\( \frac75 \)}{}{}}
\LANGpl{
\Question{
Odwrotnością wyrażenia \( \left[ \left( \frac12 \right)^{-1} -5^{-2}\right]^{\frac12} \) jest:}
{\( \frac57 \)}{\( \sqrt2-\frac15 \)}{\( \frac{25}{49} \)}{\( \frac75 \)}{}{}}
\LANGcs{
\Question{
Převrácená hodnota výrazu \( \left[ \left( \frac12 \right)^{-1} -5^{-2}\right]^{\frac12} \) je:}
{\( \frac57 \)}{\( \sqrt2-\frac15 \)}{\( \frac{25}{49} \)}{\( \frac75 \)}{}{}}
\LANGsk{
\Question{
Určte prevrátenú hodnotu výrazu \( \left[ \left( \frac12 \right)^{-1} -5^{-2}\right]^{\frac12} \).}
{\( \frac57 \)}{\( \sqrt2-\frac15 \)}{\( \frac{25}{49} \)}{\( \frac75 \)}{}{}}
\LANGes{
\Question{
El valor del inverso multiplicativo de \( \left[ \left( \frac12 \right)^{-1} -5^{-2}\right]^{\frac12} \) es:}
{\( \frac57 \)}{\( \sqrt2-\frac15 \)}{\( \frac{25}{49} \)}{\( \frac75 \)}{}{}}
\End

\Data\ID{2010004702}
\Updated{2022-05-16 11:05:43}
\Level{B}
\SubArea{Expressions with powers and roots}
\LANGen{
\Question{
Simplify the fraction \( \frac{\sqrt[12]8\cdot\sqrt[20]{16}\cdot\sqrt[5]{15}}{\sqrt[5]{30}} \).}
{\( \sqrt[4]2 \)}{\( \frac1{\sqrt[4]2} \)}{\( 1 \)}{\( 2 \)}{}{}}
\LANGpl{
\Question{
Uprość ułamek \( \frac{\sqrt[12]8\cdot\sqrt[20]{16}\cdot\sqrt[5]{15}}{\sqrt[5]{30}} \).}
{\( \sqrt[4]2 \)}{\( \frac1{\sqrt[4]2} \)}{\( 1 \)}{\( 2 \)}{}{}}
\LANGcs{
\Question{
Zjednodušte zlomek \( \frac{\sqrt[12]8\cdot\sqrt[20]{16}\cdot\sqrt[5]{15}}{\sqrt[5]{30}} \).}
{\( \sqrt[4]2 \)}{\( \frac1{\sqrt[4]2} \)}{\( 1 \)}{\( 2 \)}{}{}}
\LANGsk{
\Question{
Zjednodušte zlomok \( \frac{\sqrt[12]8\cdot\sqrt[20]{16}\cdot\sqrt[5]{15}}{\sqrt[5]{30}} \).}
{\( \sqrt[4]2 \)}{\( \frac1{\sqrt[4]2} \)}{\( 1 \)}{\( 2 \)}{}{}}
\LANGes{
\Question{
Simplifica la fracción \( \frac{\sqrt[12]8\cdot\sqrt[20]{16}\cdot\sqrt[5]{15}}{\sqrt[5]{30}} \).}
{\( \sqrt[4]2 \)}{\( \frac1{\sqrt[4]2} \)}{\( 1 \)}{\( 2 \)}{}{}}
\End

\Data\ID{2010004703}
\Updated{2022-11-02 09:11:13}
\Level{B}
\SubArea{Expressions with powers and roots}
\LANGen{
\Question{
Simplifying \( \left( \sqrt[5]{2\sqrt[3]8} \right)^{\frac52}\cdot  \sqrt{4^{-1}} \) we get:}
{\( 1 \)}{\( \sqrt2 \)}{\( \frac{1}{\sqrt{2}} \)}{\( 2 \)}{}{}}
\LANGpl{
\Question{
Upraszczając \( \left( \sqrt[5]{2\sqrt[3]8} \right)^{\frac52}\cdot  \sqrt{4^{-1}} \) otrzymamy:}
{\( 1 \)}{\( \sqrt2 \)}{\( \frac{1}{\sqrt{2}} \)}{\( 2 \)}{}{}}
\LANGcs{
\Question{
Zjednodušením \( \left( \sqrt[5]{2\sqrt[3]8} \right)^{\frac52}\cdot \sqrt{4^{-1}} \) dostaneme:}
{\( 1 \)}{\( \sqrt2 \)}{\( \frac{1}{\sqrt{2}} \)}{\( 2 \)}{}{}}
\LANGsk{
\Question{
Zjednodušením výrazu \( \left( \sqrt[5]{2\sqrt[3]8} \right)^{\frac52}\cdot  \sqrt{4^{-1}} \) dostaneme:}
{\( 1 \)}{\( \sqrt2 \)}{\( \frac{1}{\sqrt{2}} \)}{\( 2 \)}{}{}}
\LANGes{
\Question{
Simplificando \( \left( \sqrt[5]{2\sqrt[3]8} \right)^{\frac52}\cdot  \sqrt{4^{-1}} \) obtenemos:}
{\( 1 \)}{\( \sqrt2 \)}{\( \frac{1}{\sqrt{2}} \)}{\( 2 \)}{}{}}
\End

\Data\ID{2010004704}
\Updated{2022-11-02 09:11:13}
\Level{B}
\SubArea{Expressions with powers and roots}
\LANGen{
\Question{
The number \( \left(\frac{81^{-3}\cdot{16}^{-3}}{9^{-5}\cdot4^{-4}}\right)^{-2} \) equals:}
{\( 12^{4} \)}{\( 6^4 \)}{\( 6^{12} \)}{\( \frac1{3^4\cdot2^{8}} \)}{}{}}
\LANGpl{
\Question{
Liczba \( \left(\frac{81^{-3}\cdot{16}^{-3}}{9^{-5}\cdot4^{-4}}\right)^{-2} \) jest równa:}
{\( 12^{4} \)}{\( 6^4 \)}{\( 6^{12} \)}{\( \frac1{3^4\cdot2^{8}} \)}{}{}}
\LANGcs{
\Question{
Číslo \( \left(\frac{81^{-3}\cdot{16}^{-3}}{9^{-5}\cdot4^{-4}}\right)^{-2} \) se rovná:}
{\( 12^{4} \)}{\( 6^4 \)}{\( 6^{12} \)}{\( \frac1{3^4\cdot2^{8}} \)}{}{}}
\LANGsk{
\Question{
Číselný výraz \( \left(\frac{81^{-3}\cdot{16}^{-3}}{9^{-5}\cdot4^{-4}}\right)^{-2} \) sa rovná:}
{\( 12^{4} \)}{\( 6^4 \)}{\( 6^{12} \)}{\( \frac1{3^4\cdot2^{8}} \)}{}{}}
\LANGes{
\Question{
El número \( \left(\frac{81^{-3}\cdot{16}^{-3}}{9^{-5}\cdot4^{-4}}\right)^{-2} \) es igual a:}
{\( 12^{4} \)}{\( 6^4 \)}{\( 6^{12} \)}{\( \frac1{3^4\cdot2^{8}} \)}{}{}}
\End

\Data\ID{2010004705}
\Updated{2022-11-02 09:11:13}
\Level{B}
\SubArea{Expressions with powers and roots}
\LANGen{
\Question{
Express the value of the expression \( \left(\frac32-3^{-2}\right)^{-1} \) as a decimal number.}
{\( 0.72 \)}{\( 1.3\overline8\)}{\( \frac{18}{25}\)}{\( -0.1\overline3 \)}{}{}}
\LANGpl{
\Question{
Wyraź wartość wyrażenia \( \left(\frac32-3^{-2}\right)^{-1} \) jako liczbę dziesiętną.}
{\( 0{,}72 \)}{\( 1{,}3\overline8\)}{\( \frac{18}{25}\)}{\( -0{,}1\overline3 \)}{}{}}
\LANGcs{
\Question{
Vyjádřete hodnotu výrazu\( \left(\frac32-3^{-2}\right)^{-1} \) jako desetinné číslo.}
{\( 0{,}72 \)}{\( 1{,}3\overline8\)}{\( \frac{18}{25}\)}{\( -0{,}1\overline3 \)}{}{}}
\LANGsk{
\Question{
Vyjadrite hodnotu výrazu \( \left(\frac32-3^{-2}\right)^{-1} \) ako desatinné číslo.}
{\( 0{,}72 \)}{\( 1{,}3\overline8\)}{\( \frac{18}{25}\)}{\( -0{,}1\overline3 \)}{}{}}
\LANGes{
\Question{
Encuentra el valor de la expresión \( \left(\frac32-3^{-2}\right)^{-1} \) como un número decimal.}
{\( 0.72 \)}{\( 1.3\overline8\)}{\( \frac{18}{25}\)}{\( -0.1\overline3 \)}{}{}}
\End

\Data\ID{2010004706}
\Updated{2022-11-02 09:11:13}
\Level{B}
\SubArea{Expressions with powers and roots}
\LANGen{
\Question{
The multiplicative inverse of the number \( \frac{\sqrt{4^3}:8^{\frac13}}{\sqrt[3]4} \)  is:}
{\( 2^{-\frac43} \)}{\( 2^{\frac34} \)}{\( 2^{\frac43} \)}{\( 2^{-\frac13} \)}{}{}}
\LANGpl{
\Question{
Liczbą odwrotną wyrażenia \( \frac{\sqrt{4^3}:8^{\frac13}}{\sqrt[3]4} \) jest:}
{\( 2^{-\frac43} \)}{\( 2^{\frac34} \)}{\( 2^{\frac43} \)}{\( 2^{-\frac13} \)}{}{}}
\LANGcs{
\Question{
Převrácená hodnota čísla \( \frac{\sqrt{4^3}:8^{\frac13}}{\sqrt[3]4} \)  je:}
{\( 2^{-\frac43} \)}{\( 2^{\frac34} \)}{\( 2^{\frac43} \)}{\( 2^{-\frac13} \)}{}{}}
\LANGsk{
\Question{
Prevrátená hodnota ku hodnote číselného výrazu \( \frac{\sqrt{4^3}:8^{\frac13}}{\sqrt[3]4} \)  je:}
{\( 2^{-\frac43} \)}{\( 2^{\frac34} \)}{\( 2^{\frac43} \)}{\( 2^{-\frac13} \)}{}{}}
\LANGes{
\Question{
El inverso multiplicativo del número \( \frac{\sqrt{4^3}:8^{\frac13}}{\sqrt[3]4} \) es:}
{\( 2^{-\frac43} \)}{\( 2^{\frac34} \)}{\( 2^{\frac43} \)}{\( 2^{-\frac13} \)}{}{}}
\End

\Data\ID{2010004707}
\Updated{2022-11-02 09:11:13}
\Level{B}
\SubArea{Expressions with powers and roots}
\LANGen{
\Question{
The number \( \frac1{4^{2020}}\cdot(0.002)^{2020} \) equals:}
{\( (0.0005)^{2020} \)}{\( \frac1{5000^{2020}} \)}{\( (0.008)^{2020} \)}{\( (0.005)^{2020} \)}{}{}}
\LANGpl{
\Question{
Liczba \( \frac1{4^{2020}}\cdot(0{,}002)^{2020} \) jest równa:}
{\( (0{,}0005)^{2020} \)}{\( \frac1{5000^{2020}} \)}{\( (0{,}008)^{2020} \)}{\( (0{,}005)^{2020} \)}{}{}}
\LANGcs{
\Question{
Číslo \( \frac1{4^{2020}}\cdot(0{,}002)^{2020} \) je rovno:}
{\( (0{,}0005)^{2020} \)}{\( \frac1{5000^{2020}} \)}{\( (0{,}008)^{2020} \)}{\( (0{,}005)^{2020} \)}{}{}}
\LANGsk{
\Question{
Číselný výraz \( \frac1{4^{2020}}\cdot(0{,}002)^{2020} \) sa rovná:}
{\( (0{,}0005)^{2020} \)}{\( \frac1{5000^{2020}} \)}{\( (0{,}008)^{2020} \)}{\( (0{,}005)^{2020} \)}{}{}}
\LANGes{
\Question{
El número \( \frac1{4^{2020}}\cdot(0.002)^{2020} \) es igual a:}
{\( (0.0005)^{2020} \)}{\( \frac1{5000^{2020}} \)}{\( (0.008)^{2020} \)}{\( (0.005)^{2020} \)}{}{}}
\End

\Data\ID{2010004802}
\Updated{2022-08-25 10:08:52}
\Level{B}
\SubArea{Expressions with powers and roots}
\LANGen{
\Question{
The number \( \left( \sqrt{3+\sqrt5}+\sqrt{3-\sqrt5} \right)^2 \) equals:}
{\( 10 \)}{\( 6 \)}{\( 14 \)}{\( 2\sqrt5 \)}{}{}}
\LANGpl{
\Question{
Liczba \( \left( \sqrt{3+\sqrt5}+\sqrt{3-\sqrt5} \right)^2 \) jest równa:}
{\( 10 \)}{\( 6 \)}{\( 14 \)}{\( 2\sqrt5 \)}{}{}}
\LANGcs{
\Question{
Číslo \( \left( \sqrt{3+\sqrt5}+\sqrt{3-\sqrt5} \right)^2 \) je rovno:}
{\( 10 \)}{\( 6 \)}{\( 14 \)}{\( 2\sqrt5 \)}{}{}}
\LANGsk{
\Question{
Číslo \( \left( \sqrt{3+\sqrt5}+\sqrt{3-\sqrt5} \right)^2 \) je rovné:}
{\( 10 \)}{\( 6 \)}{\( 14 \)}{\( 2\sqrt5 \)}{}{}}
\LANGes{
\Question{
El número \( \left( \sqrt{3+\sqrt5}+\sqrt{3-\sqrt5} \right)^2 \) es igual a:}
{\( 10 \)}{\( 6 \)}{\( 14 \)}{\( 2\sqrt5 \)}{}{}}
\End

\Data\ID{2010004803}
\Updated{2022-11-02 09:11:13}
\Level{B}
\SubArea{Expressions with powers and roots}
\LANGen{
\Question{
For \(x\in \mathbb{R}\),
\(x > 0\),
simplify the following expression. 
\[ 
 \root{5}\of{x^{4}} : \root{3}\of{x^2}
\]}
{\(\root{15}\of{x^{2}}\)}{\(\root{15}\of{x^{22}}\)}{\(\root{5}\of{x^{9}}\)}{\(x\)}{}{}}
\LANGpl{
\Question{
Uprość wyrażenie dla \(x\in \mathbb{R}\),
\(x > 0\). 
\[ 
 \root{5}\of{x^{4}} : \root{3}\of{x^2}
\]}
{\(\root{15}\of{x^{2}}\)}{\(\root{15}\of{x^{22}}\)}{\(\root{5}\of{x^{9}}\)}{\(x\)}{}{}}
\LANGcs{
\Question{
Je-li \(x\)
kladné reálné číslo, pak je výraz
\[ 
 \root{5}\of{x^{4}} : \root{3}\of{x^2}
\] roven:}
{\(\root{15}\of{x^{2}}\)}{\(\root{15}\of{x^{22}}\)}{\(\root{5}\of{x^{9}}\)}{\(x\)}{}{}}
\LANGsk{
\Question{
Ak je \(x\) kladné reálne číslo, potom sa výraz
 \(
 \root{5}\of{x^{4}} : \root{3}\of{x^2}
\)
rovná výrazu:}
{\(\root{15}\of{x^{2}}\)}{\(\root{15}\of{x^{22}}\)}{\(\root{5}\of{x^{9}}\)}{\(x\)}{}{}}
\LANGes{
\Question{
Para \(x\in \mathbb{R}\),
\(x > 0\),
simplifica la siguiente expresión: 
\[ 
 \root{5}\of{x^{4}} : \root{3}\of{x^2}
\]}
{\(\root{15}\of{x^{2}}\)}{\(\root{15}\of{x^{22}}\)}{\(\root{5}\of{x^{9}}\)}{\(x\)}{}{}}
\End

\Data\ID{2010004804}
\Updated{2022-11-02 09:11:34}
\Level{B}
\SubArea{Expressions with powers and roots}
\LANGen{
\Question{
For \(x\in \mathbb{R}\),
\(x > 0\),
simplify the following expression. 
\[ 
 \root{3}\of{x}\cdot \root{5}\of{x^{3}}
\]}
{\(\root{15}\of{x^{14}}\)}{\(\root{5}\of{x}\)}{\(\root{15}\of{x^{4}}\)}{\(\sqrt{x}\)}{}{}}
\LANGpl{
\Question{
Uprość wyrażenie dla  \(x\in \mathbb{R}\),
\(x > 0\). 
\[ 
 \root{3}\of{x}\cdot \root{5}\of{x^{3}}
\]}
{\(\root{15}\of{x^{14}}\)}{\(\root{5}\of{x}\)}{\(\root{15}\of{x^{4}}\)}{\(\sqrt{x}\)}{}{}}
\LANGcs{
\Question{
Je-li \(x\)
kladné reálné číslo, pak je výraz 
\[ 
 \root{3}\of{x}\cdot \root{5}\of{x^{3}}
\] roven:}
{\(\root{15}\of{x^{14}}\)}{\(\root{5}\of{x}\)}{\(\root{15}\of{x^{4}}\)}{\(\sqrt{x}\)}{}{}}
\LANGsk{
\Question{
Ak je \(x\) kladné reálne číslo, potom sa výraz
 \(
 \root{3}\of{x}\cdot \root{5}\of{x^{3}}
\)
rovná výrazu:}
{\(\root{15}\of{x^{14}}\)}{\(\root{5}\of{x}\)}{\(\root{15}\of{x^{4}}\)}{\(\sqrt{x}\)}{}{}}
\LANGes{
\Question{
Para \(x\in \mathbb{R}\),
\(x > 0\),
simplifica la siguiente expresión. 
\[ 
 \root{3}\of{x}\cdot \root{5}\of{x^{3}}
\]}
{\(\root{15}\of{x^{14}}\)}{\(\root{5}\of{x}\)}{\(\root{15}\of{x^{4}}\)}{\(\sqrt{x}\)}{}{}}
\End

\Data\ID{2010004806}
\Updated{2022-11-02 09:11:34}
\Level{B}
\SubArea{Expressions with powers and roots}
\LANGen{
\Question{
Simplify the expression \(0.3^{\frac{4} 
{5} }\cdot 0.3^{-\frac{3} 
{10} }\) 
and write the result using a root.}
{\(\sqrt{0.3}\)}{\(\root{5}\of{0.3}\)}{\(\root{10}\of{0.3^{7}}\)}{\(\root{3}\of{0.3^8}\)}{}{}}
\LANGpl{
\Question{
Uprość wyrażenie \(0{,}3^{\frac{4} 
{5} }\cdot 0{,}3^{-\frac{3} 
{10} }\) 
i wynik zapisz w postaci pierwiastka.}
{\(\sqrt{0{,}3}\)}{\(\root{5}\of{0{,}3}\)}{\(\root{10}\of{0{,}3^{7}}\)}{\(\root{3}\of{0{,}3^8}\)}{}{}}
\LANGcs{
\Question{
Číslo \(0{,}3^{\frac{4} 
{5} }\cdot 0{,}3^{-\frac{3} 
{10} }\) 
zjednodušte a výsledek zapište ve tvaru odmocniny.}
{\(\sqrt{0{,}3}\)}{\(\root{5}\of{0{,}3}\)}{\(\root{10}\of{0{,}3^{7}}\)}{\(\root{3}\of{0{,}3^8}\)}{}{}}
\LANGsk{
\Question{
Číslo \(0{,}3^{\frac{4} 
{5} }\cdot 0{,}3^{-\frac{3} 
{10} }\) zjednodušte a zapíšte v tvare odmocniny.}
{\(\sqrt{0{,}3}\)}{\(\root{5}\of{0{,}3}\)}{\(\root{10}\of{0{,}3^{7}}\)}{\(\root{3}\of{0{,}3^8}\)}{}{}}
\LANGes{
\Question{
Simplifica la expresión \(0.3^{\frac{4} 
{5} }\cdot 0.3^{-\frac{3} 
{10} }\) 
y escribe el resultado usando una raíz.}
{\(\sqrt{0.3}\)}{\(\root{5}\of{0.3}\)}{\(\root{10}\of{0.3^{7}}\)}{\(\root{3}\of{0.3^8}\)}{}{}}
\End

\Data\ID{2010004807}
\Updated{2022-11-02 09:11:34}
\Level{B}
\SubArea{Expressions with powers and roots}
\LANGen{
\Question{
Write the number \(\root{12}\of{5^{-4}}\) 
as a power with a rational exponent.}
{\(5^{-\frac{1} 
{3} }\)}{\(5^{\frac{1} 
{3} }\)}{\(5^{3}\)}{\(5^{-3}\)}{}{}}
\LANGpl{
\Question{
Zapisz liczbę \(\root{12}\of{5^{-4}}\) 
jako potęgę z wykładnikiem wymiernym.}
{\(5^{-\frac{1} 
{3} }\)}{\(5^{\frac{1} 
{3} }\)}{\(5^{3}\)}{\(5^{-3}\)}{}{}}
\LANGcs{
\Question{
Číslo \(\root{12}\of{5^{-4}}\) 
zapište ve tvaru mocniny.}
{\(5^{-\frac{1} 
{3} }\)}{\(5^{\frac{1} 
{3} }\)}{\(5^{3}\)}{\(5^{-3}\)}{}{}}
\LANGsk{
\Question{
Číslo \(\root{12}\of{5^{-4}}\) 
zapíšte v tvare mocniny.}
{\(5^{-\frac{1} 
{3} }\)}{\(5^{\frac{1} 
{3} }\)}{\(5^{3}\)}{\(5^{-3}\)}{}{}}
\LANGes{
\Question{
Escribe el número \(\root{12}\of{5^{-4}}\) 
como una potencia con exponente racional.}
{\(5^{-\frac{1} 
{3} }\)}{\(5^{\frac{1} 
{3} }\)}{\(5^{3}\)}{\(5^{-3}\)}{}{}}
\End

\Data\ID{2010004808}
\Updated{2022-11-02 09:11:34}
\Level{B}
\SubArea{Expressions with powers and roots}
\LANGen{
\Question{
Assuming \(x\not \in \{-4;-1;0;4\}\),
simplify the following expression. 
\[ 
 \left(\frac{x+1} 
{x^{2} +4x}\right)^{-1}\cdot \frac{x^{2} +x}
 {x^2 - 16} 
\]}
{\(\frac{x^2}
 {x-4} \)}{\(\frac{x^2}
 {x+4} \)}{\(\frac{2x}
 {x-4} \)}{\(\frac{x}
 {x-4} \)}{}{}}
\LANGpl{
\Question{
Uprość wyrażenie dla \(x\not \in \{-4;-1;0;4\}\). 
\[ 
 \left(\frac{x+1} 
{x^{2} +4x}\right)^{-1}\cdot \frac{x^{2} +x}
 {x^2 - 16} 
\]}
{\(\frac{x^2}
 {x-4} \)}{\(\frac{x^2}
 {x+4} \)}{\(\frac{2x}
 {x-4} \)}{\(\frac{x}
 {x-4} \)}{}{}}
\LANGcs{
\Question{
Za předpokladu, že \(x\not \in \{-4;-1;0;4\}\),
upravte na co nejjednodušší tvar výraz
\[ 
 \left(\frac{x+1} 
{x^{2} +4x}\right)^{-1}\cdot \frac{x^{2} +x}
 {x^2 - 16}. 
\]}
{\(\frac{x^2}
 {x-4} \)}{\(\frac{x^2}
 {x+4} \)}{\(\frac{2x}
 {x-4} \)}{\(\frac{x}
 {x-4} \)}{}{}}
\LANGsk{
\Question{
Za predpokladu, že \(x\not \in \{-4;-1;0;4\}\),
zjednodušte daný výraz
\( 
 \left(\frac{x+1} 
{x^{2} +4x}\right)^{-1}\cdot \frac{x^{2} +x}
 {x^2 - 16} 
\).}
{\(\frac{x^2}
 {x-4} \)}{\(\frac{x^2}
 {x+4} \)}{\(\frac{2x}
 {x-4} \)}{\(\frac{x}
 {x-4} \)}{}{}}
\LANGes{
\Question{
Suponiendo que \(x\not \in \{-4;-1;0;4\}\),
simplifica la siguiente expresión. 
\[ 
 \left(\frac{x+1} 
{x^{2} +4x}\right)^{-1}\cdot \frac{x^{2} +x}
 {x^2 - 16} 
\]}
{\(\frac{x^2}
 {x-4} \)}{\(\frac{x^2}
 {x+4} \)}{\(\frac{2x}
 {x-4} \)}{\(\frac{x}
 {x-4} \)}{}{}}
\End

\Data\ID{2010005701}
\Updated{2022-11-02 09:11:36}
\Level{B}
\SubArea{Expressions with powers and roots}
\LANGen{
\Question{
The product of the numbers \( \left(\sqrt[6]{3}\cdot 81^{\frac14}\right)^{-1} \)   and  \( \sqrt[3]9\cdot \sqrt{3}\) equals:}
{\( 1 \)}{\( \frac13 \)}{\( \sqrt3 \)}{\( 3 \)}{}{}}
\LANGpl{
\Question{
Iloczyn liczb\( \left(\sqrt[6]{3}\cdot 81^{\frac14}\right)^{-1} \) i \( \sqrt[3]9\cdot \sqrt{3}\) wynosi:}
{\( 1 \)}{\( \frac13 \)}{\( \sqrt3 \)}{\( 3 \)}{}{}}
\LANGcs{
\Question{
Součin čísel \( \left(\sqrt[6]{3}\cdot 81^{\frac14}\right)^{-1} \)   a  \( \sqrt[3]9\cdot \sqrt{3}\) je roven:}
{\( 1 \)}{\( \frac13 \)}{\( \sqrt3 \)}{\( 3 \)}{}{}}
\LANGsk{
\Question{
Súčin číselných výrazov \( \left(\sqrt[6]{3}\cdot 81^{\frac14}\right)^{-1} \)   a  \( \sqrt[3]9\cdot \sqrt{3}\) je rovný:}
{\( 1 \)}{\( \frac13 \)}{\( \sqrt3 \)}{\( 3 \)}{}{}}
\LANGes{
\Question{
El producto de los números \( \left(\sqrt[6]{3}\cdot 81^{\frac14}\right)^{-1} \)   y  \( \sqrt[3]9\cdot \sqrt{3}\) es igual a:}
{\( 1 \)}{\( \frac13 \)}{\( \sqrt3 \)}{\( 3 \)}{}{}}
\End

\Data\ID{2010005702}
\Updated{2022-11-02 09:11:36}
\Level{B}
\SubArea{Expressions with powers and roots}
\LANGen{
\Question{
The number \(\sqrt[5]{(-3)^2}\cdot\left(\frac1{81}\right)^{-\frac25} \) equals:}
{\( 9 \)}{\( -9 \)}{\( 3 \)}{\( -3 \)}{}{}}
\LANGpl{
\Question{
Liczba \(\sqrt[5]{(-3)^2}\cdot\left(\frac1{81}\right)^{-\frac25} \) jest równa:}
{\( 9 \)}{\( -9 \)}{\( 3 \)}{\( -3 \)}{}{}}
\LANGcs{
\Question{
Číslo \(\sqrt[5]{(-3)^2}\cdot\left(\frac1{81}\right)^{-\frac25} \) je rovno:}
{\( 9 \)}{\( -9 \)}{\( 3 \)}{\( -3 \)}{}{}}
\LANGsk{
\Question{
Číselný výraz \(\sqrt[5]{(-3)^2}\cdot\left(\frac1{81}\right)^{-\frac25} \) sa rovná:}
{\( 9 \)}{\( -9 \)}{\( 3 \)}{\( -3 \)}{}{}}
\LANGes{
\Question{
El número \(\sqrt[5]{(-3)^2}\cdot\left(\frac1{81}\right)^{-\frac25} \) es igual a:}
{\( 9 \)}{\( -9 \)}{\( 3 \)}{\( -3 \)}{}{}}
\End

\Data\ID{2010005703}
\Updated{2022-11-02 09:11:36}
\Level{B}
\SubArea{Expressions with powers and roots}
\LANGen{
\Question{
Calculate \( \left(4{.}3\cdot10^{-6}\right)\cdot\left(3\cdot10^{9}\right) \)  and give the result in the scientific   notation.}
{\( 12\,900 = 1{.}29\cdot10^4\)}{\( 12\,900 = 12{.}9\cdot10^3\)}{\( 12\,900 = 129\cdot10^2\)}{\( 0.000\,000\,000\,000 \,001\,29= 1{.}29\cdot10^{-15}\)}{}{}}
\LANGpl{
\Question{
Oblicz \( \left(4{,}3\cdot10^{-6}\right)\cdot\left(3\cdot10^{9}\right) \)  i podaj wynik w notacji wykładniczej.}
{\( 12\,900 = 1{,}29\cdot10^4\)}{\( 12\,900 = 12{,}9\cdot10^3\)}{\( 12\,900 = 129\cdot10^2\)}{\( 0{,}000\,000\,000\,000 \,001\,29= 1{,}29\cdot10^{-15}\)}{}{}}
\LANGcs{
\Question{
Vypočítejte hodnotu výrazu \( \left(4{,}3\cdot10^{-6}\right)\cdot\left(3\cdot10^{9}\right) \)  a výsledek zapište vědeckým zápisem čísla.}
{\( 12\,900 = 1{,}29\cdot10^4\)}{\( 12\,900 = 12{,}9\cdot10^3\)}{\( 12\,900 = 129\cdot10^2\)}{\( 0{,}000\,000\,000\,000 \,001\,29= 1{,}29\cdot10^{-15}\)}{}{}}
\LANGsk{
\Question{
Vypočítajte hodnotu výrazu \( \left(4{,}3\cdot10^{-6}\right)\cdot\left(3\cdot10^{9}\right) \) a výsledok zapíšte v tvare vedeckého zápisu čísel.}
{\( 12\,900 = 1{,}29\cdot10^4\)}{\( 12\,900 = 12{,}9\cdot10^3\)}{\( 12\,900 = 129\cdot10^2\)}{\( 0{,}000\,000\,000\,000 \,001\,29= 1{,}29\cdot10^{-15}\)}{}{}}
\LANGes{
\Question{
Calcula \( \left(4{.}3\cdot10^{-6}\right)\cdot\left(3\cdot10^{9}\right) \)  y aporta el resultado en notación científica.}
{\( 12\,900 = 1{.}29\cdot10^4\)}{\( 12\,900 = 12{.}9\cdot10^3\)}{\( 12\,900 = 129\cdot10^2\)}{\( 0.000\,000\,000\,000 \,001\,29= 1{.}29\cdot10^{-15}\)}{}{}}
\End

\Data\ID{2010005704}
\Updated{2022-11-02 09:11:40}
\Level{B}
\SubArea{Expressions with powers and roots}
\LANGen{
\Question{
Writing the expression \( \frac{(0.25)^{-2}\cdot (x \colon y^2)^{-2} }{(2y)^4\cdot x^{-2}}\), \( x\neq0\), \( y\neq0 \) in a simplified form, we will get:}
{\( 1\)}{\( y^{-8} \)}{\( x^{-4} \)}{\( \frac14 \)}{}{}}
\LANGpl{
\Question{
Po uproszczeniu wyrażenia\( \frac{(0{,}25)^{-2}\cdot (x \colon y^2)^{-2} }{(2y)^4\cdot x^{-2}}\), gdzie \( x\neq0\), \( y\neq0 \) otrzymamy:}
{\( 1\)}{\( y^{-8} \)}{\( x^{-4} \)}{\( \frac14 \)}{}{}}
\LANGcs{
\Question{
Když výraz \( \frac{(0{,}25)^{-2}\cdot (x \colon y^2)^{-2} }{(2y)^4\cdot x^{-2}}\), \( x\neq0\), \( y\neq0 \) zjednodušíme, dostaneme:}
{\( 1\)}{\( y^{-8} \)}{\( x^{-4} \)}{\( \frac14 \)}{}{}}
\LANGsk{
\Question{
Zjednodušte daný výraz \( \frac{(0{,}25)^{-2}\cdot (x \colon y^2)^{-2} }{(2y)^4\cdot x^{-2}}\), ak \( x\neq0\), \( y\neq0 \).}
{\( 1\)}{\( y^{-8} \)}{\( x^{-4} \)}{\( \frac14 \)}{}{}}
\LANGes{
\Question{
Escribiendo la expresión \( \frac{(0.25)^{-2}\cdot (x \colon y^2)^{-2} }{(2y)^4\cdot x^{-2}}\), \( x\neq0\), \( y\neq0 \) de forma simplificada, obtenemos:}
{\( 1\)}{\( y^{-8} \)}{\( x^{-4} \)}{\( \frac14 \)}{}{}}
\End

\Data\ID{2010005705}
\Updated{2022-11-02 09:11:40}
\Level{B}
\SubArea{Expressions with powers and roots}
\LANGen{
\Question{
Writing the expression \( \frac{16\cdot \sqrt[4]{4}\cdot \sqrt[3]{\frac12}}{\sqrt{8}\cdot\sqrt[3]{2}\cdot 4\cdot \sqrt[6]{16}} \) in the form of a  power of \( 2 \), we will get:}
{\( 2^{-\frac13} \)}{\( 2^{\frac13} \)}{\( 2^{-\frac23} \)}{\( 2^{\frac56} \)}{}{}}
\LANGpl{
\Question{
Zapisując wyrażenie \( \frac{16\cdot \sqrt[4]{4}\cdot \sqrt[3]{\frac12}}{\sqrt{8}\cdot\sqrt[3]{2}\cdot 4\cdot \sqrt[6]{16}} \) w postaci potęgi \( 2 \), otrzymamy:}
{\( 2^{-\frac13} \)}{\( 2^{\frac13} \)}{\( 2^{-\frac23} \)}{\( 2^{\frac56} \)}{}{}}
\LANGcs{
\Question{
Zjednodušte výraz \( \frac{16\cdot \sqrt[4]{4}\cdot \sqrt[3]{\frac12}}{\sqrt{8}\cdot\sqrt[3]{2}\cdot 4\cdot \sqrt[6]{16}} \)  a výsledek zapište ve tvaru mocniny čísla \( 2 \).}
{\( 2^{-\frac13} \)}{\( 2^{\frac13} \)}{\( 2^{-\frac23} \)}{\( 2^{\frac56} \)}{}{}}
\LANGsk{
\Question{
Zjednodušte výraz \( \frac{16\cdot \sqrt[4]{4}\cdot \sqrt[3]{\frac12}}{\sqrt{8}\cdot\sqrt[3]{2}\cdot 4\cdot \sqrt[6]{16}} \) a výsledok zapíšte v tvare mocniny čísla \( 2 \).}
{\( 2^{-\frac13} \)}{\( 2^{\frac13} \)}{\( 2^{-\frac23} \)}{\( 2^{\frac56} \)}{}{}}
\LANGes{
\Question{
Escribiendo la expresión \( \frac{16\cdot \sqrt[4]{4}\cdot \sqrt[3]{\frac12}}{\sqrt{8}\cdot\sqrt[3]{2}\cdot 4\cdot \sqrt[6]{16}} \) como potencia de \( 2 \), obtenemos:}
{\( 2^{-\frac13} \)}{\( 2^{\frac13} \)}{\( 2^{-\frac23} \)}{\( 2^{\frac56} \)}{}{}}
\End

\Data\ID{2010005706}
\Updated{2022-11-02 09:11:40}
\Level{B}
\SubArea{Expressions with powers and roots}
\LANGen{
\Question{
What is the value of \( \left(3^{\sqrt6+\sqrt2}\right)^{\sqrt6-\sqrt2} \)?}
{\( 81\)}{\( 3^{8} \)}{\( 3^{\sqrt{32}} \)}{\( 3^{2\sqrt{6}} \)}{}{}}
\LANGpl{
\Question{
Oblicz wartość wyrażenia \( \left(3^{\sqrt6+\sqrt2}\right)^{\sqrt6-\sqrt2} \).}
{\( 81\)}{\( 3^{8} \)}{\( 3^{\sqrt{32}} \)}{\( 3^{2\sqrt{6}} \)}{}{}}
\LANGcs{
\Question{
Jaká je hodnota výrazu \( \left(3^{\sqrt6+\sqrt2}\right)^{\sqrt6-\sqrt2} \)?}
{\( 81\)}{\( 3^{8} \)}{\( 3^{\sqrt{32}} \)}{\( 3^{2\sqrt{6}} \)}{}{}}
\LANGsk{
\Question{
Aká je hodnota výrazu \( \left(3^{\sqrt6+\sqrt2}\right)^{\sqrt6-\sqrt2} \)?}
{\( 81\)}{\( 3^{8} \)}{\( 3^{\sqrt{32}} \)}{\( 3^{2\sqrt{6}} \)}{}{}}
\LANGes{
\Question{
¿Cuál es el valor de \( \left(3^{\sqrt6+\sqrt2}\right)^{\sqrt6-\sqrt2} \)?}
{\( 81\)}{\( 3^{8} \)}{\( 3^{\sqrt{32}} \)}{\( 3^{2\sqrt{6}} \)}{}{}}
\End

\Data\ID{2010005707}
\Updated{2022-11-02 09:11:40}
\Level{B}
\SubArea{Expressions with powers and roots}
\LANGen{
\Question{
Which of the given numbers belongs to the interval  \( [ -5;5 ]\)?}
{\( \left(\sqrt2\right)^4-\left(\sqrt3\right)^4 \)}{\( 2\left(\sqrt{0{.}1}\right)^2\cdot\left(\sqrt2\right)^{10} \)}{\( \left(2\sqrt{5}\right)^2-\left(\sqrt3\right)^6 \)}{\( 2\left(\sqrt{0{.}1}\right)^4+\left(\sqrt2\right)^8 \)}{}{}}
\LANGpl{
\Question{
Która z podanych liczb należy do przedziału  \( \langle -5;5 \rangle\)?}
{\( \left(\sqrt2\right)^4-\left(\sqrt3\right)^4 \)}{\( 2\left(\sqrt{0{,}1}\right)^2\cdot\left(\sqrt2\right)^{10} \)}{\( \left(2\sqrt{5}\right)^2-\left(\sqrt3\right)^6 \)}{\( 2\left(\sqrt{0{,}1}\right)^4+\left(\sqrt2\right)^8 \)}{}{}}
\LANGcs{
\Question{
Které z daných čísel náleží intervalu \( \langle -5;5 \rangle\)?}
{\( \left(\sqrt2\right)^4-\left(\sqrt3\right)^4 \)}{\( 2\left(\sqrt{0{,}1}\right)^2\cdot\left(\sqrt2\right)^{10} \)}{\( \left(2\sqrt{5}\right)^2-\left(\sqrt3\right)^6 \)}{\( 2\left(\sqrt{0{,}1}\right)^4+\left(\sqrt2\right)^8 \)}{}{}}
\LANGsk{
\Question{
Ktoré z daných čísel patrí do intervalu \( \langle -5;5 \rangle\)?}
{\( \left(\sqrt2\right)^4-\left(\sqrt3\right)^4 \)}{\( 2\left(\sqrt{0{,}1}\right)^2\cdot\left(\sqrt2\right)^{10} \)}{\( \left(2\sqrt{5}\right)^2-\left(\sqrt3\right)^6 \)}{\( 2\left(\sqrt{0{,}1}\right)^4+\left(\sqrt2\right)^8 \)}{}{}}
\LANGes{
\Question{
¿Cuál de los siguientes números pertenece al intervalo \( [ -5;5 ]\)?}
{\( \left(\sqrt2\right)^4-\left(\sqrt3\right)^4 \)}{\( 2\left(\sqrt{0{.}1}\right)^2\cdot\left(\sqrt2\right)^{10} \)}{\( \left(2\sqrt{5}\right)^2-\left(\sqrt3\right)^6 \)}{\( 2\left(\sqrt{0{.}1}\right)^4+\left(\sqrt2\right)^8 \)}{}{}}
\End

\Data\ID{9000010503}
\Updated{2020-08-28 01:08:30}
\Level{B}
\SubArea{Expressions with powers and roots}
\LANGen{
\Question{
For \(x\in \mathbb{R}\),
\(x > 0\),
simplify the following expression. 
\[ 
 \root{5}\of{x}\cdot \root{}\of{x}
\]}
{\(\root{10}\of{x^{7}}\)}{\(\root{10}\of{x}\)}{\(\root{5}\of{x^{2}}\)}{\(\root{10}\of{x^{2}}\)}{}{}}
\LANGpl{
\Question{
Dla \(x\in \mathbb{R}\),
\(x > 0\),
uprość następujące wyrażenie. 
\[ 
 \root{5}\of{x}\cdot \root{}\of{x}
\]}
{\(\root{10}\of{x^{7}}\)}{\(\root{10}\of{x}\)}{\(\root{5}\of{x^{2}}\)}{\(\root{10}\of{x^{2}}\)}{}{}}
\LANGcs{
\Question{
Je-li \(x\) kladné reálné
číslo, pak je výraz \(\root{5}\of{x}\cdot \root{}\of{x}\) 
roven:}
{\(\root{10}\of{x^{7}}\)}{\(\root{10}\of{x}\)}{\(\root{5}\of{x^{2}}\)}{\(\root{10}\of{x^{2}}\)}{}{}}
\LANGsk{
\Question{
Ak je \(x\) kladné reálne
číslo, potom je výraz \(\root{5}\of{x}\cdot \root{}\of{x}\) 
rovný:}
{\(\root{10}\of{x^{7}}\)}{\(\root{10}\of{x}\)}{\(\root{5}\of{x^{2}}\)}{\(\root{10}\of{x^{2}}\)}{}{}}
\LANGes{
\Question{
Para \(x\in \mathbb{R}\),
\(x > 0\),
simplifica la siguiente expresión: 
\[ 
 \root{5}\of{x}\cdot \root{}\of{x}
\]}
{\(\root{10}\of{x^{7}}\)}{\(\root{10}\of{x}\)}{\(\root{5}\of{x^{2}}\)}{\(\root{10}\of{x^{2}}\)}{}{}}
\End

\Data\ID{9000010505}
\Updated{2022-04-23 05:04:14}
\Level{B}
\SubArea{Expressions with powers and roots}
\LANGen{
\Question{
For \(x\in \mathbb{R}\),
\(x > 0\),
simplify the following expression. 
\[ 
 \root{5}\of{x^{3}} : \root{3}\of{x}
\]}
{\(\root{15}\of{x^{4}}\)}{\(\root{5}\of{x}\)}{\(\root{3}\of{x^{2}}\)}{\(\root{5}\of{x^{2}}\)}{}{}}
\LANGpl{
\Question{
Dla \(x\in \mathbb{R}\),
\(x > 0\),
uprość następujące wyrażenie. 
\[ 
 \root{5}\of{x^{3}} : \root{3}\of{x}
\]}
{\(\root{15}\of{x^{4}}\)}{\(\root{5}\of{x}\)}{\(\root{3}\of{x^{2}}\)}{\(\root{5}\of{x^{2}}\)}{}{}}
\LANGcs{
\Question{
Je-li \(x\) kladné reálné
číslo, pak je výraz \(\root{5}\of{x^{3}} : \root{3}\of{x}\) 
roven:}
{\(\root{15}\of{x^{4}}\)}{\(\root{5}\of{x}\)}{\(\root{3}\of{x^{2}}\)}{\(\root{5}\of{x^{2}}\)}{}{}}
\LANGsk{
\Question{
Ak je \(x\) kladné reálne
číslo, potom je výraz \(\root{5}\of{x^{3}} : \root{3}\of{x}\) 
rovný:}
{\(\root{15}\of{x^{4}}\)}{\(\root{5}\of{x}\)}{\(\root{3}\of{x^{2}}\)}{\(\root{5}\of{x^{2}}\)}{}{}}
\LANGes{
\Question{
Para  \(x\in \mathbb{R}\),
\(x > 0\),
simplifica la siguiente expresión.
\[ 
 \root{5}\of{x^{3}} : \root{3}\of{x}
\]}
{\(\root{15}\of{x^{4}}\)}{\(\root{5}\of{x}\)}{\(\root{3}\of{x^{2}}\)}{\(\root{5}\of{x^{2}}\)}{}{}}
\End

\Data\ID{9000010506}
\Updated{2020-08-28 01:08:57}
\Level{B}
\SubArea{Expressions with powers and roots}
\LANGen{
\Question{
For \(x\in \mathbb{R}\),
\(x > 0\),
simplify the following expression. 
\[ 
 x\cdot \root{}\of{x}\cdot \root{3}\of{x}
\]}
{\(x\root{6}\of{x^{5}}\)}{\(\root{6}\of{x^{3}}\)}{\(\root{}\of{x}\)}{\(x^{5}\root{6}\of{x^{5}}\)}{}{}}
\LANGpl{
\Question{
Dla \(x\in \mathbb{R}\),
\(x > 0\),
uprość następujące wyrażenie. 
\[ 
 x\cdot \root{}\of{x}\cdot \root{3}\of{x}
\]}
{\(x\root{6}\of{x^{5}}\)}{\(\root{6}\of{x^{3}}\)}{\(\root{}\of{x}\)}{\(x^{5}\root{6}\of{x^{5}}\)}{}{}}
\LANGcs{
\Question{
Je-li \(x\) kladné reálné
číslo, pak je výraz \(x\cdot \root{}\of{x}\cdot \root{3}\of{x}\) 
roven:}
{\(x\root{6}\of{x^{5}}\)}{\(\root{6}\of{x^{3}}\)}{\(\root{}\of{x}\)}{\(x^{5}\root{6}\of{x^{5}}\)}{}{}}
\LANGsk{
\Question{
Ak je \(x\) kladné reálne
číslo, potom je výraz \(x\cdot \root{}\of{x}\cdot \root{3}\of{x}\) 
rovný:}
{\(x\root{6}\of{x^{5}}\)}{\(\root{6}\of{x^{3}}\)}{\(\root{}\of{x}\)}{\(x^{5}\root{6}\of{x^{5}}\)}{}{}}
\LANGes{
\Question{
Para \(x\in \mathbb{R}\),
\(x > 0\),
simplifica la siguiente expresión.
\[ 
 x\cdot \root{}\of{x}\cdot \root{3}\of{x}
\]}
{\(x\root{6}\of{x^{5}}\)}{\(\root{6}\of{x^{3}}\)}{\(\root{}\of{x}\)}{\(x^{5}\root{6}\of{x^{5}}\)}{}{}}
\End

\Data\ID{9000010508}
\Updated{2022-04-23 05:04:14}
\Level{B}
\SubArea{Expressions with powers and roots}
\LANGen{
\Question{
For \(x\in \mathbb{R}\),
\(x > 0\),
simplify the following expression. 
\[ 
 \root{3}\of{x^{2}}\cdot \root{5}\of{x^{4}}
\]}
{\(\root{15}\of{x^{22}}\)}{\(\root{15}\of{x^{6}}\)}{\(\root{15}\of{x^{8}}\)}{\(x^{3}\root{15}\of{x}\)}{}{}}
\LANGpl{
\Question{
Dla \(x\in \mathbb{R}\),
\(x > 0\),
uprość następujące wyrażenie. 
\[ 
 \root{3}\of{x^{2}}\cdot \root{5}\of{x^{4}}
\]}
{\(\root{15}\of{x^{22}}\)}{\(\root{15}\of{x^{6}}\)}{\(\root{15}\of{x^{8}}\)}{\(x^{3}\root{15}\of{x}\)}{}{}}
\LANGcs{
\Question{
Je-li \(x\) kladné reálné
číslo, pak je výraz \(\root{3}\of{x^{2}}\cdot \root{5}\of{x^{4}}\) 
roven:}
{\(\root{15}\of{x^{22}}\)}{\(\root{15}\of{x^{6}}\)}{\(\root{15}\of{x^{8}}\)}{\(x^{3}\root{15}\of{x}\)}{}{}}
\LANGsk{
\Question{
Ak je \(x\) kladné reálne
číslo, potom je výraz \(\root{3}\of{x^{2}}\cdot \root{5}\of{x^{4}}\) 
rovný:}
{\(\root{15}\of{x^{22}}\)}{\(\root{15}\of{x^{6}}\)}{\(\root{15}\of{x^{8}}\)}{\(x^{3}\root{15}\of{x}\)}{}{}}
\LANGes{
\Question{
Para \(x\in \mathbb{R}\),
\(x > 0\),
simplifica la siguiente expresión. 
\[ 
 \root{3}\of{x^{2}}\cdot \root{5}\of{x^{4}}
\]}
{\(\root{15}\of{x^{22}}\)}{\(\root{15}\of{x^{6}}\)}{\(\root{15}\of{x^{8}}\)}{\(x^{3}\root{15}\of{x}\)}{}{}}
\End

\Data\ID{9000010510}
\Updated{2020-08-28 01:08:06}
\Level{B}
\SubArea{Expressions with powers and roots}
\LANGen{
\Question{
For \(x\in \mathbb{R}\),
\(x > 0\),
simplify the following expression. 
\[ 
 \root{3}\of{x} : \root{6}\of{x}
\]}
{\(\root{6}\of{x}\)}{\(\root{}\of{x}\)}{\(\root{3}\of{x^{2}}\)}{\(x\)}{}{}}
\LANGpl{
\Question{
Dla \(x\in \mathbb{R}\),
\(x > 0\),
uprość następujące wyrażenie. 
\[ 
 \root{3}\of{x} : \root{6}\of{x}
\]}
{\(\root{6}\of{x}\)}{\(\root{}\of{x}\)}{\(\root{3}\of{x^{2}}\)}{\(x\)}{}{}}
\LANGcs{
\Question{
Je-li \(x\) kladné reálné
číslo, pak je výraz \(\root{3}\of{x} : \root{6}\of{x}\) 
roven:}
{\(\root{6}\of{x}\)}{\(\root{}\of{x}\)}{\(\root{3}\of{x^{2}}\)}{\(x\)}{}{}}
\LANGsk{
\Question{
Ak je \(x\) kladné reálne
číslo, potom je výraz \(\root{3}\of{x} : \root{6}\of{x}\) 
rovný:}
{\(\root{6}\of{x}\)}{\(\root{}\of{x}\)}{\(\root{3}\of{x^{2}}\)}{\(x\)}{}{}}
\LANGes{
\Question{
Para \(x\in \mathbb{R}\),
\(x > 0\),
simplifica la siguiente expresión.
\[ 
 \root{3}\of{x} : \root{6}\of{x}
\]}
{\(\root{6}\of{x}\)}{\(\root{}\of{x}\)}{\(\root{3}\of{x^{2}}\)}{\(x\)}{}{}}
\End

\Data\ID{9000013501}
\Updated{2020-08-28 01:08:10}
\Level{B}
\SubArea{Expressions with powers and roots}
\LANGen{
\Question{
Write the expression \(2^{\frac{3} 
{4} }\) 
in an equivalent form which does not contain a rational exponent.}
{\(\root{4}\of{2^{3}}\)}{\(\root{4}\of{2}\)}{\(\root{3}\of{2^{4}}\)}{\(\root{4}\of{3^{2}}\)}{}{}}
\LANGpl{
\Question{
Zapisz wyrażenie \(2^{\frac{3} 
{4} }\) 
w odpowiedniej formie, tak by nie zawierało wykładnika wymiernego.}
{\(\root{4}\of{2^{3}}\)}{\(\root{4}\of{2}\)}{\(\root{3}\of{2^{4}}\)}{\(\root{4}\of{3^{2}}\)}{}{}}
\LANGcs{
\Question{
Číslo \(2^{\frac{3} 
{4} }\) 
zapište jako odmocninu.}
{\(\root{4}\of{2^{3}}\)}{\(\root{4}\of{2}\)}{\(\root{3}\of{2^{4}}\)}{\(\root{4}\of{3^{2}}\)}{}{}}
\LANGsk{
\Question{
Číslo \(2^{\frac{3} 
{4} }\) 
zapíšte ako odmocninu.}
{\(\root{4}\of{2^{3}}\)}{\(\root{4}\of{2}\)}{\(\root{3}\of{2^{4}}\)}{\(\root{4}\of{3^{2}}\)}{}{}}
\LANGes{
\Question{
Escribe el número \(2^{\frac{3} 
{4} }\) 
en forma de raíz.}
{\(\root{4}\of{2^{3}}\)}{\(\root{4}\of{2}\)}{\(\root{3}\of{2^{4}}\)}{\(\root{4}\of{3^{2}}\)}{}{}}
\End

\Data\ID{9000013502}
\Updated{2022-04-23 05:04:14}
\Level{B}
\SubArea{Expressions with powers and roots}
\LANGen{
\Question{
Simplify the expression \(0.5^{\frac{6} 
{7} }\cdot 0.5^{-\frac{5} 
{14} }\) 
and write the result using a root.}
{\(\sqrt{0.5}\)}{\(\root{7}\of{0.5}\)}{\(\root{14}\of{0.5^{11}}\)}{\(\root{14}\of{0.5}\)}{}{}}
\LANGpl{
\Question{
Uprość wyrażenie \(0.5^{\frac{6} 
{7} }\cdot 0.5^{-\frac{5} 
{14} }\) 
i zapisz je w odpowiedniej formie, tak by nie zawierało wykładnika wymiernego.}
{\(\sqrt{0.5}\)}{\(\root{7}\of{0.5}\)}{\(\root{14}\of{0.5^{11}}\)}{\(\root{14}\of{0.5}\)}{}{}}
\LANGcs{
\Question{
Číslo \(0{,}5^{\frac{6} 
{7} }\cdot 0{,}5^{-\frac{5} 
{14} }\) 
zjednodušte a výsledek zapište jako odmocninu.}
{\(\sqrt{0, 5}\)}{\(\root{7}\of{0{,}5}\)}{\(\root{14}\of{0{,}5^{11}}\)}{\(\root{14}\of{0{,}5}\)}{}{}}
\LANGsk{
\Question{
Číslo \(0{,}5^{\frac{6} 
{7} }\cdot 0{,}5^{-\frac{5} 
{14} }\) 
zjednodušte a výsledok zapíšte ako odmocninu.}
{\(\sqrt{0, 5}\)}{\(\root{7}\of{0{,}5}\)}{\(\root{14}\of{0{,}5^{11}}\)}{\(\root{14}\of{0{,}5}\)}{}{}}
\LANGes{
\Question{
Simplifica el número \(0{,}5^{\frac{6} 
{7} }\cdot 0{,}5^{-\frac{5} 
{14} }\) 
y escríbelo en forma de raíz.}
{\(\sqrt{0, 5}\)}{\(\root{7}\of{0{,}5}\)}{\(\root{14}\of{0{,}5^{11}}\)}{\(\root{14}\of{0{,}5}\)}{}{}}
\End

\Data\ID{9000013503}
\Updated{2022-04-23 05:04:14}
\Level{B}
\SubArea{Expressions with powers and roots}
\LANGen{
\Question{
Write the number \(\root{6}\of{3^{-3}}\) 
as a power with a rational exponent.}
{\(3^{-\frac{1} 
{2} }\)}{\(3^{\frac{1} 
{2} }\)}{\(3^{2}\)}{\(3^{-2}\)}{}{}}
\LANGpl{
\Question{
Zapisz liczbę  \(\root{6}\of{3^{-3}}\) 
jako potęgę z wykładnikiem wymiernym.}
{\(3^{-\frac{1} 
{2} }\)}{\(3^{\frac{1} 
{2} }\)}{\(3^{2}\)}{\(3^{-2}\)}{}{}}
\LANGcs{
\Question{
Číslo \(\root{6}\of{3^{-3}}\) 
zapište jako mocninu.}
{\(3^{-\frac{1} 
{2} }\)}{\(3^{\frac{1} 
{2} }\)}{\(3^{2}\)}{\(3^{-2}\)}{}{}}
\LANGsk{
\Question{
Číslo \(\root{6}\of{3^{-3}}\) 
zapíšte ako mocninu.}
{\(3^{-\frac{1} 
{2} }\)}{\(3^{\frac{1} 
{2} }\)}{\(3^{2}\)}{\(3^{-2}\)}{}{}}
\LANGes{
\Question{
Escribe el número\(\root{6}\of{3^{-3}}\) 
en forma de potencia con exponente racional.}
{\(3^{-\frac{1} 
{2} }\)}{\(3^{\frac{1} 
{2} }\)}{\(3^{2}\)}{\(3^{-2}\)}{}{}}
\End

\Data\ID{9000013504}
\Updated{2020-08-28 01:08:22}
\Level{B}
\SubArea{Expressions with powers and roots}
\LANGen{
\Question{
Simplify \(\sqrt{\root{4}\of{25}}\).}
{\(\root{4}\of{5}\)}{\(\root{8}\of{5}\)}{\(\root{4}\of{25}\)}{\(\sqrt{5}\)}{}{}}
\LANGpl{
\Question{
Uprość \(\sqrt{\root{4}\of{25}}\).}
{\(\root{4}\of{5}\)}{\(\root{8}\of{5}\)}{\(\root{4}\of{25}\)}{\(\sqrt{5}\)}{}{}}
\LANGcs{
\Question{
Číslo \(\sqrt{\root{4}\of{25}}\) 
zjednodušte a zapište jako odmocninu.}
{\(\root{4}\of{5}\)}{\(\root{8}\of{5}\)}{\(\root{4}\of{25}\)}{\(\sqrt{5}\)}{}{}}
\LANGsk{
\Question{
Číslo \(\sqrt{\root{4}\of{25}}\) 
zjednodušte a zapíšte ako odmocninu.}
{\(\root{4}\of{5}\)}{\(\root{8}\of{5}\)}{\(\root{4}\of{25}\)}{\(\sqrt{5}\)}{}{}}
\LANGes{
\Question{
Simplifica \(\sqrt{\root{4}\of{25}}\).}
{\(\root{4}\of{5}\)}{\(\root{8}\of{5}\)}{\(\root{4}\of{25}\)}{\(\sqrt{5}\)}{}{}}
\End

\Data\ID{9000013507}
\Updated{2020-08-28 01:08:07}
\Level{B}
\SubArea{Expressions with powers and roots}
\LANGen{
\Question{
Simplify \(\root{4}\of{16\cdot 81}\).}
{\(6\)}{\(36\)}{\(\sqrt{6}\)}{\(\root{4}\of{6}\)}{}{}}
\LANGpl{
\Question{
Uprość \(\root{4}\of{16\cdot 81}\).}
{\(6\)}{\(36\)}{\(\sqrt{6}\)}{\(\root{4}\of{6}\)}{}{}}
\LANGcs{
\Question{
Vypočítejte \(\root{4}\of{16\cdot 81}\).}
{\(6\)}{\(36\)}{\(\sqrt{6}\)}{\(\root{4}\of{6}\)}{}{}}
\LANGsk{
\Question{
Vypočítajte \(\root{4}\of{16\cdot 81}\).}
{\(6\)}{\(36\)}{\(\sqrt{6}\)}{\(\root{4}\of{6}\)}{}{}}
\LANGes{
\Question{
Simplifica \(\root{4}\of{16\cdot 81}\).}
{\(6\)}{\(36\)}{\(\sqrt{6}\)}{\(\root{4}\of{6}\)}{}{}}
\End

\Data\ID{9000079207}
\Updated{2022-04-23 05:04:14}
\Level{B}
\SubArea{Expressions with powers and roots}
\LANGen{
\Question{
Assuming \(x\not \in \{0;1;3\}\),
simplify the following expression. 
\[ 
 \frac{x^{2} - 9} 
{x^{2} - x}\cdot \left (\frac{x^{2} - 3x}
 {x - 1} \right )^{-1}
\]}
{\(\frac{x+3}
 {x^{2}} \)}{\(\frac{x-3}
 {x^{2}} \)}{\(\frac{x+3}
 {2x} \)}{\(\frac{x+3}
 {x} \)}{}{}}
\LANGpl{
\Question{
Zakładając, że  \(x\not \in \{0;1;3\}\),
 uprość podane wyrażenie:
\[ 
 \frac{x^{2} - 9} 
{x^{2} - x}\cdot \left (\frac{x^{2} - 3x}
 {x - 1} \right )^{-1}
\]}
{\(\frac{x+3}
 {x^{2}} \)}{\(\frac{x-3}
 {x^{2}} \)}{\(\frac{x+3}
 {2x} \)}{\(\frac{x+3}
 {x} \)}{}{}}
\LANGcs{
\Question{
Za předpokladu, že \(x\not \in \{0;1;3\}\), upravte na co
nejjednodušší tvar výraz \(\frac{x^{2}-9}
{x^{2}-x}\cdot \left (\frac{x^{2}-3x}
 {x-1} \right )^{-1}\).}
{\(\frac{x+3}
 {x^{2}} \)}{\(\frac{x-3}
 {x^{2}} \)}{\(\frac{x+3}
 {2x} \)}{\(\frac{x+3}
 {x} \)}{}{}}
\LANGsk{
\Question{
Za predpokladu, že \(x\not \in \{0;1;3\}\), zjednodušte
daný výraz \(\frac{x^{2}-9}
{x^{2}-x}\cdot \left (\frac{x^{2}-3x}
 {x-1} \right )^{-1}\).}
{\(\frac{x+3}
 {x^{2}} \)}{\(\frac{x-3}
 {x^{2}} \)}{\(\frac{x+3}
 {2x} \)}{\(\frac{x+3}
 {x} \)}{}{}}
\LANGes{
\Question{
Suponiendo \(x\not \in \{0;1;3\}\),
simplifica la siguiente expresión. 
\[ 
 \frac{x^{2} - 9} 
{x^{2} - x}\cdot \left (\frac{x^{2} - 3x}
 {x - 1} \right )^{-1}
\]}
{\(\frac{x+3}
 {x^{2}} \)}{\(\frac{x-3}
 {x^{2}} \)}{\(\frac{x+3}
 {2x} \)}{\(\frac{x+3}
 {x} \)}{}{}}
\End

\Data\ID{1003034108}
\Updated{2022-08-25 10:08:00}
\Level{C}
\SubArea{Expressions with powers and roots}
\LANGen{
\Question{
Give the result of the gradual exponentiation \(\left(\left(\left(\left(2\right)^2\right)^0\right)^3\right)^4\):}
{\( 1 \)}{\(  0 \)}{\( 4^{12} \)}{\( 2^9 \)}{}{}}
\LANGpl{
\Question{
Liczba \(\left(\left(\left(\left(2\right)^2\right)^0\right)^3\right)^4\) jest równa:}
{\( 1 \)}{\(  0 \)}{\( 4^{12} \)}{\( 2^9 \)}{}{}}
\LANGcs{
\Question{
Umocněním výrazu \(\left(\left(\left(\left(2\right)^2\right)^0\right)^3\right)^4\) dostaneme:}
{\( 1 \)}{\(  0 \)}{\( 4^{12} \)}{\( 2^9 \)}{}{}}
\LANGsk{
\Question{
Umocnením výrazu \(\left(\left(\left(\left(2\right)^2\right)^0\right)^3\right)^4\) dostaneme:}
{\( 1 \)}{\(  0 \)}{\( 4^{12} \)}{\( 2^9 \)}{}{}}
\LANGes{
\Question{
Halla el resultado de la exponenciación gradual \(\left(\left(\left(\left(2\right)^2\right)^0\right)^3\right)^4\):}
{\( 1 \)}{\(  0 \)}{\( 4^{12} \)}{\( 2^9 \)}{}{}}
\End

\Data\ID{1003109205}
\Updated{2020-10-12 10:10:40}
\Level{C}
\SubArea{Expressions with powers and roots}
\LANGen{
\Question{
The bank offers a term deposit with annual interest rate of \( 0.9\% \). We need to tax our profits with the income tax of \( 15\% \). How many Euros are we going to have after five years of savings if the original deposit is \( 200 \) Euros?}
{\( 207.77 \)}{\( 209.36 \)}{\( 371.8 \)}{\( 209.16 \)}{\( 207.65 \)}{}}
\LANGpl{
\Question{
Bank oferuje lokatę z rocznym oprocentowaniem w wysokości  \( 0{,}9\% \). Podatek dochodowy od zysku wynosi \( 15\% \). Ile Euro zyskamy po pięciu latach oszczędzania jeśli początkowy wkład wynosi  \( 200 \) Euro?}
{\( 207{,}77 \)}{\( 209{,}36 \)}{\( 371{,}8 \)}{\( 209{,}16 \)}{\( 207{,}65 \)}{}}
\LANGcs{
\Question{
Banka nabízí na termínované vklady úrok \( 0{,}9\% \) s roční úrokovou mírou. Daň z příjmu je \( 15\% \). Kolik euro budeme mít po pěti letech spoření, vložíme-li \( 200 \) euro?}
{\( 207{,}77 \)}{\( 209{,}36 \)}{\( 371{,}8 \)}{\( 209{,}16 \)}{\( 207{,}65 \)}{}}
\LANGsk{
\Question{
Banka ponúka na termínovaných vkladoch úrok \( 0{,}9\% \) s ročnou úrokovou mierou. Daň z príjmu je \( 15\% \). Koľko euro budeme mať po piatich rokoch sporenia, ak vložíme \( 200 \) euro?}
{\( 207{,}77 \)}{\( 209{,}36 \)}{\( 371{,}8 \)}{\( 209{,}16 \)}{\( 207{,}65 \)}{}}
\LANGes{
\Question{
Un banco ofrece depósitos a plazos con una tasa de interés anual del \( 0.9\% \). Necesitamos declarar nuestra ganancia con un impuesto de \( 15\% \). ¿Cuántos euros tendremos después de cinco años si el depósito original fue de \( 200 \) Euros?}
{\( 207.77 \)}{\( 209.36 \)}{\( 371.8 \)}{\( 209.16 \)}{\( 207.65 \)}{}}
\End

\Data\ID{1003118001}
\Updated{2020-08-27 08:08:19}
\Level{C}
\SubArea{Expressions with powers and roots}
\LANGen{
\Question{
The number \( 3\sqrt[3]3\cdot\sqrt[4]{3\sqrt[3]3}\cdot\sqrt[5]{3\sqrt[3]3\cdot\sqrt[4]{3\sqrt[3]3}} \) is equal to:}
{\( 9 \)}{\( 3 \)}{\( \sqrt3 \)}{\( \sqrt[3]3 \)}{}{}}
\LANGpl{
\Question{
Liczba \( 3\sqrt[3]3\cdot\sqrt[4]{3\sqrt[3]3}\cdot\sqrt[5]{3\sqrt[3]3\cdot\sqrt[4]{3\sqrt[3]3}} \) jest równa:}
{\( 9 \)}{\( 3 \)}{\( \sqrt3 \)}{\( \sqrt[3]3 \)}{}{}}
\LANGcs{
\Question{
Číslo \( 3\sqrt[3]3\cdot\sqrt[4]{3\sqrt[3]3}\cdot\sqrt[5]{3\sqrt[3]3\cdot\sqrt[4]{3\sqrt[3]3}} \) se rovná:}
{\( 9 \)}{\( 3 \)}{\( \sqrt3 \)}{\( \sqrt[3]3 \)}{}{}}
\LANGsk{
\Question{
Číslo \( 3\sqrt[3]3\cdot\sqrt[4]{3\sqrt[3]3}\cdot\sqrt[5]{3\sqrt[3]3\cdot\sqrt[4]{3\sqrt[3]3}} \) sa rovná:}
{\( 9 \)}{\( 3 \)}{\( \sqrt3 \)}{\( \sqrt[3]3 \)}{}{}}
\LANGes{
\Question{
El número \( 3\sqrt[3]3\cdot\sqrt[4]{3\sqrt[3]3}\cdot\sqrt[5]{3\sqrt[3]3\cdot\sqrt[4]{3\sqrt[3]3}} \) es igual a:}
{\( 9 \)}{\( 3 \)}{\( \sqrt3 \)}{\( \sqrt[3]3 \)}{}{}}
\End

\Data\ID{1003118002}
\Updated{2020-08-27 08:08:42}
\Level{C}
\SubArea{Expressions with powers and roots}
\LANGen{
\Question{
The number \( \sqrt{7-4\sqrt3} + \sqrt{7+4\sqrt3} \) is equal to:}
{\( 4 \)}{\( \sqrt{14} \)}{\( 16 \)}{\( 8\sqrt3 \)}{}{}}
\LANGpl{
\Question{
Liczba \( \sqrt{7-4\sqrt3} + \sqrt{7+4\sqrt3} \) jest równa:}
{\( 4 \)}{\( \sqrt{14} \)}{\( 16 \)}{\( 8\sqrt3 \)}{}{}}
\LANGcs{
\Question{
Číslo \( \sqrt{7-4\sqrt3} + \sqrt{7+4\sqrt3} \) se rovná:}
{\( 4 \)}{\( \sqrt{14} \)}{\( 16 \)}{\( 8\sqrt3 \)}{}{}}
\LANGsk{
\Question{
Číslo \( \sqrt{7-4\sqrt3} + \sqrt{7+4\sqrt3} \) sa rovná:}
{\( 4 \)}{\( \sqrt{14} \)}{\( 16 \)}{\( 8\sqrt3 \)}{}{}}
\LANGes{
\Question{
El número \( \sqrt{7-4\sqrt3} + \sqrt{7+4\sqrt3} \) es igual a:}
{\( 4 \)}{\( \sqrt{14} \)}{\( 16 \)}{\( 8\sqrt3 \)}{}{}}
\End

\Data\ID{1003118003}
\Updated{2020-10-12 10:10:13}
\Level{C}
\SubArea{Expressions with powers and roots}
\LANGen{
\Question{
Give the multiplicative inverse of \( \frac{\sqrt[3]4-\sqrt[3]2}2 \).}
{\( 2\sqrt[3]2 + 2 +\sqrt[3]4 \)}{\( \sqrt[3]{12}+2+\sqrt[3]4 \)}{\( 2\sqrt[3]{12}+2+\sqrt[3]4 \)}{\( \sqrt[3]{2}+2+\sqrt[3]4 \)}{}{}}
\LANGpl{
\Question{
Liczbą odwrotną do \( \frac{\sqrt[3]4-\sqrt[3]2}2 \) jest:}
{\( 2\sqrt[3]2 + 2 +\sqrt[3]4 \)}{\( \sqrt[3]{12}+2+\sqrt[3]4 \)}{\( 2\sqrt[3]{12}+2+\sqrt[3]4 \)}{\( \sqrt[3]{2}+2+\sqrt[3]4 \)}{}{}}
\LANGcs{
\Question{
Určete převrácený výraz k výrazu \( \frac{\sqrt[3]4-\sqrt[3]2}2 \).}
{\( 2\sqrt[3]2 + 2 +\sqrt[3]4 \)}{\( \sqrt[3]{12}+2+\sqrt[3]4 \)}{\( 2\sqrt[3]{12}+2+\sqrt[3]4 \)}{\( \sqrt[3]{2}+2+\sqrt[3]4 \)}{}{}}
\LANGsk{
\Question{
Určte prevrátený výraz k výrazu \( \frac{\sqrt[3]4-\sqrt[3]2}2 \).}
{\( 2\sqrt[3]2 + 2 +\sqrt[3]4 \)}{\( \sqrt[3]{12}+2+\sqrt[3]4 \)}{\( 2\sqrt[3]{12}+2+\sqrt[3]4 \)}{\( \sqrt[3]{2}+2+\sqrt[3]4 \)}{}{}}
\LANGes{
\Question{
Halla el inverso multiplicativo de \( \frac{\sqrt[3]4-\sqrt[3]2}2 \).}
{\( 2\sqrt[3]2 + 2 +\sqrt[3]4 \)}{\( \sqrt[3]{12}+2+\sqrt[3]4 \)}{\( 2\sqrt[3]{12}+2+\sqrt[3]4 \)}{\( \sqrt[3]{2}+2+\sqrt[3]4 \)}{}{}}
\End

\Data\ID{1003118004}
\Updated{2020-08-27 08:08:38}
\Level{C}
\SubArea{Expressions with powers and roots}
\LANGen{
\Question{
The number \( \sqrt{3-2\sqrt2} \) equals:}
{\( \sqrt2-1 \)}{\( 1-\sqrt2 \)}{\( \sqrt2 \)}{\( \sqrt3-\sqrt{2\sqrt2} \)}{}{}}
\LANGpl{
\Question{
Liczba \( \sqrt{3-2\sqrt2} \)  jest równa:}
{\( \sqrt2-1 \)}{\( 1-\sqrt2 \)}{\( \sqrt2 \)}{\( \sqrt3-\sqrt{2\sqrt2} \)}{}{}}
\LANGcs{
\Question{
Číslo\( \sqrt{3-2\sqrt2} \) se rovná:}
{\( \sqrt2-1 \)}{\( 1-\sqrt2 \)}{\( \sqrt2 \)}{\( \sqrt3-\sqrt{2\sqrt2} \)}{}{}}
\LANGsk{
\Question{
Číslo \( \sqrt{3-2\sqrt2} \) sa rovná:}
{\( \sqrt2-1 \)}{\( 1-\sqrt2 \)}{\( \sqrt2 \)}{\( \sqrt3-\sqrt{2\sqrt2} \)}{}{}}
\LANGes{
\Question{
El número \( \sqrt{3-2\sqrt2} \) es igual a:}
{\( \sqrt2-1 \)}{\( 1-\sqrt2 \)}{\( \sqrt2 \)}{\( \sqrt3-\sqrt{2\sqrt2} \)}{}{}}
\End

\Data\ID{1003118006}
\Updated{2020-08-27 09:08:09}
\Level{C}
\SubArea{Expressions with powers and roots}
\LANGen{
\Question{
Give the value of the expression \( \sqrt{\left(2-\sqrt7\right)^2}-\sqrt{\left(3+\sqrt7\right)^2} \).}
{\( -5 \)}{\( -1 \)}{\( -1-2\sqrt7 \)}{\( -5+2\sqrt7 \)}{}{}}
\LANGpl{
\Question{
Wartość wyrażenia \( \sqrt{\left(2-\sqrt7\right)^2}-\sqrt{\left(3+\sqrt7\right)^2} \) jest równa:}
{\( -5 \)}{\( -1 \)}{\( -1-2\sqrt7 \)}{\( -5+2\sqrt7 \)}{}{}}
\LANGcs{
\Question{
Určete hodnotu výrazu \( \sqrt{\left(2-\sqrt7\right)^2}-\sqrt{\left(3+\sqrt7\right)^2} \).}
{\( -5 \)}{\( -1 \)}{\( -1-2\sqrt7 \)}{\( -5+2\sqrt7 \)}{}{}}
\LANGsk{
\Question{
Určte hodnotu výrazu \( \sqrt{\left(2-\sqrt7\right)^2}-\sqrt{\left(3+\sqrt7\right)^2} \).}
{\( -5 \)}{\( -1 \)}{\( -1-2\sqrt7 \)}{\( -5+2\sqrt7 \)}{}{}}
\LANGes{
\Question{
Halla el valor de la expesión \( \sqrt{\left(2-\sqrt7\right)^2}-\sqrt{\left(3+\sqrt7\right)^2} \).}
{\( -5 \)}{\( -1 \)}{\( -1-2\sqrt7 \)}{\( -5+2\sqrt7 \)}{}{}}
\End

\Data\ID{1003118007}
\Updated{2020-08-27 09:08:11}
\Level{C}
\SubArea{Expressions with powers and roots}
\LANGen{
\Question{
Calculate the fourth power of the number \( 1-\sqrt2 \).}
{\( 17-12\sqrt2 \)}{\( 17-4\sqrt2 \)}{\( 3-2\sqrt2 \)}{\( 9-4\sqrt2 \)}{}{}}
\LANGpl{
\Question{
Czwarta potęga liczby \( 1-\sqrt2 \) jest równa:}
{\( 17-12\sqrt2 \)}{\( 17-4\sqrt2 \)}{\( 3-2\sqrt2 \)}{\( 9-4\sqrt2 \)}{}{}}
\LANGcs{
\Question{
Vypočítejte čtvrtou mocninu čísla \( 1-\sqrt2 \).}
{\( 17-12\sqrt2 \)}{\( 17-4\sqrt2 \)}{\( 3-2\sqrt2 \)}{\( 9-4\sqrt2 \)}{}{}}
\LANGsk{
\Question{
Vypočítajte štvrtú mocninu čísla \( 1-\sqrt2 \).}
{\( 17-12\sqrt2 \)}{\( 17-4\sqrt2 \)}{\( 3-2\sqrt2 \)}{\( 9-4\sqrt2 \)}{}{}}
\LANGes{
\Question{
Calcula la cuarta potencia del número \( 1-\sqrt2 \).}
{\( 17-12\sqrt2 \)}{\( 17-4\sqrt2 \)}{\( 3-2\sqrt2 \)}{\( 9-4\sqrt2 \)}{}{}}
\End

\Data\ID{1003118008}
\Updated{2020-10-12 10:10:40}
\Level{C}
\SubArea{Expressions with powers and roots}
\LANGen{
\Question{
Let \( a=\frac1{5+\sqrt3}+\frac1{44-22\sqrt3} \). Choose the correct statement about the number \( a \).}
{It is a rational number.}{It is a positive integer.}{It is an irrational number.}{It is greater than \( 1 \).}{}{}}
\LANGpl{
\Question{
Liczba \( a=\frac1{5+\sqrt3}+\frac1{44-22\sqrt3} \). Zatem liczba  \( a \) jest:}
{liczbą wymierną}{dodatnią liczbą całkowitą}{liczbą niewymierną}{liczbą większą od \( 1 \)}{}{}}
\LANGcs{
\Question{
Nechť \( a=\frac1{5+\sqrt3}+\frac1{44-22\sqrt3} \). Vyberte správné tvrzení o číslu \( a \).}
{Je to racionální číslo.}{Je to kladné celé číslo.}{Je to iracionální číslo.}{Je větší než \( 1 \).}{}{}}
\LANGsk{
\Question{
Nech \( a=\frac1{5+\sqrt3}+\frac1{44-22\sqrt3} \). Vyberte správne tvrdenie o čísle \( a \).}
{Je to racionálne číslo.}{Je to kladné celé číslo.}{Je to iracionálne číslo.}{Je väčšie ako \( 1 \).}{}{}}
\LANGes{
\Question{
Sea \( a=\frac1{5+\sqrt3}+\frac1{44-22\sqrt3} \). Elije la declaración correcta sobre el número \( a \).}
{Es un número racional.}{Es un número natural.}{Es un número irracional.}{Es mayor que \( 1 \).}{}{}}
\End

\Data\ID{1003118010}
\Updated{2020-08-27 09:08:23}
\Level{C}
\SubArea{Expressions with powers and roots}
\LANGen{
\Question{
The product of the number \( \sqrt{\sqrt2+1} \) and the multiplicative inverse  of \( \sqrt{\sqrt2-1} \) is   equal to:}
{\( 1+\sqrt2 \)}{\( 2\sqrt2 \)}{\( 1-\sqrt2 \)}{\( 1 \)}{}{}}
\LANGpl{
\Question{
Iloczyn liczby \( \sqrt{\sqrt2+1} \) i odwrotności liczby \( \sqrt{\sqrt2-1} \) jest równy:}
{\( 1+\sqrt2 \)}{\( 2\sqrt2 \)}{\( 1-\sqrt2 \)}{\( 1 \)}{}{}}
\LANGcs{
\Question{
Součin čísla \( \sqrt{\sqrt2+1} \) a převrácené hodnoty čísla \( \sqrt{\sqrt2-1} \) se rovná:}
{\( 1+\sqrt2 \)}{\( 2\sqrt2 \)}{\( 1-\sqrt2 \)}{\( 1 \)}{}{}}
\LANGsk{
\Question{
Súčin čísla \( \sqrt{\sqrt2+1} \) a prevrátenej hodnoty čísla \( \sqrt{\sqrt2-1} \) je rovný:}
{\( 1+\sqrt2 \)}{\( 2\sqrt2 \)}{\( 1-\sqrt2 \)}{\( 1 \)}{}{}}
\LANGes{
\Question{
El producto del número \( \sqrt{\sqrt2+1} \) y el inverso multiplicativo de \( \sqrt{\sqrt2-1} \) es igual a:}
{\( 1+\sqrt2 \)}{\( 2\sqrt2 \)}{\( 1-\sqrt2 \)}{\( 1 \)}{}{}}
\End

\Data\ID{1003118101}
\Updated{2020-10-12 10:10:13}
\Level{C}
\SubArea{Expressions with powers and roots}
\LANGen{
\Question{
Which of the following inequalities describes a correct relationship between the numbers \( a \), \( b \), \( c \) if \( a=4^{2^4} \), \( b=3^{4^3} \), \( c=2^{3^4} \)? (Hint: Barring parentheses, \( x^{y^z} \) is equivalent to \( x^{\left(y^z\right)} \).)}
{\( a < c < b \)}{\( b < a < c \)}{\( a < b < c \)}{\( b < c < a \)}{}{}}
\LANGpl{
\Question{
Która z poniższych  nierówności opisuje poprawną zależność między liczbami  \( a \), \( b \), \( c \) jeśli: \( a=4^{2^4} \), \( b=3^{4^3} \), \( c=2^{3^4} \)? (Wskazówka: \( x^{y^z} \) jest równe \( x^{\left(y^z\right)} \).)}
{\( a < c < b \)}{\( b < a < c \)}{\( a < b < c \)}{\( b < c < a \)}{}{}}
\LANGcs{
\Question{
Nechť \( a=4^{2^4} \), \( b=3^{4^3} \) a \( c=2^{3^4} \). Která z uvedených nerovností popisuje správně vztah mezi čísly \( a \), \( b \), \( c \)? (Nápověda:  \( x^{y^z} = x^{\left(y^z\right)} \))}
{\( a < c < b \)}{\( b < a < c \)}{\( a < b < c \)}{\( b < c < a \)}{}{}}
\LANGsk{
\Question{
Nech \( a=4^{2^4} \), \( b=3^{4^3} \), \( c=2^{3^4} \). Ktorá z uvedených nerovností popisuje správne vzťah medzi \( a \), \( b \), \( c \) ?  (Pomoc: \( x^{y^z}=x^{\left(y^z\right)} \))}
{\( a < c < b \)}{\( b < a < c \)}{\( a < b < c \)}{\( b < c < a \)}{}{}}
\LANGes{
\Question{
Sean \( a=4^{2^4} \), \( b=3^{4^3} \) y \( c=2^{3^4} \). ¿Cuál de las inecuaciones describe la relación entre \( a \), \( b \) y  \( c \)? (Pista:  \( x^{y^z} = x^{\left(y^z\right)} \))}
{\( a < c < b \)}{\( b < a < c \)}{\( a < b < c \)}{\( b < c < a \)}{}{}}
\End

\Data\ID{1003118104}
\Updated{2020-08-27 09:08:03}
\Level{C}
\SubArea{Expressions with powers and roots}
\LANGen{
\Question{
Simplify \(\sqrt[4]{\left(\sqrt3-\sqrt2\right)^4}+\sqrt[4]{\left(\sqrt2-\sqrt5\right)^4}+\sqrt[3]{\left(\sqrt3-\sqrt5\right)^3} \).}
{\( 2\sqrt3-2\sqrt2 \)}{\( 2\sqrt5 - 2\sqrt2 \)}{\( 2\sqrt3-2\sqrt5 \)}{\( 2\sqrt5-2\sqrt3 \)}{}{}}
\LANGpl{
\Question{
Liczba \(\sqrt[4]{\left(\sqrt3-\sqrt2\right)^4}+\sqrt[4]{\left(\sqrt2-\sqrt5\right)^4}+\sqrt[3]{\left(\sqrt3-\sqrt5\right)^3} \) jest równa:}
{\( 2\sqrt3-2\sqrt2 \)}{\( 2\sqrt5 - 2\sqrt2 \)}{\( 2\sqrt3-2\sqrt5 \)}{\( 2\sqrt5-2\sqrt3 \)}{}{}}
\LANGcs{
\Question{
Zjednodušte \(\sqrt[4]{\left(\sqrt3-\sqrt2\right)^4}+\sqrt[4]{\left(\sqrt2-\sqrt5\right)^4}+\sqrt[3]{\left(\sqrt3-\sqrt5\right)^3} \).}
{\( 2\sqrt3-2\sqrt2 \)}{\( 2\sqrt5 - 2\sqrt2 \)}{\( 2\sqrt3-2\sqrt5 \)}{\( 2\sqrt5-2\sqrt3 \)}{}{}}
\LANGsk{
\Question{
Zjednodušte \(\sqrt[4]{\left(\sqrt3-\sqrt2\right)^4}+\sqrt[4]{\left(\sqrt2-\sqrt5\right)^4}+\sqrt[3]{\left(\sqrt3-\sqrt5\right)^3} \).}
{\( 2\sqrt3-2\sqrt2 \)}{\( 2\sqrt5 - 2\sqrt2 \)}{\( 2\sqrt3-2\sqrt5 \)}{\( 2\sqrt5-2\sqrt3 \)}{}{}}
\LANGes{
\Question{
Simplifica \(\sqrt[4]{\left(\sqrt3-\sqrt2\right)^4}+\sqrt[4]{\left(\sqrt2-\sqrt5\right)^4}+\sqrt[3]{\left(\sqrt3-\sqrt5\right)^3} \).}
{\( 2\sqrt3-2\sqrt2 \)}{\( 2\sqrt5 - 2\sqrt2 \)}{\( 2\sqrt3-2\sqrt5 \)}{\( 2\sqrt5-2\sqrt3 \)}{}{}}
\End

\Data\ID{1003118105}
\Updated{2020-08-27 09:08:26}
\Level{C}
\SubArea{Expressions with powers and roots}
\LANGen{
\Question{
The number \( \frac1{\left(2-\sqrt3\right)^3} \) is equal to:}
{\( 26+15\sqrt3 \)}{\( 27+24\sqrt3 \)}{\( 14+7\sqrt3 \)}{\( 27+30\sqrt3 \)}{}{}}
\LANGpl{
\Question{
Liczba  \( \frac1{\left(2-\sqrt3\right)^3} \) jest równa:}
{\( 26+15\sqrt3 \)}{\( 27+24\sqrt3 \)}{\( 14+7\sqrt3 \)}{\( 27+30\sqrt3 \)}{}{}}
\LANGcs{
\Question{
Číslo \( \frac1{\left(2-\sqrt3\right)^3} \) se rovná:}
{\( 26+15\sqrt3 \)}{\( 27+24\sqrt3 \)}{\( 14+7\sqrt3 \)}{\( 27+30\sqrt3 \)}{}{}}
\LANGsk{
\Question{
Číslo \( \frac1{\left(2-\sqrt3\right)^3} \) sa rovná:}
{\( 26+15\sqrt3 \)}{\( 27+24\sqrt3 \)}{\( 14+7\sqrt3 \)}{\( 27+30\sqrt3 \)}{}{}}
\LANGes{
\Question{
El número \( \frac1{\left(2-\sqrt3\right)^3} \) es igual a:}
{\( 26+15\sqrt3 \)}{\( 27+24\sqrt3 \)}{\( 14+7\sqrt3 \)}{\( 27+30\sqrt3 \)}{}{}}
\End

\Data\ID{1003118108}
\Updated{2020-08-27 09:08:17}
\Level{C}
\SubArea{Expressions with powers and roots}
\LANGen{
\Question{
The value of \( \sqrt{5\left(222^2+111^2\right)} \) is:}
{\( 555 \)}{\( 111\sqrt5 \)}{\( \sqrt{15} \)}{\( \sqrt{1110^2-555^2} \)}{}{}}
\LANGpl{
\Question{
Wartość wyrażenia \( \sqrt{5\left(222^2+111^2\right)} \) jest równa:}
{\( 555 \)}{\( 111\sqrt5 \)}{\( \sqrt{15} \)}{\( \sqrt{1110^2-555^2} \)}{}{}}
\LANGcs{
\Question{
Hodnota \( \sqrt{5\left(222^2+111^2\right)} \) je:}
{\( 555 \)}{\( 111\sqrt5 \)}{\( \sqrt{15} \)}{\( \sqrt{1110^2-555^2} \)}{}{}}
\LANGsk{
\Question{
Hodnota \( \sqrt{5\left(222^2+111^2\right)} \) je:}
{\( 555 \)}{\( 111\sqrt5 \)}{\( \sqrt{15} \)}{\( \sqrt{1110^2-555^2} \)}{}{}}
\LANGes{
\Question{
El valor de \( \sqrt{5\left(222^2+111^2\right)} \) es:}
{\( 555 \)}{\( 111\sqrt5 \)}{\( \sqrt{15} \)}{\( \sqrt{1110^2-555^2} \)}{}{}}
\End

\Data\ID{1003118109}
\Updated{2020-10-12 10:10:04}
\Level{C}
\SubArea{Expressions with powers and roots}
\LANGen{
\Question{
Which of the following numbers has the same unit digit as the number \( 555^{555} \)?}
{\( 10^{555}+5 \)}{\( \frac{10^{555}}5 \)}{\( 5\cdot10^{555} \)}{\( 10^{555}+10^{555} \)}{}{}}
\LANGpl{
\Question{
Cyfra jedności liczby \( 555^{555} \), jest taka sama jak cyfra jedności liczby:}
{\( 10^{555}+5 \)}{\( \frac{10^{555}}5 \)}{\( 5\cdot10^{555} \)}{\( 10^{555}+10^{555} \)}{}{}}
\LANGcs{
\Question{
Které z následujících čísel má stejnou číslici na místě jednotek jako číslo \( 555^{555} \)?}
{\( 10^{555}+5 \)}{\( \frac{10^{555}}5 \)}{\( 5\cdot10^{555} \)}{\( 10^{555}+10^{555} \)}{}{}}
\LANGsk{
\Question{
Ktoré z následujúcich čísel má rovnakú číslicu na mieste jednotiek ako číslo \( 555^{555} \)?}
{\( 10^{555}+5 \)}{\( \frac{10^{555}}5 \)}{\( 5\cdot10^{555} \)}{\( 10^{555}+10^{555} \)}{}{}}
\LANGes{
\Question{
¿Cuál de los números siguientes tiene el mismo dígito en la posición de las unidades que \( 555^{555} \)?}
{\( 10^{555}+5 \)}{\( \frac{10^{555}}5 \)}{\( 5\cdot10^{555} \)}{\( 10^{555}+10^{555} \)}{}{}}
\End

\Data\ID{1003118110}
\Updated{2020-10-12 10:10:28}
\Level{C}
\SubArea{Expressions with powers and roots}
\LANGen{
\Question{
Let \( \frac{u\sqrt5}{5-\sqrt5}=\frac{\frac{5+\sqrt5}{\sqrt5+1}}{1-\sqrt5} \). What is the reciprocal of \( u \)?}
{\( \frac1u=-\frac{\sqrt5}5 \)}{\( \frac1u = \frac1{\sqrt5} \)}{\( \frac1u = \frac5{\sqrt5} \)}{\( \frac1u = -\sqrt5 \)}{}{}}
\LANGpl{
\Question{
Niech  \( \frac{u\sqrt5}{5-\sqrt5}=\frac{\frac{5+\sqrt5}{\sqrt5+1}}{1-\sqrt5} \). Jaka jest odwrotność liczby   \( u \)?}
{\( \frac1u=-\frac{\sqrt5}5 \)}{\( \frac1u = \frac1{\sqrt5} \)}{\( \frac1u = \frac5{\sqrt5} \)}{\( \frac1u = -\sqrt5 \)}{}{}}
\LANGcs{
\Question{
Nechť \( \frac{u\sqrt5}{5-\sqrt5}=\frac{\frac{5+\sqrt5}{\sqrt5+1}}{1-\sqrt5} \). Jaké je převrácené číslo k \( u \)?}
{\( \frac1u=-\frac{\sqrt5}5 \)}{\( \frac1u = \frac1{\sqrt5} \)}{\( \frac1u = \frac5{\sqrt5} \)}{\( \frac1u = -\sqrt5 \)}{}{}}
\LANGsk{
\Question{
Nech platí \( \frac{u\sqrt5}{5-\sqrt5}=\frac{\frac{5+\sqrt5}{\sqrt5+1}}{1-\sqrt5} \). Aké je prevrátené číslo k \( u \)?}
{\( \frac1u=-\frac{\sqrt5}5 \)}{\( \frac1u = \frac1{\sqrt5} \)}{\( \frac1u = \frac5{\sqrt5} \)}{\( \frac1u = -\sqrt5 \)}{}{}}
\LANGes{
\Question{
Sea \( \frac{u\sqrt5}{5-\sqrt5}=\frac{\frac{5+\sqrt5}{\sqrt5+1}}{1-\sqrt5} \). ¿Cuá es el recíproco de  \( u \)?}
{\( \frac1u=-\frac{\sqrt5}5 \)}{\( \frac1u = \frac1{\sqrt5} \)}{\( \frac1u = \frac5{\sqrt5} \)}{\( \frac1u = -\sqrt5 \)}{}{}}
\End

\Data\ID{1003118201}
\Updated{2020-08-27 10:08:51}
\Level{C}
\SubArea{Expressions with powers and roots}
\LANGen{
\Question{
Suppose that \( \sqrt[3]2\approx 1.25 \). Give the approximate value of \( \left(\frac1{81}\right)^{-\sqrt[3]2}\). Do not use a calculator.}
{\( 243 \)}{\( 729 \)}{\( \frac19 \)}{\( 0 \)}{}{}}
\LANGpl{
\Question{
Wiadomo, że \( \sqrt[3]2\approx 1{,}25 \). Podaj przybliżoną wartość potęgi \( \left(\frac1{81}\right)^{-\sqrt[3]2}\). Nie używaj kalkulatora.}
{\( 243 \)}{\( 729 \)}{\( \frac19 \)}{\( 0 \)}{}{}}
\LANGcs{
\Question{
Předpokládejme, že platí \( \sqrt[3]2\approx 1{,}25 \). Určete přibližnou hodnotu čísla \( \left(\frac1{81}\right)^{-\sqrt[3]2}\) bez použití kalkulačky.}
{\( 243 \)}{\( 729 \)}{\( \frac19 \)}{\( 0 \)}{}{}}
\LANGsk{
\Question{
Predpokladajme, že platí \( \sqrt[3]2\approx 1{,}25 \). Určte približnú hodnotu čísla \( \left(\frac1{81}\right)^{-\sqrt[3]2}\) bez použitia kalkulačky.}
{\( 243 \)}{\( 729 \)}{\( \frac19 \)}{\( 0 \)}{}{}}
\LANGes{
\Question{
Supongamos \( \sqrt[3]2\approx 1.25 \). Halla el valor aproximado de \( \left(\frac1{81}\right)^{-\sqrt[3]2}\). No uses la calculadora.}
{\( 243 \)}{\( 729 \)}{\( \frac19 \)}{\( 0 \)}{}{}}
\End

\Data\ID{1003118202}
\Updated{2020-10-12 10:10:12}
\Level{C}
\SubArea{Expressions with powers and roots}
\LANGen{
\Question{
The approximate value of the number \( 10^{-0.8} \) is \( 0.158489 \). What is the approximate value of the number \( 10^{1.2} \)? Give your answer to three decimal places. Do not use a calculator.}
{\( 15.849 \)}{\( 1.585 \)}{\( 0.015 \)}{\( 3.170 \)}{}{}}
\LANGpl{
\Question{
Przybliżenie liczby  \( 10^{-0{,}8} \) jest równe \( 0{,}158489 \). Podaj przybliżenie dziesiętne liczby  \( 10^{1{,}2} \) z dokładnością do trzeciego miejsca po przecinku. Oblicz bez użycia kalkulatora.}
{\( 15{,}849 \)}{\( 1{,}585 \)}{\( 0{,}015 \)}{\( 3{,}170 \)}{}{}}
\LANGcs{
\Question{
Přibližná hodnota čísla \( 10^{-0{,}8} \) je \( 0{,}158489 \). Určete přibližnou hodnotu čísla \( 10^{1{,}2} \) bez použití kalkulačky. Zaokrouhlete výsledek na tři desetinná místa.}
{\( 15{,}849 \)}{\( 1{,}585 \)}{\( 0{,}015 \)}{\( 3{,}170 \)}{}{}}
\LANGsk{
\Question{
Približná hodnota čísla \( 10^{-0{,}8} \) je \( 0{,}158489 \). Určte približnú hodnotu čísla \( 10^{1{,}2} \) bez použitia kalkulačky. Zaokrúhlite výsledok na tri desatinné miesta.}
{\( 15{,}849 \)}{\( 1{,}585 \)}{\( 0{,}015 \)}{\( 3{,}170 \)}{}{}}
\LANGes{
\Question{
El valor aproximado del número \( 10^{-0.8} \) es \( 0.158489 \). ¿Cuál es el valor aproximado de \( 10^{1.2} \)? Escríbelo con tres dígitos decimales. No uses la calculadora.}
{\( 15.849 \)}{\( 1.585 \)}{\( 0.015 \)}{\( 3.170 \)}{}{}}
\End

\Data\ID{1003118203}
\Updated{2021-10-11 05:10:08}
\Level{C}
\SubArea{Expressions with powers and roots}
\LANGen{
\Question{
Assuming \( 12 < \sqrt{153} < 13 \), which interval does the number \( \frac{5-\sqrt{153}}5 \) belong to? Do not use a calculator.}
{\( (-1.6; -1.4 ) \)}{\( (-1.8; -1.5 ) \)}{\( (1.4; 1.6 ) \)}{\( (1.6; 1.8 ) \)}{}{}}
\LANGpl{
\Question{
Wiadomo, że \( 12 < \sqrt{153} < 13 \). Wskaż przedział do którego należy  liczba \( \frac{5-\sqrt{153}}5 \). Oblicz bez użycia kalkulatora.}
{\( (-1{,}6; -1{,}4 ) \)}{\( (-1{,}8; -1{,}5 ) \)}{\( (1{,}4; 1{,}6 ) \)}{\( (1{,}6; 1{,}8 ) \)}{}{}}
\LANGcs{
\Question{
Předpokládejme, že platí \( 12 < \sqrt{153} < 13 \). Bez použití kalkulačky určete interval, do kterého patří číslo \( \frac{5-\sqrt{153}}5 \).}
{\( (-1{,}6; -1{,}4 ) \)}{\( (-1{,}8; -1{,}5 ) \)}{\( (1{,}4; 1{,}6 ) \)}{\( (1{,}6; 1{,}8 ) \)}{}{}}
\LANGsk{
\Question{
Predpokladajme, že platí \( 12 < \sqrt{153} < 13 \). Bez použitia kalkulačky určte interval, do ktorého patrí číslo \( \frac{5-\sqrt{153}}5 \).}
{\( (-1{,}6; -1{,}4 ) \)}{\( (-1{,}8; -1{,}5 ) \)}{\( (1{,}4; 1{,}6 ) \)}{\( (1{,}6; 1{,}8 ) \)}{}{}}
\LANGes{
\Question{
Suponiendo que \( 12 < \sqrt{153} < 13 \), ¿a qué intervalo pertenece el número  \( \frac{5-\sqrt{153}}5 \) ? No uses la calculadora.}
{\( (-1.6; -1.4 ) \)}{\( (-1.8; -1.5 ) \)}{\( (1.4; 1.6 ) \)}{\( (1.6; 1.8 ) \)}{}{}}
\End

\Data\ID{1003118204}
\Updated{2020-10-12 10:10:03}
\Level{C}
\SubArea{Expressions with powers and roots}
\LANGen{
\Question{
The approximate value of the number \( 1.3\cdot10^{-0.4} \) is \( 0.5175393 \). What is the approximate value of the number \( 39\cdot10^{0.6} \)? Give your answer to three decimal places. Do not use a calculator.}
{\( 155.262 \)}{\( 1552.618 \)}{\( 15.526 \)}{\( 1552.617 \)}{}{}}
\LANGpl{
\Question{
Przybliżenie liczby \( 1{,}3\cdot10^{-0{,}4} \) jest równe \( 0{,}5175393 \). Podaj przybliżenie liczby \( 39\cdot10^{0{,}6} \) z dokładnością do trzeciego miejsca po przecinku. Nie używaj kalkulatora.}
{\( 155{,}262 \)}{\( 1552{,}618 \)}{\( 15{,}526 \)}{\( 1552{,}617 \)}{}{}}
\LANGcs{
\Question{
Přibližná hodnota čísla \( 1{,}3\cdot10^{-0{,}4} \) je \( 0{,}5175393 \). Určete přibližnou hodnotu čísla \( 39\cdot10^{0{,}6} \) bez použití kalkulačky. Zaokrouhlete výsledek na tři desetinná místa.}
{\( 155{,}262 \)}{\( 1552{,}618 \)}{\( 15{,}526 \)}{\( 1552{,}617 \)}{}{}}
\LANGsk{
\Question{
Približná hodnota čísla \( 1{,}3\cdot10^{-0{,}4} \) je \( 0{,}5175393 \). Určte približnú hodnotu čísla \( 39\cdot10^{0{,}6} \) bez použitia kalkulačky. Zaokrúhlite výsledok na tri desatinné miesta.}
{\( 155{,}262 \)}{\( 1552{,}618 \)}{\( 15{,}526 \)}{\( 1552{,}617 \)}{}{}}
\LANGes{
\Question{
El valor aproximado del número \( 1.3\cdot10^{-0.4} \) es \( 0.5175393 \). ¿Cuál es el valor aproximado de \( 39\cdot10^{0.6} \)? Halla la respuesta con tres cifras decimales. No uses la calculadora.}
{\( 155.262 \)}{\( 1552.618 \)}{\( 15.526 \)}{\( 1552.617 \)}{}{}}
\End

\Data\ID{1003118205}
\Updated{2021-10-11 05:10:23}
\Level{C}
\SubArea{Expressions with powers and roots}
\LANGen{
\Question{
The approximate value of the number \( 10^{0.1} \) to five decimal places is \( 1.25893 \). Give the approximate value of the number \( 10^{-0.9} \) to three decimal places. Do not use a calculator.}
{\( 0.126 \)}{\( 1.259 \)}{\( 12.589 \)}{\( 7.943 \)}{}{}}
\LANGpl{
\Question{
Przybliżenie dziesiętne liczby \( 10^{0{,}1} \) z dokładnością do pięciu miejsc po przecinku jest równe \( 1{,}25893 \). Podaj przybliżenie dziesiętne liczby \( 10^{-0{,}9} \) z dokładnością do trzech  miejsc po przecinku. Oblicz bez użycia kalkulatora.}
{\( 0{,}126 \)}{\( 1{,}259 \)}{\( 12{,}589 \)}{\( 7{,}943 \)}{}{}}
\LANGcs{
\Question{
Přibližná hodnota čísla \( 10^{0{,}1} \) na pět desetinných míst je \( 1{,}25893 \). Určete přibližnou hodnotu čísla \( 10^{-0{,}9} \) bez použití kalkulačky. Výsledek zaokrouhlete na tři desetinná místa.}
{\( 0{,}126 \)}{\( 1{,}259 \)}{\( 12{,}589 \)}{\( 7{,}943 \)}{}{}}
\LANGsk{
\Question{
Približná hodnota čísla \( 10^{0{,}1} \) na päť desatinných miest je \( 1{,}25893 \). Určte približnú hodnotu čísla \( 10^{-0{,}9} \) bez použitia kalkulačky. Výsledok zaokrúhlite na tri desatinné miesta.}
{\( 0{,}126 \)}{\( 1{,}259 \)}{\( 12{,}589 \)}{\( 7{,}943 \)}{}{}}
\LANGes{
\Question{
El valor aproximado con cinco cifras decimales del número \( 10^{0.1} \) es\( 1.25893 \). Halla el valor aproximado con tres cifras decimales de \( 10^{-0.9} \). No uses la calculadora.}
{\( 0.126 \)}{\( 1.259 \)}{\( 12.589 \)}{\( 7.943 \)}{}{}}
\End

\Data\ID{1003124902}
\Updated{2020-08-27 02:08:47}
\Level{C}
\SubArea{Expressions with powers and roots}
\LANGen{
\Question{
Rationalizing the denominator \( \frac1{\left(2-\sqrt3\right)^3} \) you get:}
{\( 26+15\sqrt3 \)}{\( 27+30\sqrt3 \)}{\( 14+7\sqrt3 \)}{\( 27+24\sqrt3 \)}{}{}}
\LANGpl{
\Question{
Liczba \( \frac1{\left(2-\sqrt3\right)^3} \) jest równa:}
{\( 26+15\sqrt3 \)}{\( 27+30\sqrt3 \)}{\( 14+7\sqrt3 \)}{\( 27+24\sqrt3 \)}{}{}}
\LANGcs{
\Question{
Zjednodušením zlomku \( \frac1{\left(2-\sqrt3\right)^3} \) dostaneme:}
{\( 26+15\sqrt3 \)}{\( 27+30\sqrt3 \)}{\( 14+7\sqrt3 \)}{\( 27+24\sqrt3 \)}{}{}}
\LANGsk{
\Question{
Zjednodušením zlomku \( \frac1{\left(2-\sqrt3\right)^3} \) dostaneme:}
{\( 26+15\sqrt3 \)}{\( 27+30\sqrt3 \)}{\( 14+7\sqrt3 \)}{\( 27+24\sqrt3 \)}{}{}}
\LANGes{
\Question{
Racionalizando el denominador de la fracción \( \frac1{\left(2-\sqrt3\right)^3} \) obtienes:}
{\( 26+15\sqrt3 \)}{\( 27+30\sqrt3 \)}{\( 14+7\sqrt3 \)}{\( 27+24\sqrt3 \)}{}{}}
\End

\Data\ID{1003124903}
\Updated{2020-10-12 10:10:21}
\Level{C}
\SubArea{Expressions with powers and roots}
\LANGen{
\Question{
The number \( \left( \sqrt2+1 \right)^4 \) is greater than the number \( \left( \sqrt2-1 \right)^4 \) by:}
{\( 24\sqrt2 \)}{\( 12\sqrt2 \)}{\( 2 \)}{\( 4\sqrt2 \)}{}{}}
\LANGpl{
\Question{
Liczba \( \left( \sqrt2+1 \right)^4 \) jest większa od liczby \( \left( \sqrt2-1 \right)^4 \) o:}
{\( 24\sqrt2 \)}{\( 12\sqrt2 \)}{\( 2 \)}{\( 4\sqrt2 \)}{}{}}
\LANGcs{
\Question{
O kolik je číslo \( \left( \sqrt2+1 \right)^4 \) větší než číslo \( \left( \sqrt2-1 \right)^4 \)?}
{\( 24\sqrt2 \)}{\( 12\sqrt2 \)}{\( 2 \)}{\( 4\sqrt2 \)}{}{}}
\LANGsk{
\Question{
O koľko je číslo \( \left( \sqrt2+1 \right)^4 \) väčšie ako číslo \( \left( \sqrt2-1 \right)^4 \)?}
{\( 24\sqrt2 \)}{\( 12\sqrt2 \)}{\( 2 \)}{\( 4\sqrt2 \)}{}{}}
\LANGes{
\Question{
¿Cuántas veces es el número \( \left( \sqrt2+1 \right)^4 \) mayor que \( \left( \sqrt2-1 \right)^4 \) ?}
{\( 24\sqrt2 \)}{\( 12\sqrt2 \)}{\( 2 \)}{\( 4\sqrt2 \)}{}{}}
\End

\Data\ID{1003124904}
\Updated{2020-08-27 02:08:43}
\Level{C}
\SubArea{Expressions with powers and roots}
\LANGen{
\Question{
The number\( \sqrt[4]{\left( \sqrt3-\sqrt2 \right)^4}+\sqrt[4]{\left( \sqrt2-\sqrt5 \right)^4}+\sqrt[3]{\left( \sqrt3-\sqrt5 \right)^3} \)  equals:}
{\( 2\sqrt3 - 2\sqrt2 \)}{\( 2\sqrt5 - 2\sqrt2 \)}{\( 2\sqrt3 - 2\sqrt5 \)}{\( 2\sqrt5 - 2\sqrt3 \)}{}{}}
\LANGpl{
\Question{
Liczba \( \sqrt[4]{\left( \sqrt3-\sqrt2 \right)^4}+\sqrt[4]{\left( \sqrt2-\sqrt5 \right)^4}+\sqrt[3]{\left( \sqrt3-\sqrt5 \right)^3} \)  jest równa:}
{\( 2\sqrt3 - 2\sqrt2 \)}{\( 2\sqrt5 - 2\sqrt2 \)}{\( 2\sqrt3 - 2\sqrt5 \)}{\( 2\sqrt5 - 2\sqrt3 \)}{}{}}
\LANGcs{
\Question{
Číslo \( \sqrt[4]{\left( \sqrt3-\sqrt2 \right)^4}+\sqrt[4]{\left( \sqrt2-\sqrt5 \right)^4}+\sqrt[3]{\left( \sqrt3-\sqrt5 \right)^3} \)  je rovno:}
{\( 2\sqrt3 - 2\sqrt2 \)}{\( 2\sqrt5 - 2\sqrt2 \)}{\( 2\sqrt3 - 2\sqrt5 \)}{\( 2\sqrt5 - 2\sqrt3 \)}{}{}}
\LANGsk{
\Question{
Číslo \( \sqrt[4]{\left( \sqrt3-\sqrt2 \right)^4}+\sqrt[4]{\left( \sqrt2-\sqrt5 \right)^4}+\sqrt[3]{\left( \sqrt3-\sqrt5 \right)^3} \) sa rovná:}
{\( 2\sqrt3 - 2\sqrt2 \)}{\( 2\sqrt5 - 2\sqrt2 \)}{\( 2\sqrt3 - 2\sqrt5 \)}{\( 2\sqrt5 - 2\sqrt3 \)}{}{}}
\LANGes{
\Question{
El número \( \sqrt[4]{\left( \sqrt3-\sqrt2 \right)^4}+\sqrt[4]{\left( \sqrt2-\sqrt5 \right)^4}+\sqrt[3]{\left( \sqrt3-\sqrt5 \right)^3} \)  es igual a:}
{\( 2\sqrt3 - 2\sqrt2 \)}{\( 2\sqrt5 - 2\sqrt2 \)}{\( 2\sqrt3 - 2\sqrt5 \)}{\( 2\sqrt5 - 2\sqrt3 \)}{}{}}
\End

\Data\ID{1003124905}
\Updated{2020-10-12 10:10:56}
\Level{C}
\SubArea{Expressions with powers and roots}
\LANGen{
\Question{
Numbers  \( \left(2^2\right)^{\left(2^2\right)} \), \( \left(2\right)^{\left({2^2}^2\right)} \), \( \left({2^2}^2\right)^2 \), \( \left(2\right)^{\left(2^2\right)^2} \) were written on the board. How many different numbers are represented by these records?}
{\( 2 \)}{\( 3 \)}{\( 4 \)}{\( 1 \)}{}{}}
\LANGpl{
\Question{
Na tablicy zapisano liczby \( \left(2^2\right)^{\left(2^2\right)} \), \( \left(2\right)^{\left({2^2}^2\right)} \), \( \left({2^2}^2\right)^2 \), \( \left(2\right)^{\left(2^2\right)^2} \). Ile różnych liczb reprezentują te zapisy?}
{\( 2 \)}{\( 3 \)}{\( 4 \)}{\( 1 \)}{}{}}
\LANGcs{
\Question{
Čísla \( \left(2^2\right)^{\left(2^2\right)} \), \( \left(2\right)^{\left({2^2}^2\right)} \), \( \left({2^2}^2\right)^2 \), \( \left(2\right)^{\left(2^2\right)^2} \) byla napsaná na tabuli.  Kolik různých čísel představuje tento seznam?}
{\( 2 \)}{\( 3 \)}{\( 4 \)}{\( 1 \)}{}{}}
\LANGsk{
\Question{
Čísla \( \left(2^2\right)^{\left(2^2\right)} \), \( \left(2\right)^{\left({2^2}^2\right)} \), \( \left({2^2}^2\right)^2 \), \( \left(2\right)^{\left(2^2\right)^2} \) boli napísané na tabuľu.  Koľko rôznych čísel reprezentuje tento zoznam?}
{\( 2 \)}{\( 3 \)}{\( 4 \)}{\( 1 \)}{}{}}
\LANGes{
\Question{
Los números  \( \left(2^2\right)^{\left(2^2\right)} \), \( \left(2\right)^{\left({2^2}^2\right)} \), \( \left({2^2}^2\right)^2 \), \( \left(2\right)^{\left(2^2\right)^2} \) fueron escriton en la pizarra. ¿Cuántos números diferentes estuvieron representados por las expresiones?}
{\( 2 \)}{\( 3 \)}{\( 4 \)}{\( 1 \)}{}{}}
\End

\Data\ID{1003124906}
\Updated{2020-08-28 07:08:49}
\Level{C}
\SubArea{Expressions with powers and roots}
\LANGen{
\Question{
The number \( \left(\sqrt7+1\right)^4-\left(\sqrt7-1\right)^4 \) equals:}
{\( 64\sqrt7 \)}{\( 32\sqrt7 \)}{\( \sqrt7 \)}{\( 2 \)}{}{}}
\LANGpl{
\Question{
Liczba \( \left(\sqrt7+1\right)^4-\left(\sqrt7-1\right)^4 \) jest równa:}
{\( 64\sqrt7 \)}{\( 32\sqrt7 \)}{\( \sqrt7 \)}{\( 2 \)}{}{}}
\LANGcs{
\Question{
Číslo \( \left(\sqrt7+1\right)^4-\left(\sqrt7-1\right)^4 \) je rovno:}
{\( 64\sqrt7 \)}{\( 32\sqrt7 \)}{\( \sqrt7 \)}{\( 2 \)}{}{}}
\LANGsk{
\Question{
Číslo \( \left(\sqrt7+1\right)^4-\left(\sqrt7-1\right)^4 \) sa rovná:}
{\( 64\sqrt7 \)}{\( 32\sqrt7 \)}{\( \sqrt7 \)}{\( 2 \)}{}{}}
\LANGes{
\Question{
El número \( \left(\sqrt7+1\right)^4-\left(\sqrt7-1\right)^4 \) es igual a:}
{\( 64\sqrt7 \)}{\( 32\sqrt7 \)}{\( \sqrt7 \)}{\( 2 \)}{}{}}
\End

\Data\ID{1003134508}
\Updated{2020-08-27 02:08:43}
\Level{C}
\SubArea{Expressions with powers and roots}
\LANGen{
\Question{
The number \( \left(\sqrt2-1\right)^6 \) is equal to:}
{\( 99-70\sqrt2 \)}{\( 5\sqrt2 - 7 \)}{\( 49-35\sqrt2 \)}{\( 34-24\sqrt2 \)}{}{}}
\LANGpl{
\Question{
Liczba \( \left(\sqrt2-1\right)^6 \) jest równa:}
{\( 99-70\sqrt2 \)}{\( 5\sqrt2 - 7 \)}{\( 49-35\sqrt2 \)}{\( 34-24\sqrt2 \)}{}{}}
\LANGcs{
\Question{
Číslo \( \left(\sqrt2-1\right)^6 \) je rovno:}
{\( 99-70\sqrt2 \)}{\( 5\sqrt2 - 7 \)}{\( 49-35\sqrt2 \)}{\( 34-24\sqrt2 \)}{}{}}
\LANGsk{
\Question{
Číslo \( \left(\sqrt2-1\right)^6 \) sa rovná:}
{\( 99-70\sqrt2 \)}{\( 5\sqrt2 - 7 \)}{\( 49-35\sqrt2 \)}{\( 34-24\sqrt2 \)}{}{}}
\LANGes{
\Question{
El número \( \left(\sqrt2-1\right)^6 \) es igual a:}
{\( 99-70\sqrt2 \)}{\( 5\sqrt2 - 7 \)}{\( 49-35\sqrt2 \)}{\( 34-24\sqrt2 \)}{}{}}
\End

\Data\ID{1003134510}
\Updated{2020-08-28 08:08:33}
\Level{C}
\SubArea{Expressions with powers and roots}
\LANGen{
\Question{
Complete the following sentence to obtain a true statement. The number \( \left( \sqrt3 \right)^{\sqrt2^{\sqrt[3]{64}}} \) is ...}
{a rational number less than \( 10 \).}{an irrational number less than \( 10 \).}{an integer greater than \( 10 \).}{a rational number greater than \( 10 \).}{}{}}
\LANGpl{
\Question{
Liczba \( \left( \sqrt3 \right)^{\sqrt2^{\sqrt[3]{64}}} \) jest liczbą}
{wymierną mniejszą od  \( 10 \).}{niewymierną mniejszą od  \( 10 \).}{całkowitą większą od  \( 10 \).}{wymierną większą od  \( 10 \).}{}{}}
\LANGcs{
\Question{
Doplňte následující tvrzení, aby byl výrok pravdivý. Číslo \( \left( \sqrt3 \right)^{\sqrt2^{\sqrt[3]{64}}} \) je ...}
{racionální číslo menší než \( 10 \).}{iracionální číslo menší než \( 10 \).}{celé číslo větší než \( 10 \).}{racionální číslo větší než \( 10 \).}{}{}}
\LANGsk{
\Question{
Doplňte nasledujúce tvrdenie, aby bol výrok pravdivý. Číslo \( \left( \sqrt3 \right)^{\sqrt2^{\sqrt[3]{64}}} \) je ...}
{racionálne číslo menšie ako \( 10 \).}{iracionálne číslo menšie ako \( 10 \).}{celé číslo väčšie ako \( 10 \).}{racionálne číslo väčšie ako \( 10 \).}{}{}}
\LANGes{
\Question{
Completa la frase siguiente para que sea correcta. El número \( \left( \sqrt3 \right)^{\sqrt2^{\sqrt[3]{64}}} \) es ...}
{un número racional menor que \( 10 \).}{un número irracional menor que \( 10 \).}{un número entero mayor que \( 10 \).}{un número racional mayor que \( 10 \).}{}{}}
\End

\Data\ID{1003158601}
\Updated{2020-08-28 08:08:01}
\Level{C}
\SubArea{Expressions with powers and roots}
\LANGen{
\Question{
The number \( \sqrt[5]9 \) is bigger than:}
{\( 3^{0.3} \)}{\( \sqrt3 \)}{\( \sqrt[9]{81} \)}{\( 3 \)}{}{}}
\LANGpl{
\Question{
Liczba  \( \sqrt[5]9 \) jest większa od:}
{\( 3^{0{,}3} \)}{\( \sqrt3 \)}{\( \sqrt[9]{81} \)}{\( 3 \)}{}{}}
\LANGcs{
\Question{
Číslo \( \sqrt[5]9 \) je větší než:}
{\( 3^{0{,}3} \)}{\( \sqrt3 \)}{\( \sqrt[9]{81} \)}{\( 3 \)}{}{}}
\LANGsk{
\Question{
Číslo \( \sqrt[5]9 \) je väčšie ako:}
{\( 3^{0{,}3} \)}{\( \sqrt3 \)}{\( \sqrt[9]{81} \)}{\( 3 \)}{}{}}
\LANGes{
\Question{
El número \( \sqrt[5]9 \) es mayor que:}
{\( 3^{0.3} \)}{\( \sqrt3 \)}{\( \sqrt[9]{81} \)}{\( 3 \)}{}{}}
\End

\Data\ID{1003158602}
\Updated{2020-10-12 10:10:59}
\Level{C}
\SubArea{Expressions with powers and roots}
\LANGen{
\Question{
If \( a=\frac1{\sqrt3-1} \) and \( b=\frac1{\sqrt5-1}-\frac2{\sqrt5}+1 \), then:}
{\( a > b \)}{\( a < b \)}{\( a = b \)}{\( a = -b \)}{}{}}
\LANGpl{
\Question{
Jeśli \( a=\frac1{\sqrt3-1} \) i \( b=\frac1{\sqrt5-1}-\frac2{\sqrt5}+1 \), to:}
{\( a > b \)}{\( a < b \)}{\( a = b \)}{\( a = -b \)}{}{}}
\LANGcs{
\Question{
Jestliže \( a=\frac1{\sqrt3-1} \) a \( b=\frac1{\sqrt5-1}-\frac2{\sqrt5}+1 \), pak:}
{\( a > b \)}{\( a < b \)}{\( a = b \)}{\( a = -b \)}{}{}}
\LANGsk{
\Question{
Ak \( a=\frac1{\sqrt3-1} \) a \( b=\frac1{\sqrt5-1}-\frac2{\sqrt5}+1 \), potom:}
{\( a > b \)}{\( a < b \)}{\( a = b \)}{\( a = -b \)}{}{}}
\LANGes{
\Question{
Si \( a=\frac1{\sqrt3-1} \) y \( b=\frac1{\sqrt5-1}-\frac2{\sqrt5}+1 \), entonces:}
{\( a > b \)}{\( a < b \)}{\( a = b \)}{\( a = -b \)}{}{}}
\End

\Data\ID{1003158603}
\Updated{2020-08-28 08:08:59}
\Level{C}
\SubArea{Expressions with powers and roots}
\LANGen{
\Question{
After rationalizing the denominator \( \frac{\sqrt{\sqrt{13}+\sqrt{12}}}{\sqrt{\sqrt{13}-\sqrt{12}}} \) we get:}
{\( \sqrt{13}+2\sqrt3 \)}{\( \sqrt{13}-\sqrt{12} \)}{\( 1 \)}{\( 25 \)}{}{}}
\LANGpl{
\Question{
Po usunięciu niewymierności z mianownika  \( \frac{\sqrt{\sqrt{13}+\sqrt{12}}}{\sqrt{\sqrt{13}-\sqrt{12}}} \) otrzymamy:}
{\( \sqrt{13}+2\sqrt3 \)}{\( \sqrt{13}-\sqrt{12} \)}{\( 1 \)}{\( 25 \)}{}{}}
\LANGcs{
\Question{
Usměrněním zlomku \( \frac{\sqrt{\sqrt{13}+\sqrt{12}}}{\sqrt{\sqrt{13}-\sqrt{12}}} \) dostaneme:}
{\( \sqrt{13}+2\sqrt3 \)}{\( \sqrt{13}-\sqrt{12} \)}{\( 1 \)}{\( 25 \)}{}{}}
\LANGsk{
\Question{
Usmernením zlomku \( \frac{\sqrt{\sqrt{13}+\sqrt{12}}}{\sqrt{\sqrt{13}-\sqrt{12}}} \) dostaneme:}
{\( \sqrt{13}+2\sqrt3 \)}{\( \sqrt{13}-\sqrt{12} \)}{\( 1 \)}{\( 25 \)}{}{}}
\LANGes{
\Question{
Racionalizando el denominador de \( \frac{\sqrt{\sqrt{13}+\sqrt{12}}}{\sqrt{\sqrt{13}-\sqrt{12}}} \) obtenemos:}
{\( \sqrt{13}+2\sqrt3 \)}{\( \sqrt{13}-\sqrt{12} \)}{\( 1 \)}{\( 25 \)}{}{}}
\End

\Data\ID{1003158604}
\Updated{2020-08-28 08:08:35}
\Level{C}
\SubArea{Expressions with powers and roots}
\LANGen{
\Question{
After rationalizing the denominator \( \frac4{1-\sqrt2+\sqrt3} \) we get:}
{\( \sqrt6-\sqrt2+2 \)}{\( -4 \)}{\( \sqrt6-\sqrt2 \)}{\( 1+\sqrt2-\sqrt3 \)}{}{}}
\LANGpl{
\Question{
Po usunięciu niewymierności z mianownika  \( \frac4{1-\sqrt2+\sqrt3} \) otrzymamy:}
{\( \sqrt6-\sqrt2+2 \)}{\( -4 \)}{\( \sqrt6-\sqrt2 \)}{\( 1+\sqrt2-\sqrt3 \)}{}{}}
\LANGcs{
\Question{
Usměrněním zlomku \( \frac4{1-\sqrt2+\sqrt3} \) dostaneme:}
{\( \sqrt6-\sqrt2+2 \)}{\( -4 \)}{\( \sqrt6-\sqrt2 \)}{\( 1+\sqrt2-\sqrt3 \)}{}{}}
\LANGsk{
\Question{
Usmernením zlomku \( \frac4{1-\sqrt2+\sqrt3} \) dostaneme:}
{\( \sqrt6-\sqrt2+2 \)}{\( -4 \)}{\( \sqrt6-\sqrt2 \)}{\( 1+\sqrt2-\sqrt3 \)}{}{}}
\LANGes{
\Question{
Racionalizando el denominador de la fracción \( \frac4{1-\sqrt2+\sqrt3} \) obtenemos:}
{\( \sqrt6-\sqrt2+2 \)}{\( -4 \)}{\( \sqrt6-\sqrt2 \)}{\( 1+\sqrt2-\sqrt3 \)}{}{}}
\End

\Data\ID{1003158605}
\Updated{2020-08-28 08:08:32}
\Level{C}
\SubArea{Expressions with powers and roots}
\LANGen{
\Question{
Using the formula for factoring \( a^3-b^3 \) to rationalize the denominator of \( \frac2{\sqrt[3]{25}+\sqrt[3]{15}+\sqrt[3]{9}} \) you get:}
{\( \sqrt[3]5-\sqrt[3]3 \)}{\( 2 \)}{\( 2\sqrt[3]{49} \)}{\( 2\left( \sqrt[3]{25}-\sqrt[3]{15-\sqrt[3]9}\right)  \)}{}{}}
\LANGpl{
\Question{
Używając wzoru  \( a^3-b^3 \) do usunięcia niewymierności z mianownika \( \frac2{\sqrt[3]{25}+\sqrt[3]{15}+\sqrt[3]{9}} \) otrzymamy:}
{\( \sqrt[3]5-\sqrt[3]3 \)}{\( 2 \)}{\( 2\sqrt[3]{49} \)}{\( 2\left( \sqrt[3]{25}-\sqrt[3]{15-\sqrt[3]9}\right)  \)}{}{}}
\LANGcs{
\Question{
Usměrněním zlomku \( \frac2{\sqrt[3]{25}+\sqrt[3]{15}+\sqrt[3]{9}} \) s využitím vzorce \( a^3-b^3 \) dostaneme:}
{\( \sqrt[3]5-\sqrt[3]3 \)}{\( 2 \)}{\( 2\sqrt[3]{49} \)}{\( 2\left( \sqrt[3]{25}-\sqrt[3]{15-\sqrt[3]9}\right)  \)}{}{}}
\LANGsk{
\Question{
Usmernením zlomku \( \frac2{\sqrt[3]{25}+\sqrt[3]{15}+\sqrt[3]{9}} \) s využitím vzorca \( a^3-b^3 \) dostaneme:}
{\( \sqrt[3]5-\sqrt[3]3 \)}{\( 2 \)}{\( 2\sqrt[3]{49} \)}{\( 2\left( \sqrt[3]{25}-\sqrt[3]{15-\sqrt[3]9}\right)  \)}{}{}}
\LANGes{
\Question{
Usando la fórmula \( a^3-b^3 \) para racionalizar el denominador de \( \frac2{\sqrt[3]{25}+\sqrt[3]{15}+\sqrt[3]{9}} \) obtienes:}
{\( \sqrt[3]5-\sqrt[3]3 \)}{\( 2 \)}{\( 2\sqrt[3]{49} \)}{\( 2\left( \sqrt[3]{25}-\sqrt[3]{15-\sqrt[3]9}\right)  \)}{}{}}
\End

\Data\ID{1003158606}
\Updated{2020-08-28 12:08:41}
\Level{C}
\SubArea{Expressions with powers and roots}
\LANGen{
\Question{
Using the formula for factoring \( a^3-b^3 \) to find the multiplicative inverse of \( x=\sqrt[3]3-\sqrt[3]4 \) you get:}
{\( -\sqrt[3]9-\sqrt[3]{12}-\sqrt[3]{16} \)}{\( \sqrt[3]9-\sqrt[3]{12}+\sqrt[3]{16} \)}{\( \sqrt[3]3-\sqrt[3]4 \)}{\( \sqrt[3]4-\sqrt[3]3 \)}{}{}}
\LANGpl{
\Question{
Używając wzoru \( a^3-b^3 \) do utworzenia odwrotności liczby \( x=\sqrt[3]3-\sqrt[3]4 \) otrzymamy:}
{\( -\sqrt[3]9-\sqrt[3]{12}-\sqrt[3]{16} \)}{\( \sqrt[3]9-\sqrt[3]{12}+\sqrt[3]{16} \)}{\( \sqrt[3]3-\sqrt[3]4 \)}{\( \sqrt[3]4-\sqrt[3]3 \)}{}{}}
\LANGcs{
\Question{
Nelezením inverzního prvku k číslu \( x=\sqrt[3]3-\sqrt[3]4 \) využitím vzorce \( a^3-b^3 \) dostaneme:}
{\( -\sqrt[3]9-\sqrt[3]{12}-\sqrt[3]{16} \)}{\( \sqrt[3]9-\sqrt[3]{12}+\sqrt[3]{16} \)}{\( \sqrt[3]3-\sqrt[3]4 \)}{\( \sqrt[3]4-\sqrt[3]3 \)}{}{}}
\LANGsk{
\Question{
Nájdením inverzného prvku k číslu \( x=\sqrt[3]3-\sqrt[3]4 \) využitím vzorca \( a^3-b^3 \) dostaneme:}
{\( -\sqrt[3]9-\sqrt[3]{12}-\sqrt[3]{16} \)}{\( \sqrt[3]9-\sqrt[3]{12}+\sqrt[3]{16} \)}{\( \sqrt[3]3-\sqrt[3]4 \)}{\( \sqrt[3]4-\sqrt[3]3 \)}{}{}}
\LANGes{
\Question{
Usando la fórmula \( a^3-b^3 \) para encontrar el inverso multiplicativo de  \( x=\sqrt[3]3-\sqrt[3]4 \) obtenemos:}
{\( -\sqrt[3]9-\sqrt[3]{12}-\sqrt[3]{16} \)}{\( \sqrt[3]9-\sqrt[3]{12}+\sqrt[3]{16} \)}{\( \sqrt[3]3-\sqrt[3]4 \)}{\( \sqrt[3]4-\sqrt[3]3 \)}{}{}}
\End

\Data\ID{1003158607}
\Updated{2020-08-28 12:08:39}
\Level{C}
\SubArea{Expressions with powers and roots}
\LANGen{
\Question{
As a result of this mathematical action \( \sqrt[3]{\left[\left(\sqrt2+1\right)^3-\left(\sqrt2-1\right)^3\right]^2-13^2} \) we get:}
{\( 3 \)}{\( \sqrt[3]{-170} \)}{\( \sqrt[3]{13} \)}{\( \sqrt[3]{-168} \)}{}{}}
\LANGpl{
\Question{
W wyniku działania \( \sqrt[3]{\left[\left(\sqrt2+1\right)^3-\left(\sqrt2-1\right)^3\right]^2-13^2} \) otrzymamy:}
{\( 3 \)}{\( \sqrt[3]{-170} \)}{\( \sqrt[3]{13} \)}{\( \sqrt[3]{-168} \)}{}{}}
\LANGcs{
\Question{
Zjednodušením výrazu \( \sqrt[3]{\left[\left(\sqrt2+1\right)^3-\left(\sqrt2-1\right)^3\right]^2-13^2} \) dostaneme:}
{\( 3 \)}{\( \sqrt[3]{-170} \)}{\( \sqrt[3]{13} \)}{\( \sqrt[3]{-168} \)}{}{}}
\LANGsk{
\Question{
Zjednodušením výrazu \( \sqrt[3]{\left[\left(\sqrt2+1\right)^3-\left(\sqrt2-1\right)^3\right]^2-13^2} \) dostaneme:}
{\( 3 \)}{\( \sqrt[3]{-170} \)}{\( \sqrt[3]{13} \)}{\( \sqrt[3]{-168} \)}{}{}}
\LANGes{
\Question{
Simplificando la expresión \( \sqrt[3]{\left[\left(\sqrt2+1\right)^3-\left(\sqrt2-1\right)^3\right]^2-13^2} \) obtenemos:}
{\( 3 \)}{\( \sqrt[3]{-170} \)}{\( \sqrt[3]{13} \)}{\( \sqrt[3]{-168} \)}{}{}}
\End

\Data\ID{1003158608}
\Updated{2020-08-28 12:08:51}
\Level{C}
\SubArea{Expressions with powers and roots}
\LANGen{
\Question{
As a result of this mathematical action \( \frac1{\sqrt2+1}+\frac1{\sqrt3+\sqrt2}+\frac1{\sqrt4+\sqrt3}+\dots+\frac1{\sqrt{100}+\sqrt{99}} \) we get:}
{\( 9 \)}{\( 12 \)}{\( 100 \)}{\( 99 \)}{}{}}
\LANGpl{
\Question{
W wyniku działania  \( \frac1{\sqrt2+1}+\frac1{\sqrt3+\sqrt2}+\frac1{\sqrt4+\sqrt3}+\dots+\frac1{\sqrt{100}+\sqrt{99}} \) otrzymamy:}
{\( 9 \)}{\( 12 \)}{\( 100 \)}{\( 99 \)}{}{}}
\LANGcs{
\Question{
Úpravou výrazu \( \frac1{\sqrt2+1}+\frac1{\sqrt3+\sqrt2}+\frac1{\sqrt4+\sqrt3}+\dots+\frac1{\sqrt{100}+\sqrt{99}} \) dostaneme:}
{\( 9 \)}{\( 12 \)}{\( 100 \)}{\( 99 \)}{}{}}
\LANGsk{
\Question{
Úpravou výrazu \( \frac1{\sqrt2+1}+\frac1{\sqrt3+\sqrt2}+\frac1{\sqrt4+\sqrt3}+\dots+\frac1{\sqrt{100}+\sqrt{99}} \) dostaneme:}
{\( 9 \)}{\( 12 \)}{\( 100 \)}{\( 99 \)}{}{}}
\LANGes{
\Question{
Como resultado de esta expresión \( \frac1{\sqrt2+1}+\frac1{\sqrt3+\sqrt2}+\frac1{\sqrt4+\sqrt3}+\dots+\frac1{\sqrt{100}+\sqrt{99}} \) obtenemos:}
{\( 9 \)}{\( 12 \)}{\( 100 \)}{\( 99 \)}{}{}}
\End

\Data\ID{1003158609}
\Updated{2020-08-28 01:08:08}
\Level{C}
\SubArea{Expressions with powers and roots}
\LANGen{
\Question{
Choose the correct relation.}
{\( \frac{4^{2\sqrt3}\cdot8}{16^{\sqrt3+1}} =\frac12 \)}{\( \frac{9^{2\sqrt5}\cdot27^{\frac13}}{81^{1-\sqrt5}} =\frac1{27} \)}{\( \frac{\pi^{1-\sqrt2}\cdot\pi^{1+\sqrt2}}{\pi^2} < \frac1{\pi^3} \)}{\( \frac{2^{2\pi}\cdot9^{\pi}}{36^{\pi-1}} < 6 \)}{}{}}
\LANGpl{
\Question{
Wskaż zdanie prawdziwe.}
{\( \frac{4^{2\sqrt3}\cdot8}{16^{\sqrt3+1}} =\frac12 \)}{\( \frac{9^{2\sqrt5}\cdot27^{\frac13}}{81^{1-\sqrt5}} =\frac1{27} \)}{\( \frac{\pi^{1-\sqrt2}\cdot\pi^{1+\sqrt2}}{\pi^2} < \frac1{\pi^3} \)}{\( \frac{2^{2\pi}\cdot9^{\pi}}{36^{\pi-1}} < 6 \)}{}{}}
\LANGcs{
\Question{
Vyberte pravdivé tvrzení.}
{\( \frac{4^{2\sqrt3}\cdot8}{16^{\sqrt3+1}} =\frac12 \)}{\( \frac{9^{2\sqrt5}\cdot27^{\frac13}}{81^{1-\sqrt5}} =\frac1{27} \)}{\( \frac{\pi^{1-\sqrt2}\cdot\pi^{1+\sqrt2}}{\pi^2} < \frac1{\pi^3} \)}{\( \frac{2^{2\pi}\cdot9^{\pi}}{36^{\pi-1}} < 6 \)}{}{}}
\LANGsk{
\Question{
Vyberte pravdivé tvrdenie.}
{\( \frac{4^{2\sqrt3}\cdot8}{16^{\sqrt3+1}} =\frac12 \)}{\( \frac{9^{2\sqrt5}\cdot27^{\frac13}}{81^{1-\sqrt5}} =\frac1{27} \)}{\( \frac{\pi^{1-\sqrt2}\cdot\pi^{1+\sqrt2}}{\pi^2} < \frac1{\pi^3} \)}{\( \frac{2^{2\pi}\cdot9^{\pi}}{36^{\pi-1}} < 6 \)}{}{}}
\LANGes{
\Question{
Elije la relación correcta.}
{\( \frac{4^{2\sqrt3}\cdot8}{16^{\sqrt3+1}} =\frac12 \)}{\( \frac{9^{2\sqrt5}\cdot27^{\frac13}}{81^{1-\sqrt5}} =\frac1{27} \)}{\( \frac{\pi^{1-\sqrt2}\cdot\pi^{1+\sqrt2}}{\pi^2} < \frac1{\pi^3} \)}{\( \frac{2^{2\pi}\cdot9^{\pi}}{36^{\pi-1}} < 6 \)}{}{}}
\End

\Data\ID{1003164001}
\Updated{2020-08-28 01:08:21}
\Level{C}
\SubArea{Expressions with powers and roots}
\LANGen{
\Question{
After simplifying the expression 
\[ \sqrt[3]{5\sqrt2+7}-\sqrt[3]{5\sqrt2-7} \]
we get:}
{\( 2 \)}{\( 2\sqrt[3]{57} \)}{\( 1 \)}{\( 2\sqrt[3]{97} \)}{}{}}
\LANGpl{
\Question{
Po uproszczeniu wyrażenia   
\[ \sqrt[3]{5\sqrt2+7}-\sqrt[3]{5\sqrt2-7} \]
otrzymamy:}
{\( 2 \)}{\( 2\sqrt[3]{57} \)}{\( 1 \)}{\( 2\sqrt[3]{97} \)}{}{}}
\LANGcs{
\Question{
Zjednodušením výrazu 
\[ \sqrt[3]{5\sqrt2+7}-\sqrt[3]{5\sqrt2-7} \]
dostaneme:}
{\( 2 \)}{\( 2\sqrt[3]{57} \)}{\( 1 \)}{\( 2\sqrt[3]{97} \)}{}{}}
\LANGsk{
\Question{
Zjednodušením výrazu 
\[ \sqrt[3]{5\sqrt2+7}-\sqrt[3]{5\sqrt2-7} \]
dostaneme:}
{\( 2 \)}{\( 2\sqrt[3]{57} \)}{\( 1 \)}{\( 2\sqrt[3]{97} \)}{}{}}
\LANGes{
\Question{
Simplificando la expresión 
\[ \sqrt[3]{5\sqrt2+7}-\sqrt[3]{5\sqrt2-7} \]
obtenemos:}
{\( 2 \)}{\( 2\sqrt[3]{57} \)}{\( 1 \)}{\( 2\sqrt[3]{97} \)}{}{}}
\End

\Data\ID{1003164002}
\Updated{2020-08-28 01:08:01}
\Level{C}
\SubArea{Expressions with powers and roots}
\LANGen{
\Question{
After simplifying the expression
\[ a=\sqrt{6-2\sqrt5}-\sqrt5 \]
we get:}
{\( -1 \)}{\( 6-\sqrt5 \)}{\( 6+\sqrt5 \)}{\( 1-2\sqrt5 \)}{}{}}
\LANGpl{
\Question{
Liczba
\[ a=\sqrt{6-2\sqrt5}-\sqrt5 \]
po uproszczeniu jest równa:}
{\( -1 \)}{\( 6-\sqrt5 \)}{\( 6+\sqrt5 \)}{\( 1-2\sqrt5 \)}{}{}}
\LANGcs{
\Question{
Zjednodušením výrazu
\[ a=\sqrt{6-2\sqrt5}-\sqrt5 \]
dostaneme:}
{\( -1 \)}{\( 6-\sqrt5 \)}{\( 6+\sqrt5 \)}{\( 1-2\sqrt5 \)}{}{}}
\LANGsk{
\Question{
Zjednodušením výrazu
\[ a=\sqrt{6-2\sqrt5}-\sqrt5 \]
dostaneme:}
{\( -1 \)}{\( 6-\sqrt5 \)}{\( 6+\sqrt5 \)}{\( 1-2\sqrt5 \)}{}{}}
\LANGes{
\Question{
Simplificando la expresión
\[ a=\sqrt{6-2\sqrt5}-\sqrt5 \]
obtenemos:}
{\( -1 \)}{\( 6-\sqrt5 \)}{\( 6+\sqrt5 \)}{\( 1-2\sqrt5 \)}{}{}}
\End

\Data\ID{1003164003}
\Updated{2020-08-28 01:08:39}
\Level{C}
\SubArea{Expressions with powers and roots}
\LANGen{
\Question{
Let
\[ a=\left[\left(2-\sqrt3\right)^{\frac12}+\left(2+\sqrt3\right)^{\frac12}\right]^2,\ b=\frac{81^{-1}\cdot\sqrt3}{27^{-2}\cdot\sqrt[4]9}.\]
Comparing the numbers \( a^b \) and \( b^a \) you get:}
{\( a^b > b^a \)}{\( a^b < b^a \)}{\( a^b \leq b^a \)}{\( a^b = b^a \)}{}{}}
\LANGpl{
\Question{
Niech
\[ a=\left[\left(2-\sqrt3\right)^{\frac12}+\left(2+\sqrt3\right)^{\frac12}\right]^2,\ b=\frac{81^{-1}\cdot\sqrt3}{27^{-2}\cdot\sqrt[4]9}.\]
Porównując liczby  \( a^b \) i \( b^a \) otrzymamy:}
{\( a^b > b^a \)}{\( a^b < b^a \)}{\( a^b \leq b^a \)}{\( a^b = b^a \)}{}{}}
\LANGcs{
\Question{
Nechť
\[ a=\left[\left(2-\sqrt3\right)^{\frac12}+\left(2+\sqrt3\right)^{\frac12}\right]^2,\  b=\frac{81^{-1}\cdot\sqrt3}{27^{-2}\cdot\sqrt[4]9}.\]
Porovnáním čísel \( a^b \) a \( b^a \) dostaneme:}
{\( a^b > b^a \)}{\( a^b < b^a \)}{\( a^b \leq b^a \)}{\( a^b = b^a \)}{}{}}
\LANGsk{
\Question{
Nech
\[ a=\left[\left(2-\sqrt3\right)^{\frac12}+\left(2+\sqrt3\right)^{\frac12}\right]^2,\ b=\frac{81^{-1}\cdot\sqrt3}{27^{-2}\cdot\sqrt[4]9}.\]
Porovnaním čísel  \( a^b \) a \( b^a \) dostaneme:}
{\( a^b > b^a \)}{\( a^b < b^a \)}{\( a^b \leq b^a \)}{\( a^b = b^a \)}{}{}}
\LANGes{
\Question{
Sea
\[ a=\left[\left(2-\sqrt3\right)^{\frac12}+\left(2+\sqrt3\right)^{\frac12}\right]^2,\ b=\frac{81^{-1}\cdot\sqrt3}{27^{-2}\cdot\sqrt[4]9}.\]
Comparando los números \( a^b \) y \( b^a \) te sale:}
{\( a^b > b^a \)}{\( a^b < b^a \)}{\( a^b \leq b^a \)}{\( a^b = b^a \)}{}{}}
\End

\Data\ID{1003164004}
\Updated{2020-10-12 10:10:05}
\Level{C}
\SubArea{Expressions with powers and roots}
\LANGen{
\Question{
Evaluate the expression 
\[ \frac{5x^2+10x+10}{(x+1)^4-1}\]
for \( x=\sqrt5-2 \) and write the result in the form \( a+b\sqrt c \), where \( a \), \( b \), \( c \) are natural numbers.}
{\( 5+2\sqrt5 \)}{\( 2+5\sqrt2 \)}{\( 5+5\sqrt2 \)}{\( 5+5\sqrt5 \)}{}{}}
\LANGpl{
\Question{
Oblicz wartość wyrażenia  
\[ \frac{5x^2+10x+10}{(x+1)^4-1}\]
dla \( x=\sqrt5-2 \) i wynik zapisz w postaci  \( a+b\sqrt c \), gdzie \( a \), \( b \), \( c \) są liczbami naturalnymi.}
{\( 5+2\sqrt5 \)}{\( 2+5\sqrt2 \)}{\( 5+5\sqrt2 \)}{\( 5+5\sqrt5 \)}{}{}}
\LANGcs{
\Question{
Určete hodnotu výrazu 
\[ \frac{5x^2+10x+10}{(x+1)^4-1}\]
pro \( x=\sqrt5-2 \) a výsledek zapiště v tvaru \( a+b\sqrt c \), kde \( a \), \( b \), \( c \) jsou přirozená čísla.}
{\( 5+2\sqrt5 \)}{\( 2+5\sqrt2 \)}{\( 5+5\sqrt2 \)}{\( 5+5\sqrt5 \)}{}{}}
\LANGsk{
\Question{
Určte hodnotu výrazu 
\[ \frac{5x^2+10x+10}{(x+1)^4-1}\]
pre \( x=\sqrt5-2 \) a výsledok zapíšte v tvare \( a+b\sqrt c \), kde \( a \), \( b \), \( c \) sú prirodzené čísla.}
{\( 5+2\sqrt5 \)}{\( 2+5\sqrt2 \)}{\( 5+5\sqrt2 \)}{\( 5+5\sqrt5 \)}{}{}}
\LANGes{
\Question{
Evalúa la expresión
\[ \frac{5x^2+10x+10}{(x+1)^4-1}\]
para \( x=\sqrt5-2 \) y escribe el resultado en la forma \( a+b\sqrt c \), dónde \( a \), \( b \), \( c \) son números naturales.}
{\( 5+2\sqrt5 \)}{\( 2+5\sqrt2 \)}{\( 5+5\sqrt2 \)}{\( 5+5\sqrt5 \)}{}{}}
\End

\Data\ID{1003164005}
\Updated{2020-08-28 01:08:16}
\Level{C}
\SubArea{Expressions with powers and roots}
\LANGen{
\Question{
Evaluate the expression
\[ \sqrt[3]{\sqrt[3]{10}-\sqrt[3]2}\cdot\sqrt[3]{\sqrt[3]{100}+\sqrt[3]{20}+\sqrt[3]4}. \]}
{\( 2 \)}{\( 8 \)}{\( \sqrt[3]{10}-\sqrt[3]2 \)}{\( \sqrt[3]{100}-\sqrt[3]4 \)}{}{}}
\LANGpl{
\Question{
Liczba 
\[ \sqrt[3]{\sqrt[3]{10}-\sqrt[3]2}\cdot\sqrt[3]{\sqrt[3]{100}+\sqrt[3]{20}+\sqrt[3]4} \]
jest równa:}
{\( 2 \)}{\( 8 \)}{\( \sqrt[3]{10}-\sqrt[3]2 \)}{\( \sqrt[3]{100}-\sqrt[3]4 \)}{}{}}
\LANGcs{
\Question{
Určete hodnotu výrazu
\[ \sqrt[3]{\sqrt[3]{10}-\sqrt[3]2}\cdot\sqrt[3]{\sqrt[3]{100}+\sqrt[3]{20}+\sqrt[3]4}. \]}
{\( 2 \)}{\( 8 \)}{\( \sqrt[3]{10}-\sqrt[3]2 \)}{\( \sqrt[3]{100}-\sqrt[3]4 \)}{}{}}
\LANGsk{
\Question{
Určte hodnotu výrazu
\[ \sqrt[3]{\sqrt[3]{10}-\sqrt[3]2}\cdot\sqrt[3]{\sqrt[3]{100}+\sqrt[3]{20}+\sqrt[3]4}.\]}
{\( 2 \)}{\( 8 \)}{\( \sqrt[3]{10}-\sqrt[3]2 \)}{\( \sqrt[3]{100}-\sqrt[3]4 \)}{}{}}
\LANGes{
\Question{
Evalúa la expresión
\[ \sqrt[3]{\sqrt[3]{10}-\sqrt[3]2}\cdot\sqrt[3]{\sqrt[3]{100}+\sqrt[3]{20}+\sqrt[3]4}. \]}
{\( 2 \)}{\( 8 \)}{\( \sqrt[3]{10}-\sqrt[3]2 \)}{\( \sqrt[3]{100}-\sqrt[3]4 \)}{}{}}
\End

\Data\ID{2010004809}
\Updated{2022-11-02 09:11:34}
\Level{C}
\SubArea{Expressions with powers and roots}
\LANGen{
\Question{
Evaluate the following expression at \(x = 9\).
\[ 
 \frac{x^{-2}} 
{x^{-\frac{1}{2}} - x^{-1}}
\]}
{\(\frac{1} 
{18}\)}{\(-\frac{1}
 {18} \)}{\(\frac{27}
{2}\)}{\(-\frac{27}
{8}\)}{}{}}
\LANGpl{
\Question{
Oblicz wyrażenie, gdzie \(x = 9\).
\[ 
 \frac{x^{-2}} 
{x^{-\frac{1}{2}} - x^{-1}}
\]}
{\(\frac{1} 
{18}\)}{\(-\frac{1}
 {18} \)}{\(\frac{27}
{2}\)}{\(-\frac{27}
{8}\)}{}{}}
\LANGcs{
\Question{
Určete hodnotu  výrazu
\[ 
 \frac{x^{-2}} 
{x^{-\frac{1}{2}} - x^{-1}}
\]  pro \(x = 9\).}
{\(\frac{1} 
{18}\)}{\(-\frac{1}
 {18} \)}{\(\frac{27}
{2}\)}{\(-\frac{27}
{8}\)}{}{}}
\LANGsk{
\Question{
Vypočítajte hodnotu výrazu pre \(x = 9\).
\[ 
 \frac{x^{-2}} 
{x^{-\frac{1}{2}} - x^{-1}}
\]}
{\(\frac{1} 
{18}\)}{\(-\frac{1}
 {18} \)}{\(\frac{27}
{2}\)}{\(-\frac{27}
{8}\)}{}{}}
\LANGes{
\Question{
Evalúa la siguiente expresión en \(x = 9\).
\[ 
 \frac{x^{-2}} 
{x^{-\frac{1}{2}} - x^{-1}}
\]}
{\(\frac{1} 
{18}\)}{\(-\frac{1}
 {18} \)}{\(\frac{27}
{2}\)}{\(-\frac{27}
{8}\)}{}{}}
\End

\Data\ID{9000079209}
\Updated{2022-04-23 05:04:14}
\Level{C}
\SubArea{Expressions with powers and roots}
\LANGen{
\Question{
Evaluate the following expression at \(x = 4\).
\[ 
 \frac{x^{-\frac{1} 
{2} }} 
{x^{-2} - x^{-1}}
\]}
{\(-\frac{8} 
{3}\)}{\(\frac{31}
 {3} \)}{\(\frac{8}
{3}\)}{\(6\)}{}{}}
\LANGpl{
\Question{
Oblicz następujące wyrażenie dla \(x = 4\).
\[ 
 \frac{x^{-\frac{1} 
{2} }} 
{x^{-2} - x^{-1}}
\]}
{\(-\frac{8} 
{3}\)}{\(\frac{31}
 {3} \)}{\(\frac{8}
{3}\)}{\(6\)}{}{}}
\LANGcs{
\Question{
Určete hodnotu výrazu \(\frac{x^{-\frac{1} {2} }} 
{x^{-2}-x^{-1}} \) 
pro \(x = 4\).}
{\(-\frac{8} 
{3}\)}{\(\frac{31}
 {3} \)}{\(\frac{8}
{3}\)}{\(6\)}{}{}}
\LANGsk{
\Question{
Určte hodnotu výrazu \(\frac{x^{-\frac{1} {2} }} 
{x^{-2}-x^{-1}} \) 
pre \(x = 4\).}
{\(-\frac{8} 
{3}\)}{\(\frac{31}
 {3} \)}{\(\frac{8}
{3}\)}{\(6\)}{}{}}
\LANGes{
\Question{
Evalúa la expresión en \(x = 4\).
\[ 
 \frac{x^{-\frac{1} 
{2} }} 
{x^{-2} - x^{-1}}
\]}
{\(-\frac{8} 
{3}\)}{\(\frac{31}
 {3} \)}{\(\frac{8}
{3}\)}{\(6\)}{}{}}
\End

\Data\ID{9000085602}
\Updated{2020-10-12 10:10:08}
\Level{C}
\SubArea{Expressions with powers and roots}
\LANGen{
\Question{
Evaluate the following numbers and round to the nearest tens. 
\[ 
 \left [(2^{2})^{2}\right ]^{2}
\]}
{\(260\)}{\(510\)}{\(120\)}{\(60\)}{}{}}
\LANGpl{
\Question{
Liczba 
\[ 
 \left [(2^{2})^{2}\right ]^{2}
\] po zaokrągleniu do dziesiątek jest równa:}
{\(260\)}{\(510\)}{\(120\)}{\(60\)}{}{}}
\LANGcs{
\Question{
Číslo \(\left [(2^{2})^{2}\right ]^{2}\) 
je po zaokrouhlení na desítky rovno:}
{\(260\)}{\(510\)}{\(120\)}{\(60\)}{}{}}
\LANGsk{
\Question{
Číslo \(\left [(2^{2})^{2}\right ]^{2}\) 
sa po zaokrúhlení na desiatky rovná:}
{\(260\)}{\(510\)}{\(120\)}{\(60\)}{}{}}
\LANGes{
\Question{
Evalúa el número y aproximalo a la decena más cercana.
\[ 
 \left [(2^{2})^{2}\right ]^{2}
\]}
{\(260\)}{\(510\)}{\(120\)}{\(60\)}{}{}}
\End
